\documentclass[11pt,a4paper]{article}
\usepackage[T1]{fontenc}\usepackage[utf8]{inputenc}\usepackage{lmodern}
\usepackage[margin=23mm,headheight=16pt]{geometry}
\usepackage{amsmath,amssymb,amsthm,mathtools,bm,booktabs,tabularx,longtable,microtype,graphicx,xcolor}
\usepackage{hyperref,fancyhdr,enumitem,slashed,mathrsfs}
\definecolor{navy}{HTML}{18344F}\definecolor{amber}{HTML}{825100}
\hypersetup{colorlinks=true,linkcolor=navy,urlcolor=navy,pdftitle={QGD Chapter 6 v2.6: Action-first theory, quantum kernels, natural-invariant response},pdfauthor={QGD Research Programme}}
\pagestyle{fancy}\fancyhf{}\lhead{\small QGD $\mid$ New Chapter 6}\rhead{\small Action-first / quantum \& dark energy v2.7}\cfoot{\thepage}
\setlength{\parindent}{0pt}\setlength{\parskip}{5pt}
\emergencystretch=2em
\newtheorem{proposition}{Proposition}[section]\newtheorem{lemma}[proposition]{Lemma}\newtheorem{definition}[proposition]{Definition}\newtheorem{remark}[proposition]{Remark}
\newcommand{\Gam}{\Gamma}\newcommand{\cD}{\mathcal D}\newcommand{\dd}{\mathrm d}\newcommand{\Tr}{\operatorname{Tr}}\newcommand{\cE}{\mathcal E}\newcommand{\cK}{\mathcal K}\newcommand{\cQ}{\mathcal Q}\newcommand{\cS}{\mathcal S}\newcommand{\cR}{\mathcal R}\newcommand{\sig}{\sigma}
\begin{document}
\begin{center}
{\color{navy}\Large\bfseries Quantum Gravity Dynamics}\\[5pt]
{\LARGE\bfseries Chapter 6}\\[4pt]
{\large Quantum Action, Natural-Invariant Dynamics and Historical Correspondence Tests}\\[6pt]
{\small Research edition v2.7 (action-first; QM043 + QN001--QN005 + Q052--Q054) \quad $\mid$ \quad 5 November 2026}\\[3pt]
{\small Action-first synthesis; Chapters 1--5 define the effective gravitational theory}
\end{center}
\vspace{5pt}\hrule\vspace{8pt}
\section*{Action-first charter: quantum response and the status of the historical tests (5 November 2026)}
\addcontentsline{toc}{section}{Action-first charter: quantum response and historical tests}
\noindent\textbf{Revised role of this chapter.} QGD is an action-defined gravitational alternative formulated with independent natural-invariant one-forms. Chapter~1 states \emph{Axiom I}, the complete founding master metric, and \emph{Axiom II}, the unproved existence of an independent microscopic counterpart to the phenomenological map. Neither Axiom~II nor a successful $\mathcal F_{\rm micro}$ construction is needed to calculate QGD's regular metric-mediated gravitational equations and classical solutions. The core dynamics comes from the declared Einstein--Hilbert--matter action composed with $g(\Sigma)$ and its renormalized quantum-effective extension; Chapter~4 gives its defining variational law.

\noindent\textbf{Primary current result.} Section~40 (QN001--QN005) contains the four \emph{explicitly computed} proper-time Dirac/Hartree quadratic kernels $\mathcal F_C^\Lambda$, $\mathcal F_R^\Lambda$, $\mathcal V_\Lambda$ and $K_\Lambda$, the Schur-complement curvature response $\mathcal F_R^\Lambda-\mathcal V_\Lambda^2/(2K_\Lambda)$, and the exact pullback of the corresponding metric Hessian to natural-invariant perturbations. Those results are conditional on the stated Euclidean flat-background quantum model; they do not define a complete causal Lorentzian nonlocal functional on all geometries. The independent finite local curvature couplings remain renormalization data.

\noindent\textbf{Status of QM001--QM043.} The remaining sections of this research edition preserve earlier and continuing studies of possible realizations of Axiom~II. They are valuable \emph{optional tests}, not missing prerequisites for the validity of the effective QGD gravitational action. The phenomenological map $\mathcal F_{\rm phen}$ is the source/state/boundary-to-$[\Sigma]$ solution correspondence of the specified action, modulo metric-preserving representation redundancies. The microscopic map must be independent of that field-equation solver to count as an actual realization, and its existence has not been established.

\noindent\textbf{Mathematical hygiene.} An eight-sector regular background has a ten-dimensional metric image and 22 representation directions. Near the all-zero sector, the quadratic metric map has vanishing first derivative; the projected Euler equations may then hold even where the Einstein residual does not. Physical equivalence with GR is therefore conditional on a declared admissible, sufficiently regular sector domain. The historical $\lambda\mathcal Q_{ab}$ slot remains in Axiom~I but is set to zero in the controlled QN004--QN005 metric-only calculations.


\medskip\noindent\textbf{Historical and research-edition synopsis (preserved).}\par\smallskip
\noindent\textbf{Editorial update (QN001--QN005, 4 November 2026).} An additional dedicated section consolidates the analytical quantum effective action, common-regulator Dirac/Hartree response, and the exact natural-invariant pullback (flat signed sectors and weak-static spherical branch). It distinguishes calculated Euclidean kernels from unresolved Lorentzian causal, finite-renormalization, and microscopic-emergence problems. All QM001--QM043 content is preserved.\par\medskip
\noindent\textbf{Earlier editorial update (QM043).} This edition adds a conditional coarse-graining, large-number, and collective-basis refinement to the inverse matching map while retaining all earlier material. It does not claim microscopic closure.\par\medskip
\textbf{Earlier consolidation (QM042).} This v2.3 edition preserves the entire QM001--QM039 chapter and appends the independent QM032--QM036 audits previously archived outside it, QM040--QM041 microscopic field calculations, and a definitive two-map/inverse-matching status ledger. \emph{Phenomenological $\mathcal F_{\rm phen}$ is conditionally closed under the adopted effective action; derived microscopic $\mathcal F_{\rm QM}$ is not established.} The original abstract follows unchanged, describing the historical sequence of calculations.
\vspace{7pt}
\textbf{Abstract.} Chapter 5 proves a conditional static matter-wave-to-QGD-geometric correspondence when a stable source-coupled collective response is supplied as premise MW2. This chapter begins to test that physical premise with the actual interacting quantum-matter model of Chapter 1, using the independently developed microscopic spectral and WKB benchmarks QM001--QM005. We retain the exact leading Dirac spinor ratio, the conditional WKB phase transport, and the stationary current bridge $r^2j^r=C_J$ rather than replacing them. We specify the full physical source-response matrix and its infrared-pole criterion; re-evaluate the scalar-contact Hartree saddle using precisely Chapter 5's cutoff convention; and verify an analytically finite density response and a zero source-visible long-range residue in the tested scalar channel. We additionally prove a transverse-traceless scalar-mixing selection rule for an isotropic state and a local-map exterior obstruction. The earlier tensor increment independently verifies the collisionless TT shear continuum and a transverse-momentum kinetic Ward identity with its necessary local contact. These negative tests identify why a physically selected Dirac quasiparticle is not yet a universally gravitating collective field. The present extension constructs the full ten-indexed leading intraband stress matrix, defines the covariantly regulated fermion-determinant bubble and genuine metric seagull, computes its local vacuum heat-kernel terms, verifies a finite-coordinate-deformation Ward cancellation, refits the Hartree saddle under the changed regulator, and tests the scalar-resummed TT channel. It then specifies the remaining genuinely dynamical tensor and universality tests. Neither a massless graviton nor Einstein dynamics is assumed in the microscopic kernel. The founding complete coframe/signed-sector/$\cQ$ metric remains unchanged. The v0.5 extension proves an exact null-sector Clifford/coframe bridge, preserves the covariant Dirac/WKB/current equation, and derives the source-to-kinetic functional kernel that an independently microscopic promotion law would have to reproduce. It explicitly distinguishes matching within the adopted gravitational action from deriving that action microscopically. All new physical emergence claims remain open. The v2.2 QM039 addition exactly solves a small interacting derivative-bond model, computes nonzero causal cross-species response, and tests its self-consistent Hartree kinetic frame and RPA collective pole. Its uniform saddle fails the infinite-chain staggered bond stability test; no universal geometric wave is claimed. The v2.1 QM038 addition tests the full-master-metric Diracization, the Clifford-compatible connection, a conformal Hermitian WKB transport coefficient, inverse-microscopic ambiguity and long-distance frame locking. It introduces no unproved claim of a dynamically selected collective phase. The v2.0 QM037 extension adds exact Clifford/principal-symbol tests for curved Dirac fields, a conformal spin-connection current counterexample, constrained Fierz bilinears, and a conditional shared kinetic composite channel; no microscopic exterior geometry is claimed. The v1.9 QM034 extension exactly solves a unitary two-site four-flavour determinant-density model, exhibiting true causal cross-species response and a fourth-order soft collective excitation, but also a vanishing soft-pole numerator, a gapped homogeneous composite, spectator non-universality and vacuum-energy shift ambiguity; this remains a projected analogue rather than a derivation of geometry from the full four-dimensional microscopic model. The v0.5 QM009 extension derives the inhomogeneous scalar and number-density susceptibility, independently checks matched smooth regulators and sharp-cutoff artifacts, and proves the absence of a Hartree--Fock source-to-kinetic deformation even for spatially varying mean fields. The v0.7 QM010 extension sharpens the central research objective: the conditional Schwarzschild source-to-natural-invariant bridge is already closed under the adopted QGD field action, while the genuinely microscopic spinor/state-to-invariant selection map remains open. Motivated by Hamiltonian stationary action and Fermat optics, it derives the weak-field $H$ profile as a stationary source-conditioned field functional, proves its spherical uniqueness under stated boundary conditions, separates path variation from field variation, and gives the exact conditional Hessian/source-response form required for a microscopic selection principle. No graviton or independent tensor mode is assumed as a prerequisite. QM011 tests the proposed microscopic variational-selection functional directly: its Hubbard--Stratonovich stationary scalar is physically determined, but the measured positive curvature of the scalar Hessian prevents the required QGD Schwarzschild exterior selection in the Hartree--Gaussian regime. This test does not reopen the conditional QGD correspondence. The v0.8 research synthesis separately audits QM012--QM014, including the nonlinear stiffness of the regular $H=\sigma_{\rm amp}^2$ coordinate, the original versus induced-action research arrows, exact rank/gauge checks, and the status of the proposed but unadopted premise (Q). It preserves the effective static milestone and the open microscopic selection claim. The v0.9 QM015 extension adds a source-overlap pole criterion, two gapless-but-source-invisible controls, and an exactly solvable mixed-channel stationary response which reproduces the conditional $1/r$ exterior despite the physical source coupling only to a massive channel. The mixed-channel parameters and symmetry are explicitly assumed in the toy calculation; they have not been derived from the original interacting Dirac theory. The v1.0 QM016 extension now evaluates the self-consistent original scalar-contact Dirac model and proves an exact positive-curvature identity on the specified Hartree cutoff domain, computes a full two-coordinate static scalar--density thermodynamic Hessian and independent susceptibility checks, and establishes that coupling tuning in this regime does not provide the missing source-visible massless channel. The full momentum-dependent stress/kinetic 2PI problem remains open. QM017--QM019 now follow the newly proposed kinetic route through an explicitly scoped, nonlocal screened-exchange test, a positive kinetic derivative accompanied by a short-range failure, an exact cross-sector retarded Kubo obstruction for a decoupled spectator, an isotropic scalar-to-TT resummation obstruction, and an algebraic Clifford/signed-lift check proving that a species-specific medium metric is not universal gravitational emergence. All results retain their declared approximation scopes; no full Ward-completed stress--kinetic 2PI computation has been claimed. QM022 establishes an exact projective-spinor representation of the twisting Kerr null direction, its flat-reference ray transport and optical expansion, and a restricted optical-divergence identity for the Kerr/Schwarzschild gravitational weight. It proves why conserved chiral current is insufficient, why averaged mixed spinors do not pick a unique null congruence, and why an unrestricted divergence rule fails the vacuum Einstein equations. These are constructive geometric/kinematic advances, not a microscopic derivation of a common physical metric. The v1.3 QM023 consolidation incorporates the explicit working hypothesis of an actual collective quantum field extending into exterior vacuum; it retains a complete proof-status ledger, separates the frame-direction and frame-weight tasks, and poses the precise central question for $\mathcal F_{\rm micro}$. No new microscopic field operator is claimed to have been derived.

\textbf{Proof-status contract.} \textbf{[D]} denotes an exact identity, a conditional mathematical lemma with its hypotheses stated, or a reproduced numerical result within a named model. \textbf{[C]} denotes a conclusion conditional on that model, approximation, or an adopted infrared action. \textbf{[A]} denotes an independently assumed physical ingredient. \textbf{[O]} denotes an unproved microscopic QGD claim. No ``derived gravity'' is inferred from static response finiteness, WKB transport, algebraic factorization, or an Einsteinian metric pullback.
\begin{center}
{\color{navy}\bfseries QM023 RESEARCH COMPASS --- NEW PHYSICAL HYPOTHESIS}
\end{center}
\noindent\textbf{Working ontology [A]:} The natural-invariant-producing structure is to be investigated as an \emph{actual quantum collective field extending through the exterior vacuum}, rather than only a label reconstructed from a source state or a known Einstein solution. Its operator definition, source coupling, causal field equation, and physical link to $[\Sigma]$ remain open.\\[4pt]
\noindent\textbf{Central question [O]:} Can the \emph{original microscopic quantum-matter action plus source state and boundary conditions} generate an exterior collective field which \emph{selects} the weighted directions and weights of a physical natural-invariant class and thereby a \emph{common} effective metric, without assuming that metric or the Einstein equations as inputs?\\[4pt]
\noindent\textbf{Result to date:} Mathematical metric representation and conditional Einsteinian dynamics: established. Spinor-optical and real microscopic response results: established within stated model and approximation domains. Physical exterior-field selection and universality: \emph{not derived}. QM023--QM038 at the end of this chapter consolidates the achievements and establishes an explicit fitted effective selection law as a target for backward micro-derivation.

\tableofcontents\newpage
\section{Why a new Chapter 6 belongs here}
The logical sequence is now: Chapter 1 (microscopic fields, Dirac spinors, large-$N$, WKB and conserved currents); Chapter 2 (complete natural-invariant metric and selected geometric solutions); Chapter 3 (gravitational empirical normalization); Chapter 4 (adopted effective QGD action, quantum stress and full metric pullback); Chapter 5 (backward ladder and conditional static correspondence); \emph{this Chapter 6} (test of microscopic origin and source-visible response); formerly Chapter 6, now provisionally Chapter 7 (vacuum and physical energy); formerly Chapter 7, now provisionally Chapter 8 (cosmology). No original chapter is superseded by this research edition.

The Chapter 5 static milestone has a clearly identified physical premise: a collective coordinate $q$ with
\begin{equation}\label{eq:MW2}
\boxed{\cE[q;\rho_m]=\int\dd^3x\left[\frac Z2|\nabla q|^2-g_cq\rho_m\right],\qquad Z>0,}\end{equation}
which conditionally yields $G_{\rm eff}=g_c^2/(4\pi Z)$. Here $q$ is a gravitationally \emph{candidate} collective coordinate, not the Dirac quasiparticle, the Pauli-block ratio, or the historical symmetric tensor $\cQ_{ab}$. Equation~\eqref{eq:MW2} is a \textbf{[A]} premise until derived from microscopic QFT.

\section{The precise missing intersection: a physical source-to-operator-to-invariant map}\label{sec:physicalF}
\textbf{Status:} This section sharpens the research question of this chapter. It does not replace the Chapter~1 spinor construction or the Chapter~2 dynamically selected exterior natural invariant. The operative distinction is between a \emph{quantum source state} $\hat\varrho_{\rm src}$ and a \emph{probe spinor} $\Psi_{\rm pr}$ propagating through the effective background. Their state-dependent and kinematic data must not be conflated.

\paragraph{What is already present.}
For a specified effective Dirac operator $\mathcal D_{\rm eff}$, solving $\mathcal D_{\rm eff}\Psi_{\rm pr}=0$ yields the spinor ratios $\mathsf R$, WKB phase $S$, and current $j^\mu$. Separately, the Chapter~2 adopted Einstein--Hilbert QGD action dynamically selects the Schwarzschild natural invariant within its declared independent-sector ansatz. What has \emph{not} been derived is an independent physical map
\begin{equation}\label{eq:F-gap}
 \boxed{\mathcal F:\quad\bigl(\hat\varrho_{\rm src},\mathcal D_{\rm eff},S,\mathsf R,j,\text{boundary data}\bigr)
 \longmapsto[\Sigma]_{\rm physical},}
\end{equation}
where $[\Sigma]$ denotes physically indistinguishable signed-sector representations after the reference coframe and allowed redundancies are specified. A particular positive-energy lower/upper spinor ratio is a \emph{probe-dependent} kinematic quantity, not automatically the source's gravitational one-form.

\paragraph{Two logically distinct ways to complete Eq.~\eqref{eq:F-gap}.}
\begin{enumerate}[label=\textbf{Route \Alph*.},leftmargin=14mm]
\item \textbf{Conditional effective QGD matching [C].} Compute $T^{\mu\nu}_{\rm src}$ from the quantum source; solve the \emph{already adopted} QGD gravitational equations with boundary conditions; obtain $[\Sigma]$; construct the compatible Dirac operator; then test its spinors, WKB and conserved currents. This is a noncircular \emph{effective-model calculation} if the gravitational equation is explicitly listed as an independent premise. It does not derive that equation from QFT.
\item \textbf{Independent microscopic source-to-kinetic map [O].} Compute a probe-independent effective kinetic principal symbol from microscopic matter/source dynamics, without inserting $g$ or $\Sigma$ as background inputs. Demonstrate that this symbol determines a common effective geometry and physically selects the QGD natural-invariant sectors (with their coframe and $\cQ$ histories), including exterior memory of a compact source. This is the stronger physical-emergence claim.
\end{enumerate}
These routes have different burdens of proof. An independently propagated composite spin-two particle is \emph{one possible} mechanism for Route B, not a necessary ingredient of Route A.

\subsection{An exact null-sector operator matching lemma [D]}
Fix the reference Lorentz frame, take the classical $\lambda\cQ=0$ specialization of the \emph{complete} QGD metric, and select one positive null sector. Let $\ell_\mu$ be a flat-null covector, $\eta^{\mu\nu}\ell_\mu\ell_\nu=0$, and $H(x)\ge0$. Define
\begin{equation}\label{eq:KS-op}
 \sigma_\mu=\sqrt H\,\ell_\mu,\quad
 g_{\mu\nu}=\eta_{\mu\nu}+H\ell_\mu\ell_\nu,\quad
 \Gamma^\mu=\gamma^\mu-\frac H2\ell^\mu\not\!\ell,\quad
 \not\!\ell:=\ell_\nu\gamma^\nu.
\end{equation}
The complete founding metric still reads
\begin{equation}\label{eq:full-preserved}
 \boxed{g_{\mu\nu}=T^a{}_{\mu}T^b{}_{\nu}
 [\eta_{ab}-\sum_s\varepsilon_s\sigma_a^{(s)}\sigma_b^{(s)}-\lambda\cQ_{ab}],\quad s_s=-\varepsilon_s.}
\end{equation}
Equation~\eqref{eq:KS-op} is a \emph{specialization}, not a replacement for Eq.~\eqref{eq:full-preserved}.

\begin{proposition}[Exact principal-symbol and coframe correspondence]\label{prop:F-KS}
With the Chapter~1 convention $\{\gamma^a,\gamma^b\}=-2\eta^{ab}I$, the matrices of Eq.~\eqref{eq:KS-op} satisfy
\begin{equation}\label{eq:cliffF}
 \boxed{\{\Gamma^\mu,\Gamma^\nu\}=-2g^{\mu\nu}I,
 \qquad g^{\mu\nu}=\eta^{\mu\nu}-H\ell^\mu\ell^\nu.}
\end{equation}
An explicit inverse coframe is $e_a{}^\mu=\delta_a{}^\mu-\tfrac H2\ell_a\ell^\mu$ and $\theta^a{}_{\mu}=\delta^a{}_{\mu}+\tfrac H2\ell^a\ell_\mu$ in the specified flat reference coordinates. The metric determinant obeys $\det g=\det\eta$ in Cartesian Kerr--Schild coordinates.
\end{proposition}
\begin{proof}
The nilpotent rank-one perturbation $N_a{}^b=\ell_a\ell^b$ satisfies $N^2=0$ because $\ell^2=0$, hence $(I-\tfrac H2N)^{-1}=I+\tfrac H2N$. Expanding $\{\Gamma^\mu,\Gamma^\nu\}$ with $\{\gamma^\mu,\not\!\ell\}=-2\ell^\mu I$ gives the linear term $+2H\ell^\mu\ell^\nu I$; the quadratic term vanishes since $\not\!\ell^2=-\ell^2 I=0$. This equals $-2(\eta^{\mu\nu}-H\ell^\mu\ell^\nu)I$. The determinant lemma gives $\det(\eta+H\ell\ell^\top)=\det\eta\,(1+H\ell^2)=\det\eta$.
\end{proof}
\textbf{Meaning.} This is an exact \emph{operator representation} of a chosen QGD null-sector geometry. It is not a derivation of $H$ from quantum matter, nor a proof that $\Gamma^\mu$ is generated by the original scalar-contact QFT. The operator's connection and all other species' propagation must also be supplied consistently.

\subsection{Curved Dirac dynamics, WKB, and the original bridge [C]}
To obtain a covariant wave equation, the matrices in Eq.~\eqref{eq:KS-op} are not by themselves sufficient. Using the torsion-free spin connection $\Omega_\mu[e]$ of the specified coframe, take
\begin{equation}\label{eq:matched-Dirac}
 \boxed{\mathcal D_g\Psi
 :=\bigl[i\hbar\Gamma^\mu(\partial_\mu+\Omega_\mu)-mc\bigr]\Psi=0.}
\end{equation}
The connection term is required for local Lorentz covariance and the full current equation. Writing $\Psi=e^{-iS/\hbar}(a_0+O(\hbar))$ gives $p_\mu=\partial_\mu S$ with the adopted phase convention and, at leading order,
\begin{equation}\label{eq:HJ-KS}
 \boxed{\bigl(g^{\mu\nu}p_\mu p_\nu+m^2c^2\bigr)a_0=0\quad\text{as a scalar mass shell}.}
\end{equation}
The notation means $g^{\mu\nu}p_\mu p_\nu+m^2c^2=0$ wherever $a_0\ne0$; it is not a new matrix equation beyond the Dirac principal symbol. Spin transport and the spin connection enter beyond leading eikonal order.

Let $J^\mu=\bar\Psi\Gamma^\mu\Psi$ in the convention of Chapter~1. For a Hermitian compatible curved Dirac action and a real mass, $\nabla_\mu J^\mu=0$. In static spherical symmetry $\sqrt{-g}=r^2\sin\theta$ in these coordinates; consequently
\begin{equation}\label{eq:curved-bridge}
 \boxed{\partial_r(r^2J^r)=0,\qquad r^2J^r=C_J.}
\end{equation}
Thus the \emph{exact integrated current bridge survives} the null-sector geometric matching. The local relation $J^r\simeq\rho p_r c^2/E$ must, however, be replaced by its appropriate tetrad/observer-dependent curved-space WKB closure; copying the flat radial velocity formula unchanged into Kerr--Schild coordinates is generally incorrect. The four-mode spinor relations remain local algebraic ingredients after a compatible tetrad is fixed.

\subsection{An explicit independent Hamilton--Jacobi verification [D/C]}
Set $x^0=ct$ and $c=1$ temporarily. For $\ell_\mu dx^\mu=dx^0+dr$ in a radial exterior patch,
\begin{equation}\label{eq:KS-components}
 ds^2=-(1-H)(dx^0)^2+2Hdx^0\dd r+(1+H)\dd r^2+r^2\dd\Omega^2.
\end{equation}
For a radial eikonal $S=-Ex^0+W(r)$, $p_r=W'(r)$, the Dirac principal symbol gives
\begin{equation}\label{eq:radial-shell}
 \boxed{(1-H)p_r^2-2HEp_r-(1+H)E^2+m^2=0,}
\end{equation}
with two branches, in the exterior $H<1$,
\begin{equation}\label{eq:radial-roots}
 \boxed{p_r^{\pm}=\frac{HE\pm\sqrt{E^2-(1-H)m^2}}{1-H}.}
\end{equation}
In the massless limit $p_r^+=E(1+H)/(1-H)$ and $p_r^-=-E$; both satisfy the same null characteristic cone. As $H\to0$, $p_r^\pm\to\pm\sqrt{E^2-m^2}$. At $r=4r_s$ ($H=1/4$), $E=3/2$ and $m=1$, the outgoing/ingoing roots are approximately $2.133$ and $-1.133$, respectively. These expressions were checked independently by substitution in Eq.~\eqref{eq:radial-shell}. The resulting momentum change is geometric \emph{once $H$ has been provided}; it is not evidence that the microscopic matter theory independently selected $H$.


\section{Why a gravitational natural invariant is selected: Hamilton, Fermat, and QGD (QM010)}\label{sec:selection-principle}
\textbf{Research question.} Chapter~2 already selects the Schwarzschild natural-invariant amplitudes \emph{dynamically under the adopted Einstein--Hilbert QGD action}. Chapter~1 supplies microscopic and effective Dirac structure, its WKB mass shell, the component ratio $u_B=\mathsf Ru_A$, and the current bridge $r^2j^r=C_J$. The remaining question is not whether those effective results are invalid, but whether the quantum matter theory itself supplies the \emph{selection functional} relating a specified source state to a physical $[\Sigma]$. This section makes the distinction precise and records the conditional selection result as closed.

\subsection{Stationary optical paths and stationary gravitational fields are distinct problems}
In a prescribed isotropic optical medium of refractive index $n(\bm x)$, Fermat's principle is
\begin{equation}\label{eq:qm010-fermat}
\delta_{\gamma}\mathfrak T[\gamma;n]
=\delta_{\gamma}\left(\frac1c\int_{\gamma} n(\bm x)\,\dd s\right)=0,
\qquad\text{with $n$ fixed}.\end{equation}
Its Euler--Lagrange equations determine ray paths $\gamma$ for that medium, including Snell's law. They do \emph{not} determine $n(\bm x)$ from the rays alone. Likewise, $\delta_\gamma\int\dd s_g=0$ gives geodesics for a \emph{specified} metric $g$, not a field equation for $g$. For an individual physical history, Hamilton's principle reads $\delta\mathfrak S=0$ under the appropriate endpoint and dynamical constraints; it is generally \emph{not} the statement $\delta H_{\rm Hamiltonian}=0$, and stationarity does not imply an absolute minimum. The WKB phase $S_{\rm eik}$ must not be confused with the action functional $\mathfrak S$. See also the general discussion of stationary-action principles in \url{https://www.scholarpedia.org/article/Principle_of_least_action}.

Thus the natural analogy for the \emph{selection of the QGD medium itself} is not only ray stationarity but a separate \emph{field} stationarity condition. In the adopted QGD effective theory it is
\begin{equation}\label{eq:qm010-field-selection}
\boxed{\left.\frac{\delta S_{\rm QGD}[g[\Sigma],\psi,A]}{\delta\sigma^{(s)}_\mu(x)}\right|_{\substack{\text{specified source},\,\text{boundary data}\\\text{independently varied sectors}}}=0.}\end{equation}
This condition is a \textbf{[C]} selection equation because its gravitational action is already supplied. It is not an independently derived matter-to-$[\Sigma]$ law.

\subsection{An explicit source-conditioned stationary functional for the Schwarzschild amplitude [C]}\label{sec:qm010-H}
Let $H_{\rm KS}(\bm x)$ denote the \emph{dimensionless gravitational Kerr--Schild amplitude}, not the archival integrated matter energy $H_{\rm orig}$ of Chapter~1. In the weak-field static benchmark set $H_{\rm KS}=-2\Phi/c^2$. For a specified compact mass-density distribution $\rho_m$ computed from the source quantum state, consider the \emph{adopted Newtonian-limit field functional}
\begin{equation}\label{eq:qm010-energy}
\boxed{\mathfrak E_{\rm N}[H_{\rm KS};\rho_m]
=\int_{\mathbb R^3}\!\dd^3x\,
\left[\frac{c^4}{32\pi G_N}|\nabla H_{\rm KS}|^2
-\frac{c^2}{2}\rho_m H_{\rm KS}\right].}\end{equation}
For variations $\delta H_{\rm KS}$ vanishing at the boundary (or decaying sufficiently rapidly), integration by parts gives
\begin{align}
\delta\mathfrak E_{\rm N}
&=\int\dd^3x\left[-\frac{c^4}{16\pi G_N}\nabla^2H_{\rm KS}-\frac{c^2}{2}\rho_m\right]\delta H_{\rm KS},\label{eq:qm010-firstvar}\\
\therefore\quad\boxed{\nabla^2H_{\rm KS}=-\frac{8\pi G_N}{c^2}\rho_m.}\label{eq:qm010-poisson}
\end{align}
This gives the conditional source-to-amplitude map
\begin{equation}\label{eq:qm010-solution}
\boxed{H_{\rm KS}(\bm x)=\frac{2G_N}{c^2}\int_{\mathbb R^3}
\frac{\rho_m(\bm y)}{|\bm x-\bm y|}\,\dd^3y,}
\end{equation}
with $H_{\rm KS}\to0$ at infinity. For a compact spherical source with total mass $M_s=\int\rho_m\,\dd^3x$, Gauss' law fixes the unique spherically symmetric exterior solution
\begin{equation}\label{eq:qm010-schwarzschild}
\boxed{H_{\rm KS}(r)=\frac{2G_NM_s}{c^2r}\quad(r>R),\qquad
\sigma_\mu\dd x^\mu=\sqrt{H_{\rm KS}}\,[\dd(ct)+\dd r].}
\end{equation}
\begin{proposition}[Conditional source-and-boundary selection, with its domain]\label{prop:qm010-unique}
For $G_N>0$, a regular compact nonnegative source $\rho_m$ with the stated decay at infinity and the Euclidean Newtonian-limit energy \eqref{eq:qm010-energy}, the stationary solution $H_{\rm KS}$ is unique within the permitted regular field class. In the spherically symmetric exterior, Eq.~\eqref{eq:qm010-schwarzschild} follows without separately postulating an exterior profile. This proves \emph{conditional} variational selection within the adopted weak-field theory, not microscopic emergence of its field functional.
\end{proposition}
\begin{proof}
Stationarity gives Eq.~\eqref{eq:qm010-poisson}, whose Green-function solution with the stipulated boundary condition is Eq.~\eqref{eq:qm010-solution}. If $H_1,H_2$ are two such solutions, their difference $h=H_1-H_2$ satisfies $\nabla^2h=0$ with the imposed regularity and decay; integration of $h\nabla^2h$ by parts gives $\int|\nabla h|^2=0$, hence $h=0$. Moreover, the second variation is $\delta^2\mathfrak E_{\rm N}=(c^4/16\pi G_N)\int|\nabla\delta H|^2\ge0$ under the convention $\delta^2\mathfrak E=\dd^2\mathfrak E[H+t\delta H]/\dd t^2|_{t=0}$. For an exterior spherical harmonic function, $H=a+b/r$; decay requires $a=0$, and integration of Eq.~\eqref{eq:qm010-poisson} fixes $b=2G_NM_s/c^2$. The sign of $\sigma=\pm\sqrt H\,\ell$ and the signed-sector/coframe redundancy require additional choice or quotient, even when $H$ is unique.
\end{proof}
\textbf{Scope.} The functional \eqref{eq:qm010-energy} is a rewriting of the \emph{adopted Newtonian gravitational dynamics}; it was not obtained by integrating out Chapter~1's fermions. Its source density and Schwarzschild exterior do not turn this approximation into the full relativistic stellar field equations. Chapter~2's separately checked Einstein--Hilbert sector variation provides the appropriate exact classical vacuum exterior once the action and its variational completeness conditions are accepted. Note also that varying $H$ is not automatically identical to varying $\sigma=\sqrt H\,\ell$ at $\sigma=0$, where the change of variables can be singular.

\subsection{The complete QGD Euler equation is the actual field-selection equation [C]}\label{sec:qm010-optionC}
For the founding metric, retained in Eq.~\eqref{eq:full-preserved}, define the weight-adjusted metric Euler tensor $\Pi^{ab}$ from the adopted full $S_{\rm QGD}$. Holding the coframe and other independent variables fixed as appropriate, the chain rule gives the general pullback condition
\begin{equation}\label{eq:qm010-optionF}
\boxed{(D_{\sigma^{(s)}}g)^*_{\sqrt{-g}}\bigl(\Pi_{\mathfrak G}+\Pi_{\mathfrak T}+\Pi_{\mathfrak Q}\bigr)=0.}\end{equation}
In the Chapter~6 recognizable Option-C notation (with $s_s=-\varepsilon_s$), it is
\begin{equation}\label{eq:qm010-optionC}
\boxed{2s_s\Pi^{ab}\sigma_b^{(s)}
=\lambda\big[(D_s\mathcal Q)^*_{\sqrt{-g}}\Pi\big]^a.}\end{equation}
For the declared classical specialization $\lambda\mathcal Q=0$, this becomes the projected condition $\Pi^{ab}\sigma_b^{(s)}=0$ for a nonzero sector. The full tensor equation $\Pi^{ab}=0$ implies the projected equations, but the converse is not automatic unless the independently varied metric-map directions have sufficient rank. The variational principle can select one or more $[\Sigma]$ subject to admissibility and boundary data, but the source-dependent quantum wavefunction alone has not been proven to select a unique representative. The classical \emph{Schwarzschild natural invariant is dynamically selected within the adopted action}; it is not merely an algebraic guess.

\subsection{What would make stationary selection genuinely microscopic? [O]}\label{sec:qm010-micro}
A microscopic theory could close the stronger original bridge without introducing a separate elementary graviton, but must specify the quantum collective variables and derive their constrained effective functional from the underlying state and action. Let $\varphi^I$ be independent physical collective coordinates obtained from source-state quantum observables (or equivalence classes of composite operators). A \emph{hypothetical, to-be-derived} effective functional is $\Gamma_{\rm mic}[\varphi;J]$, where $J_A$ couples to independently defined source operators. Only if this functional is actually obtained from microscopic quantum matter may we assert a physical selection law:
\begin{equation}\label{eq:qm010-microstationarity}
\boxed{\frac{\delta\Gamma_{\rm mic}}{\delta\varphi^I}=0,
\qquad [\Sigma]_{\rm physical}=\mathcal P[\varphi_*;\,\text{frame and boundary data}].}\end{equation}
The projection $\mathcal P$ from physical collective coordinates to natural-invariant sectors must be independently defined; a four-column spinor basis, an arbitrary probe velocity or a factorization of an already known metric is not such a proof. For an isolated regular stationary branch, differentiate \eqref{eq:qm010-microstationarity} with respect to the source $J_A$ to obtain
\begin{equation}\label{eq:qm010-hessian-response}
\boxed{\frac{\delta\varphi^I}{\delta J_A}
=-[\mathcal K_{\rm phys}^{-1}]^{IJ}
\frac{\delta^2\Gamma_{\rm mic}}{\delta\varphi^J\delta J_A},\qquad
\mathcal K_{\rm phys, IJ}:=\frac{\delta^2\Gamma_{\rm mic}}{\delta\varphi^I\delta\varphi^J},}
\end{equation}
where inverses are only taken after treating constraints, gauge directions, boundary conditions and any zero modes. This identity is a \emph{conditional implicit-function theorem}, not a microscopic calculation of the kernel. It exhibits precisely where source dependence and nonlocal exterior information would enter. It does \emph{not} prescribe a spin-two particle, and it does not require that every collective variable be an independent propagating field.

An equivalent weak-field matching diagnostic is obtained by differentiating \eqref{eq:qm010-energy}: in Fourier space the \emph{adopted} quadratic amplitude kernel and source vertex are
\begin{equation}\label{eq:qm010-kernel-target}
\boxed{\mathcal K_{HH}^{\rm QGD}(\bm q)
=\frac{c^4}{16\pi G_N}|\bm q|^2,
\qquad\mathcal B_{H\rho}^{\rm QGD}=-\frac{c^2}{2},\qquad
\frac{\delta H_{\rm KS}(\bm q)}{\delta\rho_m(\bm q)}
=\frac{8\pi G_N}{c^2|\bm q|^2}.}
\end{equation}
This is an \emph{effective functional matching target}, not a result obtained from the contact interaction. Matching its observable consequences is essential for independent microscopic emergence of the same exterior source law, but computing a quasiparticle self-energy or a graviton pole is \emph{not} automatically the correct first research step for the original physical spinor-to-invariant map.

\subsection{Practical decision and falsifiable next step}\label{sec:qm010-decision}
\textbf{Closed [C]:} For compact spherical sources, $\hat\varrho_{\rm src}\mapsto\rho_m\mapsto H_{\rm KS}\mapsto\sigma$ follows by stationary field selection from the adopted QGD weak-field action and its boundary conditions. The exact external null-branch Chapter~2 Schwarzschild profile, compatible Dirac Clifford algebra, WKB shell and conserved-current bridge are retained. No new force or graviton calculation is required to recognize this \emph{conditional} closure.

\textbf{Open [O]:} Starting from the \emph{actual macroscopic Dirac solution and its source-state input}, derive a basis-independent physical rule $\mathcal F_{\rm micro}(\Psi,S,\mathsf R,j,\hat\varrho_{\rm src};\mathrm{BC})=[\Sigma]$ or derive the functional whose stationary points define that rule. Test source dependence, observer/probe independence, exterior matching, sector redundancy and complete-metric consistency. This is the research priority. A source-to-kinetic Dirac vertex, a tensor collective response and the prior QM001--QM009 spectral calculations remain \emph{optional diagnostic routes}, not prerequisites to maintaining the closed conditional bridge.

\textbf{Conceptual conclusion.} Nature does not select a field merely because a test ray takes a stationary path. A dynamical principle and source/boundary data select the \emph{field}; the selected field then governs its rays and spinor propagation. In effective QGD the relevant principle is already specified. The microscopic origin and uniqueness of its physical natural-invariant selection remain the one unresolved foundational connection.

\section{The decisive source-to-kinetic response target}\label{sec:target-kernel}
\subsection{What the effective QGD source law supplies [C]}
For a static, compact, nonrelativistic quantum source with independently computed mass density $\rho_m(\bm y)$, adopt the weak-field QGD/GR source response, including its empirically normalized $G_N$:
\begin{equation}\label{eq:Newton-source}
 \boxed{H[\rho_m](\bm x)=\frac{2G_N}{c^2}
 \int\frac{\rho_m(\bm y)}{|\bm x-\bm y|}\,\dd^3y,
 \qquad \nabla^2H=-\frac{8\pi G_N}{c^2}\rho_m.}
\end{equation}
For $r>R$ and a spherical source of total mass $M_s=\int\rho_m\dd^3y$, $H=r_s/r$, $r_s=2G_NM_s/c^2$. Chapter~2's independent Euler calculation, with the chosen QGD vacuum null branch, gives the same \emph{exact exterior} profile when that Einstein--Hilbert action is adopted. The interior Poisson relation here is only a weak-field approximation, not an exact general-relativistic stellar solution.

\subsection{A functional derivative that any stronger microscopic map must match [D/C]}
Keeping the specified exterior null orientation $\ell_\mu(\bm x)$ fixed while perturbing the source density, differentiate Eqs.~\eqref{eq:KS-op} and \eqref{eq:Newton-source}:
\begin{equation}\label{eq:source-vertex}
 \boxed{
 \frac{\delta\Gamma^\mu(\bm x)}{\delta\rho_m(\bm y)}
 =-\frac{G_N}{c^2}
 \frac{\ell^\mu(\bm x)\not\!\ell(\bm x)}{|\bm x-\bm y|}.}
\end{equation}
In a local patch with fixed $\ell$, the nonzero-momentum static Fourier-space target reads
\begin{equation}\label{eq:target-k2}
 \boxed{\delta\Gamma^\mu(\bm q)
 =-\frac{4\pi G_N}{c^2}\,
 \ell^\mu\not\!\ell\,
 \frac{\delta\rho_m(\bm q)}{|\bm q|^2}.}
\end{equation}
Here $\bm q$ denotes Fourier momentum (\emph{not} the null one-form $\ell$). Equations~\eqref{eq:source-vertex}--\eqref{eq:target-k2} are an \emph{exact functional derivative of the stipulated effective matching rule}. The $1/|\bm q|^2$ kernel was supplied by the adopted gravitational equation; it has \emph{not} been obtained from the tested microscopic contact interaction. This is now a quantitatively sharp benchmark for Route B. It does not prescribe how a fully nonlinear multi-sector, nonzero-$\lambda\cQ$ QGD construction should behave.

\subsection{What the original microscopic Dirac model can and cannot supply [D/C/O]}
Use Chapter~1's regulated contact theory $\mathcal L_{\rm micro}=\bar\psi(i\gamma^\mu\partial_\mu-m_0)\psi+G_4(\bar\psi\psi)^2$ in its tested homogeneous, isotropic Hartree phase. Its self-energy is a scalar mean-field mass shift; first-order local Fock exchange adds constant scalar and temporal terms. Those terms change the Dirac eigenspinors but \emph{do not} generate the particular position-dependent derivative matrices in Eq.~\eqref{eq:KS-op}. The tested stable scalar Gaussian source response has a finite static Hessian $K_0>0$ and, for a regular small-$\bm q$ continuation, a source-visible response analytic at $\bm q=0$; it therefore does not provide the $1/|\bm q|^2$ kernel of Eq.~\eqref{eq:target-k2}.

More sharply, at Hartree order in momentum space the defined effective Dirac operator takes the form $\mathcal D_H(p;\rho)=\gamma^\mu p_\mu-M[\rho]$ (up to the adopted energy origin). Thus its source derivative of the \emph{kinetic} principal symbol is exactly
\begin{equation}\label{eq:Hartree-mismatch}
 \boxed{\frac{\delta}{\delta\rho_m(y)}\left[\frac{\partial\mathcal D_H}{\partial p_\mu}\right]=\frac{\delta\gamma^\mu}{\delta\rho_m(y)}=0.}
\end{equation}
The first-order local Hartree--Fock corrections remain algebraic in momentum in the tested homogeneous phase, so the same zero derivative result holds at that order. It demonstrably disagrees with the nonzero effective target Eq.~\eqref{eq:source-vertex}. This is the first directly calculated \emph{operator-vertex mismatch} for the chosen microscopic phase: it is stronger than observing merely that a spinor ratio changes. It does not exclude a higher-order or different-phase state-dependent kinetic vertex.

This is a \emph{restricted negative result}, not an all-orders impossibility theorem for contact QFT: different phases, additional composite operators, non-Gaussian collective physics, dynamical gauge interactions or nonperturbative effects can change the infrared theory. The QM001--QM005 single-fermion dispersion calculations are likewise not by themselves evidence of the required universal source-to-kinetic response. A momentum-dependent quasiparticle self-energy need not be a universal curved-space tetrad.

\begin{proposition}[Probe and local-source obstructions, with domain stated]\label{prop:local-obstruction}
Within a prescribed static source and a fixed effective gravitational background: (i) no universal gravitational amplitude $H(\bm x)$ can be defined by equating it to the speed or lower/upper ratio of \emph{an arbitrarily chosen} test spinor, because different permitted probes at the same event have different velocities and spinor ratios; (ii) no strictly local function $H(\bm x)=F(\rho_m(\bm x),\ldots)$ that vanishes for vacuum-valued local source observables can reproduce a nonzero compact-source exterior field in a vacuum shell where those observables take their vacuum values.
\end{proposition}
\begin{proof}
For (i), at $H=1/4$ the designated free-fall speed is $\beta_{\rm ff}=1/2$ and $\zeta_{\rm ff}=1/(2+\sqrt3)\simeq0.267949$, whereas a different test particle at the same event may have a distinct speed $\beta$. Thus $H=\beta_{\rm arbitrary}^2$ cannot define a unique source geometry. For (ii), outside a compact source $\rho_m=0$ and the assumed other local source observables are vacuum-valued, giving $F=0$ by hypothesis, while Chapter~2's exterior amplitude $H=r_s/r>0$. A nonlocal, boundary-data-dependent or independent field-equation matching rule is necessary; Eq.~\eqref{eq:Newton-source} is one explicitly declared conditional example.
\end{proof}

\subsection{What counts as completion of the original QGD bridge}
The research should now use two separate completion criteria. \textbf{Effective correspondence [C]:} for a specified quantum source, recover $[\Sigma]$ by the adopted QGD Euler equations, verify its full master metric \eqref{eq:full-preserved} or declared null-sector specialization, build $\mathcal D_g$ with its spin connection, and check WKB and $r^2J^r=C_J$. This route is mathematically closed in the declared spherical exterior benchmark; it must not be advertised as microscopic genesis of the gravitational field. \textbf{Physical origin [O]:} compute a source-dependent kinetic symbol or a physical collective constitutive law without first using the target metric/action, and show that its source response produces the required exterior field and a species-independent effective propagation geometry. Its functional target is Eq.~\eqref{eq:source-vertex}; the long-range/nonlocal feature is not optional if the source is compact and the exterior field nonzero, but it need not be described as a separate fundamental force carrier.

\paragraph{Practical next microscopic calculation.}
Compute the \emph{source-to-probe derivative kinetic vertex}
\begin{equation}\label{eq:actual-target}
 \boxed{\mathcal V^{\mu}(x,y;p)
 =\frac{\delta}{\delta\rho_m(y)}
 \left[\frac{\partial\mathcal D_{\rm eff}(x,p)}{\partial p_\mu}\right]_{\rm physical\ state},}
\end{equation}
starting from a specified source density matrix (not one arbitrarily selected probe mode), with matched regulators and conserved source vertices. Its null-sector target, in the Dirac part of the symbol, is Eq.~\eqref{eq:source-vertex}. A failure of the tested phase to produce the required nonlocal state dependence would be a scientifically useful result; a successful calculation would provide the missing physics for the original state-to-natural-invariant map. This test is more targeted than searching for a massless tensor particle in isolation.


\section{Preserve the WKB and wavefunction--momentum bridge}
Within a coherent local positive-energy Dirac branch with an effective scalar mass $\mathcal M(x)$ and gauge-invariant spatial momentum $\bm\pi$, the leading local Dirac equation gives
\begin{equation}\label{eq:spinor}
\boxed{u_B=\mathsf R u_A,\qquad \mathsf R=\frac{c\bm\tau\cdot\bm\pi}{\mathcal E_++\mathcal Mc^2},\quad \bm\zeta=\frac{c\bm\pi}{\mathcal E_++\mathcal Mc^2}.}
\end{equation}
This is a \textbf{[D]} local eigenspinor relation under its stated assumptions, not yet the gravitational covector $\sig_\mu$. The full four-mode QGD construction may be written $\Psi=U_fc$; since alternative invertible $U_f$ yield the same $\Psi$ by a compensating change in $c$, the unrestrained four-mode basis does not uniquely identify the physical gravitational profile.

For a coherent branch, $\Psi\simeq Au\exp(iS/\hbar)$ with $\bm\pi=\bm\nabla S-q\bm A$ when the eikonal and gauge conventions allow. The leading spinor current is $\bm j\simeq\rho_D c^2\bm\pi/\mathcal E_+$, while the exact global-number Ward identity of the interacting microscopic theory is $\partial_\mu j^\mu=0$. Hence in a stationary, spherically symmetric source-free shell,
\begin{equation}\label{eq:bridge}
\boxed{r^2j^r=C_J,\qquad j^r\simeq\rho_D\frac{\pi_rc^2}{\mathcal E_+}\quad\text{for an isolated positive-energy WKB stream}.}
\end{equation}
The first equation concerns the conserved \emph{total} current; a separately identified collisionally coupled stream can have a source in its own kinetic transport equation. A zero net flow of a bound or isotropic source does not mean zero energy density or zero stress. The WKB and flux equations are \textbf{[D/C]} matter-transport statements; no $1/r$ gravity law follows solely from them.

The QM001--QM005 restricted finite-density calculations demonstrate state-dependent effective spinor projectors and a self-energy/damping programme, but do not identify a universal gravitational mode. They are \emph{input benchmarks}, not independent microscopic gravity proofs.

\section{A microscopic response kernel without an inserted Einstein propagator}
Let $\widehat{\mathcal O}_I$ be a specified set of renormalized or cutoff-defined microscopic observables, including total stress-energy, currents and candidate scalar/tensor composites. With physical probe sources $J_A$, define $W[J]$, expectation coordinates $\varphi^I$, and a Legendre effective action $\Gamma_{\rm coll}[\varphi]$ on the nondegenerate constrained response subspace. At a stable stationary state,
\begin{equation}\label{eq:K}
\cK_{IJ}(k):=\frac{\delta^2\Gamma_{\rm coll}}{\delta\varphi^I(-k)\delta\varphi^J(k)},\qquad V_A{}^I(k):=\frac{\delta^2\Gamma_{\rm coll}}{\delta J_A(-k)\delta\varphi^I(k)}.
\end{equation}
The physical linear response is schematically
\begin{equation}\label{eq:chi}
\boxed{\chi^{\rm phys}_{AB}(k)=V_A{}^I(k)[\cK^{-1}(k)]_{IJ}V_B{}^J(-k)+\chi^{\rm contact}_{AB}(k).}
\end{equation}
The exact signs of the exchange term depend on source/Legendre conventions, but the pole criterion does not. $V$ and $\cK$ must be computed from the \emph{same} microscopic generating functional. In a theory with gauge or constraint directions the inverse is understood only on the physical response quotient, with the Ward seagull/contact pieces retained. Defining $\cK$ via a Legendre transform does not presuppose independent physical propagation of every coordinate $\varphi^I$.

\begin{definition}[Necessary static source-visible criterion \textbf{[C]}]
For a specified pair of physical source channels in three spatial dimensions, a genuine long-range inverse-distance linear response requires a nonzero, noncancelled coefficient of $1/|\bm k|^2$ in the \emph{full} static physical response. Define the diagnostic
\begin{equation}\label{eq:residue}
\boxed{\cR_{AB}:=\lim_{|\bm k|\to0}|\bm k|^2\chi^{\rm phys}_{AB}(0,\bm k),}
\end{equation}
when the limit and specified order of frequency/momentum limits exist. $\cR_{AB}\ne0$ is necessary for a simple Coulomb/Newton-like static pole with regular subleading terms; it is not sufficient for universally coupled relativistic gravity.
\end{definition}
The source-visible condition is stronger than ``the spectrum is gapless.'' A sound-like density mode or a particle--hole continuum need not mediate an independent static $1/r$ interaction between generic sources. One must identify the physical eigenmode, stability, causal dynamics, coupling and cancellation structure without placing an Einstein/Fierz--Pauli $1/k^2$ propagator inside $\cK$ by assumption.

\section{Explicit Chapter 1 scalar-contact test: a finite response}
Choose the original Chapter 1 model with vanishing background electromagnetic field,
\begin{equation}\label{eq:micro}
\boxed{\mathcal L=\bar\psi(i\gamma^\mu\partial_\mu-m_0)\psi+G_4(\bar\psi\psi)^2,\qquad G_4>0.}
\end{equation}
In natural units at zero temperature, the prescribed homogeneous Hartree saddle satisfies
\begin{equation}\label{eq:gap}
0=\frac{M-m_0}{2G_4}-\frac M{\pi^2}\int_{p_F}^{\Lambda}\frac{p^2\,\dd p}{\sqrt{p^2+M^2}},\qquad p_F=\sqrt{\mu^2-M^2}.
\end{equation}
This is a fixed-cutoff effective model, not a regulator-independent ultraviolet completion. Its scalar inverse static response, density-mixing vertex, and direct density response are
\begin{align}
\cK_{MM}(0,\bm0)=K_0&=\frac1{2G_4}-\frac1{\pi^2}\left[\int_{p_F}^{\Lambda}\frac{p^4\,\dd p}{(p^2+M^2)^{3/2}}+\frac{M^2p_F}{\mu}\right],\label{eq:K0}\\
V_0=\Omega_{M\mu}&=\frac{Mp_F}{\pi^2},\qquad \chi^{\rm direct}_{nn}=\frac{\mu p_F}{\pi^2}.\label{eq:V0}
\end{align}
For $\Lambda=4$, $m_0=0.12$, $G_4=0.5654255716$, and $\mu=1.2$, independent quadrature and root-solving give
\begin{equation}\label{eq:bench}
\boxed{\begin{gathered}
M=0.693435535520,\qquad p_F=0.979360586341,\\
K_0=0.168741734136>0,\qquad V_0=0.068809590036,\\
\chi^{\rm direct}_{nn}=0.119075968585,\quad \frac{V_0^2}{K_0}=0.028059209568,\\
\chi_{nn}^{\rm total}=0.147135178153.
\end{gathered}}
\end{equation}
The total susceptibility was independently checked by differentiating the density after re-solving the nonlinear Hartree gap equation at nearby chemical potentials. This numerical reproduction agrees with Chapter 5 v3.7 and does not independently infer an exterior gravitational potential.

\begin{proposition}[Regular scalar channel does not generate the required static pole \textbf{[C]}]\label{prop:nopole}
Assume the scalar Hessian and density vertex are analytic near zero static momentum, with $\cK_{MM}(0,\bm k)=K_0+O(|\bm k|^2)$, $V(0,\bm k)=V_0+O(|\bm k|^2)$, and direct/contact terms regular there. If $K_0>0$, then
\begin{equation}\label{eq:nopole}
\boxed{\lim_{|\bm k|\to0}|\bm k|^2\chi_{nn}^{\rm phys}(0,\bm k)=0.}
\end{equation}
Therefore the tested stable scalar-density channel cannot realize the simple source-visible $1/|\bm k|^2$ static MW2 response under these hypotheses.
\end{proposition}
\begin{proof}
Since $K_0>0$, the inverse scalar Hessian is bounded in a neighborhood of zero momentum and $V\cK^{-1}V=V_0^2/K_0+O(|\bm k|^2)$. The regular direct/contact contribution is bounded too. Multiplication by $|\bm k|^2$ gives zero. The result is restricted to the named channel and analytic, stable phase; it is not a statement about all tensor or exotic collective excitations of QFT.
\end{proof}
For visual illustration only, insert a \emph{trial} positive static gradient coefficient $Z_s$ and set $x=|\bm k|\sqrt{Z_s/K_0}$. Then $\chi_{\rm exch}=V_0^2/[K_0(1+x^2)]$: its static limit remains finite while $x^2\chi_{\rm exch}\to0$. The stiffness $Z_s$ is \emph{not} claimed computed by this plotted model.
\begin{figure}[htbp]\centering
\includegraphics[width=.88\linewidth]{static_scalar_pole_diagnostic.pdf}
\caption{Reproducible scalar pole diagnostic. The positive gradient coefficient is a declared illustrative regular continuation; the $K_0,V_0$ inputs are calculated from the named microscopic Hartree saddle. No gravitational response has been inserted.}\label{fig:nopole}
\end{figure}

\section{Second obstruction: scalar matter channels do not become tensor gravity by relabeling}
Consider an isotropic equilibrium state and a single rotational scalar excitation $q$. For spatial external momentum $\bm k$, the most general parity-even symmetric scalar--stress two-point vertex has the tensor form
\begin{equation}\label{eq:rotvertex}
V_{ij,q}(0,\bm k)=a(|\bm k|^2)\delta_{ij}+b(|\bm k|^2)k_i k_j,
\end{equation}
up to contact terms with the same rotational structure.
\begin{proposition}[Vanishing transverse-traceless scalar mixing \textbf{[D/C]}]\label{prop:TT}
For either physical spatial TT polarization $e_{ij}^{\rm TT}(\bm k)$ satisfying $e_{ii}^{\rm TT}=0$ and $k_ie_{ij}^{\rm TT}=0$, one has
\begin{equation}\label{eq:TTzero}
\boxed{e_{ij}^{\rm TT}V_{ij,q}=0.}
\end{equation}
Thus an isotropic scalar-only linear exchange does not supply a helicity-two pole through its mixed stress vertex. Additional tensor, vector, anisotropic, symmetry-breaking or nonperturbative channels require their own independent tests.
\end{proposition}
\begin{proof}
Contract Eq.~\eqref{eq:rotvertex} with the TT tensor. The trace term vanishes because $e_{ii}=0$ and the momentum term because $k_ie_{ij}=0$.
\end{proof}
In a Lorentz-invariant stationary vacuum, after satisfying the complete stress Ward identity including contacts, a regular scalar--metric vertex takes the compatible form $V_{\mu\nu,q}\propto(k^2\eta_{\mu\nu}-k_\mu k_\nu)F(k^2)$ in the tested scalar channel. If $F$ is finite and $\cK_{qq}(0)>0$, exchange contributions generated solely by this vertex scale as $O(k^4)$, not as $1/k^2$. This complements Chapter 5's explicitly calculated finite-momentum stress--mass form factor. A medium can have different temporal/spatial form factors, but the isotropic spatial TT contraction above remains zero for a scalar two-point vertex.

\section{Third obstruction: a strictly local source-to-field rule cannot explain the exterior}
\begin{proposition}[Exterior obstruction for a vanishing local matter order parameter \textbf{[D/C]}]\label{prop:exterior}
Let a compact stationary matter source have coarse-grained matter data $\varphi^I(x)$ equal to fixed vacuum values $\varphi^I_{\rm vac}$ in an open exterior region $U$. Suppose a proposed natural-invariant prescription is strictly local and algebraic, $\Sigma(x)=F(\varphi^I(x))$, with no derivatives, boundary response, independent field degrees of freedom, or state-dependent nonlocal kernel. Then $\Sigma(x)=F(\varphi_{\rm vac})$ throughout $U$. It cannot represent a \emph{source-dependent} nonconstant Schwarzschild exterior in that region unless the vacuum data or map itself carry nonlocal source information.
\end{proposition}
\begin{proof}
At every exterior point the arguments of $F$ equal the same vacuum constants, so its output is constant in the chosen fixed reference frame. A universal map has no local variable that encodes source mass or radial distance there. A source-dependent exterior therefore requires additional nonlocal response, boundary-value dynamics or genuinely exterior quantum correlations.
\end{proof}
The statement does not exclude vacuum correlations extending outside a compact classical source or a collective field living in the exterior. Rather, it identifies exactly what a successful microscopic realization must compute. A local Dirac momentum ratio is not automatically such an exterior response.


\section{Incremental tensor calculation: collisionless stress spectrum and kinetic Ward completion (QGD--QM006)}\label{sec:qm006-tensor}
\noindent\textbf{Scope and status [D/C].} This first new tensor increment is evaluated in the homogeneous, zero-temperature Hartree Fermi sea of Chapter~1, retaining only its leading \emph{positive-energy intraband} particle--hole response. It is an independently reproducible refinement of that chapter's shear formula, not the complete renormalized metric-source polarization. Vacuum/antiparticle bands, full stress vertices, interaction-resummed tensor ladders, and the gravitational metric seagull remain open and are not inferred from this result.

\subsection{Derivation of the transverse-traceless Fermi-surface response}
Select external wavevector $\bm q=q\bm e_z$. For a local positive-energy on-shell Dirac quasiparticle, the leading kinetic shear insertion is $V_{xy}(\bm p)=p_xp_y/E_{\bm p}$ and the Fermi occupation $n_{\bm p}=\Theta(p_F-|\bm p|)$. With spin degeneracy two, define the retarded intraband response by
\begin{equation}\label{eq:qm006-bubble}
 \chi^R_{xy,xy}(\omega,q)=2\int\frac{\dd^3p}{(2\pi)^3}\,\frac{n_{\bm p-\bm q/2}-n_{\bm p+\bm q/2}}{\omega+i0^++E_{\bm p-\bm q/2}-E_{\bm p+\bm q/2}}\,V_{xy}(\bm p)^2+O(q/p_F)_{\rm vertex}.
\end{equation}
Taking $q\to0$ with $z=(\omega+i0^+)/(v_Fq)$ fixed, using $n_- - n_+=qv_Fu\,\delta(E_p-\mu)+O(q^3)$ and $u=\cos\theta$ gives
\begin{equation}\label{eq:qm006-J}
 \boxed{\begin{aligned}
 \chi^R_{xy,xy}(\omega,q)&=C_TJ(z)+O(q/p_F),\qquad C_T=\frac{p_F^5}{16\pi^2\mu},\\
 J(z)&=\int_{-1}^{1}\frac{u(1-u^2)^2}{z-u}\,\dd u\\
 &=-2z^4+\frac{10}{3}z^2-\frac{16}{15}+z(1-z^2)^2\log\frac{z+1}{z-1}.
 \end{aligned}}
\end{equation}
The logarithm has the retarded branch. Its static limit and its continuum discontinuity are
\begin{equation}\label{eq:qm006-static-cut}
 \boxed{\begin{aligned}
 \chi^R_{xy,xy}(0,q\to0)&=-\frac{p_F^5}{15\pi^2\mu}=-0.005071551738,\\
 \operatorname{Im}\chi^R_{xy,xy}(\omega,q)&=-\pi C_T\,z(1-z^2)^2\,\Theta(1-|z|)\quad(z\in\mathbb R),
 \end{aligned}}
\end{equation}
where the second expression assumes the low-$q$ intraband limit. At $z=0.5$, the model gives
\[\chi^R_{xy,xy}=-0.0002346320-0.0042010180\,i.\] These are finite continuum responses; they do not identify an isolated $1/q^2$ pole or a physical massless spin-two particle. An additive analytic metric-source contact cannot generate such a nonlocal pole by itself.

\noindent\textbf{Both helicity-two components.} For the orthogonal shear $V_D=(p_x^2-p_y^2)/(2E)$ the two azimuthal averages coincide:
\[\int_0^{2\pi} (\cos\phi\sin\phi)^2\,\dd\phi=\int_0^{2\pi}\tfrac14\cos^2(2\phi)\,\dd\phi=\frac\pi4.\]
Their mixed angular integral vanishes. Thus $\chi^R_{D,D}=\chi^R_{xy,xy}$ and $\chi^R_{D,xy}=0$ at this isotropic intraband order. An isotropic scalar $\bar\psi\psi$ has zero mixed susceptibility to either TT component at the same order.

\begin{center}\includegraphics[width=.87\linewidth]{tensor_shear_continuum.png}\\[-1pt]
{\footnotesize\itshape Leading positive-frequency retarded intraband tensor response. The imaginary part is the particle--hole continuum, not an isolated graviton.}
\end{center}

\subsection{Transverse-momentum Ward identity and the indispensable kinetic contact}
The TT stress is not itself longitudinally constrained by momentum conservation. To independently check conservation, use $T^{0x}$ and $T^{zx}$, with on-shell insertions $V_{0x}=p_x$ and $V_{zx}=p_zp_x/E_p$. At leading Fermi-surface order, their correlators obey
\begin{equation}\label{eq:qm006-kinetic-Ward}
 \boxed{\omega\chi^R_{0x,zx}(\omega,q)-q\big[\chi^R_{zx,zx}(\omega,q)+C_{\rm kin}\big]=0,\qquad C_{\rm kin}=\frac{p_F^5}{15\pi^2\mu}.}
\end{equation}
To verify, substitute the two vertices into Eq.~\eqref{eq:qm006-bubble}. In the combination on the left, the factor $(z-u)$ cancels the denominator, leaving the angular integral $\int_{-1}^1u^2(1-u^2)\dd u=4/15$. The resulting nonzero local term is exactly $qC_{\rm kin}$ and must be retained to satisfy the leading kinetic momentum Ward relation. Numerically, at $z=1.5$ we obtain $\chi_{0x,zx}=0.005207179181$, $\chi_{zx,zx}=0.001303080831$, and a normalized Ward residual below $2\times10^{-18}$.

\noindent\textbf{Crucial qualification.} $C_{\rm kin}$ is the contact needed to complete this \emph{linearized collisionless kinetic} stress identity. It is \emph{not yet} the full local metric seagull/contact obtained by twice varying the regulated quantum effective action. A regulator-consistent microscopic gravity-source Ward identity must include that genuine seagull, vacuum terms, exchange and all conserved interaction vertices; the present identity is a restricted nontrivial consistency test, not a proof of complete diffeomorphism covariance.

\subsection{Independent numerical checks and next increment}
A second integration with the finite spatial-momentum occupation boundaries recovers $\chi_{xy,xy}(z=1.5,q\to0)=0.000382522300$. At $q=0.005$ it gives $0.000382519086$, a relative deviation $8.40\times10^{-6}$. Independent Cauchy-principal-value integration reproduces Eq.~\eqref{eq:qm006-J} inside the Landau cut to better than $4\times10^{-11}$ at the sampled $z$ values. The code and data are packaged as \texttt{verify\_tensor.py} and \texttt{tensor\_shear\_continuum.png}.

\noindent\textbf{Next narrowly bounded increment [O].} Calculate the actual \emph{metric-source} $T^{\mu\nu}$ bubble and second-variation seagull from the same sharp-cutoff Dirac effective action; evaluate its full stress Ward residual before and after contact terms. Only after that check should the interaction-resummed tensor Hessian be tested for a source-visible isolated pole. The WKB/Dirac ratio and stationary $r^2j^r=C_J$ remain independent upstream matter identities.


% ================== QGD QM007 2026-11-02 ==================
\section{Covariant stress Hessian, genuine metric seagull and tensor-pole decision (QM007)}\label{sec:qm007}
\noindent\textbf{Scope and scheme [D/C].} The earlier scalar benchmark employs a sharp \emph{three-momentum} regulator in the preferred medium rest frame. Such a cutoff must not silently be treated as a covariant geometric regulator. In this section the vacuum determinant is separately defined by a covariant proper-time cutoff; the occupied positive-energy states contribute a finite density correction, with the chemical-potential one-form transformed under background coordinate changes. When matching the same homogeneous mass, the contact coupling is \emph{refitted} and the two regulated theories are not identified coefficient by coefficient. The complete finite-$q$ retarded Dirac Hessian, including all interaction-consistent vertices, has not been evaluated here. The exact functional bubble--seagull construction and the local covariant part have been evaluated and checked explicitly.

\subsection{Metric-source definition and exact second-variation identity}
Put the scalar Hubbard--Stratonovich field at a specified saddle $M$ and couple the Dirac field to a \emph{bookkeeping external metric} $g$. In Euclidean signature, let $\mathcal D_g=\gamma^a e_a{}^{\mu}\nabla_\mu[g]+M$, with the required spin connection. Define the positive Laplace-type operator
\begin{equation}\label{eq:qm007-H}
 \boxed{H_g=\mathcal D_g^\dagger\mathcal D_g=-\nabla_g^2+\frac14 R[g]+M^2,\qquad
 W_{\Lambda}[g;M]=\frac12\int_{\Lambda^{-2}}^\infty\frac{\dd s}{s}\Tr\,e^{-sH_g}.}
\end{equation}
Here the equality to the minimal Laplace-type form refers to the torsion-free, zero-gauge-field, constant-$M$ vacuum benchmark; nonconstant $M$ or gauge fields produce additional operator terms. The path-integral fermionic sign and the first-order determinant are incorporated in the stated parity-even Euclidean proper-time convention. $W_\Lambda$ is a \emph{matter response to an external metric}, not a dynamical graviton action. The independent HS contribution $(M-m_0)^2/(4G_4)\int\sqrt g$ must be included in the full mean-field stress and varied consistently when $M$ responds.

For independent metric perturbations $h$ and $k$, denote $H_h=\delta_h H$, $H_{hk}=\delta_h\delta_k H$. Differentiating the \emph{regulated trace itself} yields the exact identity
\begin{equation}\label{eq:qm007-seagull}
 \boxed{\begin{aligned}
 \delta_h\delta_k W_\Lambda={}&\frac12\int_{\Lambda^{-2}}^\infty\dd s\int_0^s\dd\tau\,
 \Tr\!\left(e^{-(s-\tau)H}H_k e^{-\tau H}H_h\right)\\
 &-\frac12\int_{\Lambda^{-2}}^\infty\dd s\,\Tr\!\left(e^{-sH}H_{hk}\right).
 \end{aligned}}
\end{equation}
The first term is the regulated \emph{two-linear-vertex bubble}; the second is the \emph{genuine metric-source seagull}, including second variations of the coframe, spin connection and density weight implicit in the covariant trace. Crucially, neither term is defined by adding an arbitrary local tensor after integration. A finite-dimensional regularization or a covariant bundle trace can make the cyclic functional-derivative argument rigorous. Equation~\eqref{eq:qm007-seagull} is exact for the chosen proper-time functional; the actual finite-$q$ numerical integrals are a separate task.

\subsection{Exact metric Ward identity, including contact and tadpole terms}
Define $T^{\mu\nu}=2(\sqrt g)^{-1}\delta W/\delta g_{\mu\nu}$ and $K^{\mu\nu,\rho\sigma}=\delta T^{\mu\nu}/\delta g_{\rho\sigma}$ at a homogeneous flat background, with symmetric-source differentiation conventions. Varying $\nabla_\mu T^{\mu\nu}=0$ gives, in Fourier space,
\begin{equation}\label{eq:qm007-Ward}
\boxed{\begin{aligned}
 q_\mu K^{\mu\nu,\rho\sigma}(q)={}&-\frac12 q_\lambda\overline T^{\lambda\nu}\eta^{\rho\sigma}
 -\frac12\big(q^\rho\overline T^{\sigma\nu}+q^\sigma\overline T^{\rho\nu}\big)\\
 &+\frac12q^\nu\overline T^{\rho\sigma}.
\end{aligned}}
\end{equation}
The expectation value $\overline T$ is the \emph{complete} background stress and all metric-coupled sources are transformed. At nonzero pressure or vacuum tadpole, the Hessian need not be simply transverse. A purported check of $q_\mu\langle TT\rangle=0$ before including the contacts is generally incorrect. For $q$ along $z$ and spatial shear indices $(\nu;\rho\sigma)=(x;zx)$, Eq.~\eqref{eq:qm007-Ward} becomes $K^{zx,zx}=-\overline P/2$ in the homogeneous static local pressure response convention. The occupation-only free-quasiparticle diagnostic has $n=0.03172538246$, $\varepsilon_{\rm occ}=0.03235666932$, $P_{\rm occ}=0.005713789638$, hence $-P_{\rm occ}/2=-0.002856894819$. This is a \emph{kinetic subset illustration}; the total interacting $\overline P$ also includes the HS potential, sea subtraction and chosen finite counterterms.

\begin{proposition}[Nonlinear-coordinate Ward/seagull cancellation \textbf{[D]}]\label{prop:qm007-ward-shear}
On a periodic flat Euclidean coordinate patch, choose the finite coordinate change $x'=x+\epsilon\sin(qz)/q$ and leave $z$ and the other coordinates unchanged. Its pulled-back metric has
\begin{equation}\label{eq:qm007-pullback}
 g_{xz}=\epsilon\cos(qz),\quad g_{zz}=1+\epsilon^2\cos^2(qz),\quad
 \det g=1,\quad R[g]=0.
\end{equation}
Therefore $W_\Lambda[g]$ is identical to its flat value under this diffeomorphism when all source data and boundary conditions are transformed. For a local $C_0\int\sqrt g$ term, the order-$\epsilon^2$ quadratic Hessian from the first-order shear is $-C_0\epsilon^2/4$ per unit coordinate volume, while the second-order metric displacement times the one-point tadpole is $+C_0\epsilon^2/4$. The sum vanishes exactly. The seagull/tadpole terms cannot be dropped or replaced by the transverse kinetic contact of Section~\ref{sec:qm006-tensor}.
\end{proposition}
\noindent Independent numerical sampling gives a maximum determinant-volume deviation of $1.2\times10^{-16}$ and zero cancellation residual at the displayed precision. This is an explicit nonlinear coordinate Ward check of the local covariant functional, \emph{not} a numerical evaluation of every nonlocal stress-bubble component.

\subsection{Vacuum Dirac determinant: calculated local terms and tensor power counting}
For $q^2\ll M^2$ and a constant mass, the heat-kernel coefficients of the operator in Eq.~\eqref{eq:qm007-H} are $\operatorname{tr}a_0=4$ and $\operatorname{tr}a_1=4(R/6-R/4)=-R/3$. Thus
\begin{equation}\label{eq:qm007-vac}
\boxed{W_\Lambda[g;M]=\int\dd^4x\sqrt g\big[C_0+C_RR+O(R^2,\nabla^2R)\big],\quad
 C_0=\frac{I_0}{8\pi^2},\quad C_R=-\frac{I_1}{96\pi^2},}
\end{equation}
where $I_0=\int_{\Lambda^{-2}}^\infty\dd s\,s^{-3}e^{-sM^2}$ and $I_1=\int_{\Lambda^{-2}}^\infty\dd s\,s^{-2}e^{-sM^2}$. For $M=0.693435535520$, $\Lambda=4$ in this \emph{proper-time} scheme,
\begin{equation}\label{eq:qm007-vacnumbers}
\boxed{I_0=120.81937817,\quad I_1=14.10422862,\quad C_0=1.53019530,\quad C_R=-0.01488601186.}
\end{equation}
All four coefficients are scheme-dependent; an Einstein--Hilbert term with a fixed measured Newton constant has \emph{not} been microscopically predicted. In particular, background metric dependence does not turn the metric into an integrated physical quantum degree of freedom.

For a transverse-traceless Euclidean shear $h_{xy}=\epsilon\cos(qz)$, its quadratic two-derivative action after cancelling the flat vacuum tadpole is
\begin{equation}\label{eq:qm007-TTq2}
\boxed{\frac{\Delta W_{{\rm vac},TT}^{(2)}}{V}
 =-\frac{C_R}{4}\epsilon^2q^2+O(\epsilon^2q^4,\epsilon^3)
 =0.003721502965\,\epsilon^2q^2+\cdots.}
\end{equation}
This follows from the standard quadratic expansion $\int\sqrt g R=-\frac14\int(\partial h_{ij}^{\rm TT})^2$ in the stated Euclidean convention. Its $q^2$ dependence is the induced \emph{inverse-kernel/metric-source action}, not an already observed $1/q^2$ stress susceptibility. Turning it into a massless propagator would require a physically justified, integrated metric or collective field and its normalized kinetic term. A source-only Dirac determinant does \emph{not} provide this by itself.

\subsection{Regulator matching: the same physical Hartree mass, different coupling}
The earlier sharp spatial regulator $\Lambda=4$ had $G_4^{(3D)}=0.5654255716$. In the covariant proper-time vacuum prescription and finite occupied-medium contribution, the homogeneous saddle equation is
\begin{equation}\label{eq:qm007-gap}
0=\frac{M-m_0}{2G_4^{(PT)}}-\frac{M I_1(M)}{4\pi^2}
+\frac{M}{\pi^2}\int_0^{p_F}\frac{p^2\dd p}{\sqrt{p^2+M^2}},\qquad p_F=\sqrt{\mu^2-M^2}.
\end{equation}
Matching $M=0.693435535520$ at $m_0=0.12,\,\mu=1.2$ now requires
\begin{equation}\label{eq:qm007-Gpt}
\boxed{G_4^{(PT)}=1.2697985402\ne G_4^{(3D)}.}
\end{equation}
The scalar Hessian and source vertex of this \emph{separately matched} proper-time scheme evaluate to
\begin{equation}\label{eq:qm007-PTresponse}
\boxed{K_0^{(PT)}=0.08439298667>0,\quad V_0=0.06880959004,\quad
 \chi_{nn}^{(PT)}(0,0)=0.17517967895.}
\end{equation}
Two independent finite differences of the gap equation reproduce $K_0^{(PT)}$ and $V_0$ within $5\times10^{-10}$. Neither scheme has a scalar static $1/q^2$ pole on this stable regular branch. The different numerical susceptibility illustrates \emph{regulator dependence}, not a physical change of gravitational strength.

\subsection{Full ten-channel leading Fermi-surface stress response and scalar-resummed TT result}
For spin degeneracy two at $T=0$, write $p^\mu=(\mu,\bm p)$ on the Fermi surface and define the ten kinetic source vertices $v^{\mu\nu}(\widehat{\bm p})=p^\mu p^\nu/\mu$. The complete \emph{leading intraband} symmetric response is the $10\times10$ matrix
\begin{equation}\label{eq:qm007-matrix}
 \boxed{\chi^{\rm intra}_{AB}(\omega,q)
 =\frac{\mu p_F}{\pi^2}\int\frac{\dd\Omega_{\widehat p}}{4\pi}
 \frac{v_F q\widehat p_z}{\omega-v_Fq\widehat p_z+i0^+}
 v_A(\widehat p)v_B(\widehat p),\quad A,B=(\mu\nu).}
\end{equation}
At strictly static $\omega=0$ this is $\chi^{\rm intra}(0,q\to0)=-\mathsf A$, where $\mathsf A$ is a positive-semidefinite angular Gram matrix. Direct independent Gaussian angular quadrature gives eigenvalues (multiplicity indicated)
\begin{equation}\label{eq:qm007-eigs}
\boxed{0\ (1),\ \ 0.005071551738\ (3),\ \ 0.01014310348\ (2),\quad
0.03807045895\ (3),\ \ 0.19682715345\ (1).}
\end{equation}
The rank is nine. Its null direction is the exact on-shell vertex identity
$\sum_i v^{ii}=(p_F^2/\mu^2)v^{00}$; it is not a propagating tensor gauge degree of freedom. The numerical null residual divided by the matrix norm is below $10^{-14}$. In addition, all 40 leading collisionless longitudinal kinetic identities (four momentum-conservation components tested against ten symmetric-stress sources) have a maximum numerical residual $5.76\times10^{-16}$ at $z=1.5$. These are distinct from the genuine metric-source seagull Ward check above. The two independent helicity-two shear combinations remain degenerate, and the static correlator has $\chi_{xy,xy}=-0.005071551738$ in this convention. At $z=\omega/(v_Fq)=0.5$ the analytic tensor cut has $\operatorname{Im}\chi_{xy,xy}=-0.004201017969$. The response is a continuum with a finite static limit, not a tensor $1/q^2$ exchange pole.

For the \emph{same isotropic scalar-contact Hartree/HS Gaussian approximation}, $S=\bar\psi\psi$ is helicity zero, whereas $T_{xy}$ for $\bm q\parallel\widehat z$ is helicity two. Therefore $\chi^0_{T_{xy},S}=0$ by rotational symmetry. The exact Gaussian Schur complement consequently gives
\begin{equation}\label{eq:qm007-schur}
\boxed{\Gamma_{hh,TT}^{\rm eff}
 =\Gamma_{hh,TT}-\Gamma_{hM,TT}\,K_s^{-1}\Gamma_{Mh,TT}
 =\Gamma_{hh,TT},\qquad \Gamma_{hM,TT}=0.}
\end{equation}
A vanishing scalar--TT vertex cannot generate the missing tensor collective mode by resumming only the existing scalar auxiliary field. Combined with a finite scalar $K_0$ and the massive vacuum threshold, this establishes a \emph{phase-specific negative result for the tested scalar-Hartree Gaussian tensor channel}, not a no-go theorem for all interactions, non-Gaussian exchange, critical phases, gauge sectors or other quantum matter. This is a more direct constraint on QGD's Chapter 5 premise MW2 than quasiparticle lifetime estimates alone.

\begin{figure}[htbp]
\centering
\includegraphics[width=.48\linewidth]{stress_rank9.png}\hfill
\includegraphics[width=.48\linewidth]{tensor_continuum.png}
\caption{Independent numerical evidence in the leading finite-density collisionless sector: the symmetric-stress matrix has one on-shell null relation and nine finite eigenvalues (left), while the helicity-two retarded shear function has a particle--hole cut rather than an isolated pole (right). These are matter responses, not derived gravitons.}
\label{fig:qm007-stress}
\end{figure}
\begin{figure}[htbp]
\centering
\includegraphics[width=.71\linewidth]{induced_TT_q2.png}
\caption{The local covariant vacuum Dirac determinant after flat-tadpole subtraction produces a regular induced transverse-traceless $q^2$ quadratic action (fixed small test amplitude), not a demonstrated physical $1/q^2$ pole. The displayed range is within $q<M$; higher derivative corrections are not included.}
\label{fig:qm007-induced}
\end{figure}

\subsection{Why the bridge still matters, and the actual QGD stopping point}
The microscopic branch may still satisfy the selected Dirac ratio
$u_B=\mathsf R u_A$ and a controlled WKB transport equation, with $r^2j^r=C_J$ for the stationary total current. Those are already validated \emph{matter-side} statements. The complete founding metric
$g_{\mu\nu}=T^a{}_{\mu}T^b{}_{\nu}[\eta_{ab}-\sum_s\varepsilon_s\sigma_a^{(s)}\sigma_b^{(s)}-\lambda\mathcal Q_{ab}]$
remains an explicit \emph{physical promotion and metric postulate}, not a consequence of the scalar-pole test. A state-to-$[\Sigma]$ map requires a stable source-visible collective sector, universality across different kinds of matter, and a causal common metric for probes. No such sector has been obtained by merely generating the covariant $C_RR$ term in an external-source determinant.

\paragraph{Next decisive scientific move [O].} Extend the same generating functional to the interaction-complete \emph{retarded} stress matrix with proper seagulls and dynamical Fock/crossed vertices (or an explicitly specified conserving nonperturbative scheme), and test helicity-two poles independently of the scalar HS field. If the original scalar-contact phase continues to fail, investigate a physically motivated distinct microscopic phase or operator channel rather than inserting $Z(\nabla q)^2$ or Einstein's action as the desired answer. The existence and source coupling of $q$ are the actual Chapter 5--6 missing links.
% ================== END QM007 ==================

\section{The required physical operator channels and their tests}
A subsequent, more general source matrix should not be built from $\bar\psi\psi$ alone. In the Chapter 1 matter model, consider a renormalized, regulator-consistent operator family containing $T^{\mu\nu}_{\rm total}$, the conserved vector current, the scalar auxiliary channel, and if justified independent tensor composites. The full kernel is a matrix, not an assumed Einstein propagator:
\begin{equation}\label{eq:matrix}
\cK(k)=\begin{pmatrix}\cK_{ss}&\cK_{sT}&\cdots\\ \cK_{Ts}&\cK_{TT}&\cdots\\ \vdots&\vdots&\ddots\end{pmatrix},\qquad
\chi^{\rm phys}_{AB}=V_A\cK^{-1}V_B^{\!*}+\chi_{AB}^{\rm contact}.
\end{equation}
One must calculate the \emph{retarded} frequency and spatial-momentum dependence in a phase with explicitly stated occupation, vacuum subtraction and cutoff. The physical stress response has multiple helicity channels; direct/seagull pieces are retained before projecting onto spin-zero and helicity-two structures. The prior QM001--QM005 single-fermion residues are helpful input to dressed fermion lines and vertices, but they do not by themselves give an induced collective massless pole.

\begin{enumerate}[label=\textbf{R\arabic*.},leftmargin=1.8em,itemsep=3pt]
\item \textbf{Microscopic closure.} Specify the approximation (e.g., conserving large-$N_f$ or a regulated ladder), exact state, counterterms, all source vertices and the physical quotient. Do not identify constituent count $N$, flavor count $N_f$, and signed metric-sector count.
\item \textbf{Source-visible infrared spectrum.} Compute the full physical response and test $\cR_{AB}\ne0$ for energy and stress probes, including rest matter, kinetic pressure and radiation. Check whether any candidate pole is isolated, stable and not a gauge artifact or cancelled by contacts.
\item \textbf{Universality.} Repeat with at least two genuinely different probes and verify whether one common response normalizes their couplings to full energy--momentum, not only particle density or charge.
\item \textbf{Operational geometry.} Compute the highest-derivative probe propagation operators from the same state, and test for a shared Lorentzian principal symbol and compatible mass/redshift response. An acoustic metric of one medium is insufficient by itself.
\item \textbf{State-to-natural-invariant map.} After a physical geometry/response is identified, construct a frame-covariant state-dependent class $[\Sigma]$ and prove invariance under microscopic spin-basis changes and any genuine signed-sector equivalence. No arbitrary Gram factorization may count as independent selection.
\item \textbf{Relativistic effective action.} Only then test whether the induced action has an acceptable two-derivative Einsteinian sector, correct signs and physical radiation spectrum; the adopted GR action cannot be reused as a microscopic input to its own derivation.
\end{enumerate}

The proposed route remains open to non-Dirac quantum matter. What counts is the existence of the correct physical collective response and its universal operational consequences, not insisting that the fundamental excitation must be a particular microscopic spinor or a postulated elementary graviton.

\section{Essential theoretical constraints, not replacements for calculation}
Soft massless spin-two scattering in a Lorentz-invariant relativistic theory imposes stringent universal-coupling consistency conditions; see Weinberg (1964), \url{https://doi.org/10.1103/PhysRev.135.B1049}. The Weinberg--Witten theorem has important assumptions about a Lorentz-covariant conserved stress-energy tensor and matrix elements for massless higher-spin asymptotic states; see \url{https://doi.org/10.1016/0370-2693(80)90212-9}. These constraints must be checked rather than declaring either that no composite mechanism can ever work or that an arbitrary scalar condensate evades them. Composite-gravity models with nonstandard microscopic constraint structures have been explicitly studied by Carone, Erlich and Vaman (2017), \url{https://doi.org/10.1007/JHEP03(2017)134}, and in a fermionic version \url{https://arxiv.org/abs/1812.08201}. These are comparison examples, not a derivation from QGD's scalar-contact Dirac model.

\section{Full founding QGD metric and the explicit stopping point}
The complete founding QGD master metric is retained without redefinition:
\begin{equation}\label{eq:master}
\boxed{g_{\mu\nu}=T^a{}_{\mu}T^b{}_{\nu}\left[\eta_{ab}-\sum_{s=1}^{N_\sigma}\varepsilon_s\sig_a^{(s)}\sig_b^{(s)}-\lambda\cQ_{ab}\right],\quad s_s=-\varepsilon_s.}
\end{equation}
Here $N_\sigma$ labels signed metric sectors and does not count microscopic particles or fermion flavors. The independent operator meaning and off-shell variation of $\cQ_{ab}$ remain open. The simplified signed-sector Gram map is only a declared $\lambda\cQ=0$ classical specialization, not the complete founding metric.

For a full metric-dependent effective action, the recognizable Option C equation remains
\begin{equation}\label{eq:OptionC}
\boxed{2s_s\Pi^{ab}\sig_b^{(s)}=\lambda\big[(D_s\cQ)^*_{\sqrt{-g}}\Pi\big]^a.}
\end{equation}
This exact formal sector Euler equation under the stated metric-map hypotheses is not a proof that the microscopic QFT generates the metric-only Einstein action. Chapter 2's independently varied spherical $A,B$ fields dynamically select $A=B$ and $A^2=B^2=2G_NM/(c^2r)$ \emph{within the adopted Einstein--Hilbert sector}; Chapter 1's WKB ratios have not by themselves selected this gravitational profile.

\section{An inhomogeneous quantum-source test of the missing kinetic bridge (QM009)}\label{sec:qm009}
\textbf{Research question.} Given the Chapter~1 local scalar-contact theory, can a spatially varying source state create the source-to-\emph{kinetic} response required by the independent-microscopic route of Section~\ref{sec:physicalF}? This is a stronger requirement than changing a quasiparticle's Dirac mass or lower/upper spinor ratio. The test below never inserts a precomputed Schwarzschild field into the microscopic Hamiltonian.

\subsection{Source, observable and the exact Hartree--Fock principal-symbol statement}
Work in the flat reference frame, with $c=\hbar=1$, and let a static number-source $\delta\mu(\bm x)$ couple as $-\delta\mu\,\psi^\dagger\psi$ in the single-particle Hamiltonian. This is a response probe, \emph{not} a gravitational field. The microscopic matter interaction is $G_4(\bar\psi\psi)^2$, and the local Hartree inverse propagator in Wigner variables is
\begin{equation}\label{eq:qm009-dirac}
\mathcal D_H(x,p)=\gamma^\mu p_\mu-M(x),\qquad
 \frac{\partial\mathcal D_H}{\partial p_\mu}=\gamma^\mu.
\end{equation}
The local four-fermion Hartree and Fock self-energies depend on the coincident state density matrix $S(x,x)$; at first interaction order their coordinate-space kernels are proportional to $\delta^{(4)}(x-y)$ without derivatives. This remains true when $S(x,x)$ varies, and when the source has a nonzero spin polarization or current. Consequently
\begin{equation}\label{eq:qm009-zero}
\boxed{\left.\frac{\delta}{\delta\mu(y)}\left[\frac{\partial\mathcal D_{\rm HF}(x,p)}{\partial p_\mu}\right]\right|_{\rm local\ HF}=0.}
\end{equation}
This is an exact statement \emph{within the specified local Hartree--Fock truncation}; it is not a theorem for higher-loop, nonlocal self-energies, alternative collective phases or the full QGD master metric. In particular a source-induced scalar mass can bend massive test-particle trajectories without changing the characteristic massless cone of the principal Dirac symbol. Its vanishing does not imply the spinor ratio is unchanged.

\subsection{Finite-wave-number scalar response with two independently matched regulators}
Define the static quadratic grand-potential Hessian at a homogeneous self-consistent state by
\begin{equation}\label{eq:qm009-quadratic}
\delta^2\Omega=\frac12\int_{\bm q}\left[K_s(q)|\delta M_q|^2
+2V_s(q)\delta M_{-q}\delta\mu_q
-D_{nn}(q)|\delta\mu_q|^2\right],
\quad \int_{\bm q}=\int\frac{\dd^3q}{(2\pi)^3}.
\end{equation}
With $H_M(\bm p)=\bm\alpha\cdot\bm p+\beta M$, $E_p=\sqrt{p^2+M^2}$, $P_s(\bm p)=(I+sH_M/E_p)/2$, and $s=\pm1$, the regulated static loop is evaluated from the particle/antiparticle spectral representation
\begin{equation}\label{eq:qm009-loop}
\Pi_{AB}(0,\bm q)=\int\frac{\dd^3p}{(2\pi)^3}\sum_{s,t=\pm1}
\frac{f_s(\bm p)-f_t(\bm p+\bm q)}{sE_p-tE_{p+q}}
\Tr\big[P_s(\bm p)\mathsf V_A P_t(\bm p+\bm q)\mathsf V_B\big],
\end{equation}
where $\mathsf V_M=\beta$, $\mathsf V_\mu=-I$, $f_-(\bm p)=1$ and $f_+(\bm p)=n_F(E_p-\mu)$. The vacuum interband piece has the explicitly stated ultraviolet regulator; the finite-density occupation correction is kept uncut. For the scalar vertex the Dirac trace has the closed form
\begin{equation}\label{eq:qm009-traces}
\Tr[P_s(\bm p)\beta P_t(\bm p')\beta]
=1+st\frac{M^2-\bm p\cdot\bm p'}{E_pE_{p'}},\quad
\Tr[P_s(\bm p)\beta P_t(\bm p')]=M\left(\frac{s}{E_p}+\frac{t}{E_{p'}}\right).
\end{equation}
The corresponding density--density trace is $1+st(M^2+\bm p\cdot\bm p')/(E_pE_{p'})$. All four energy-sign channels are retained, including the vacuum interband channel. With signs fixed by $\mathsf V_\mu=-I$, the static Hessians are $K_s=1/(2G_4)+\Pi_{MM}$, $V_s=\Pi_{M\mu}$, and $D_{nn}=-\Pi_{\mu\mu}$.

Minimizing over the scalar fluctuation gives the \emph{actual inhomogeneous source response} in this approximation:
\begin{equation}\label{eq:qm009-response}
\boxed{\delta M(\bm q)=-\frac{V_s(q)}{K_s(q)}\delta\mu(\bm q),\quad
 \chi_{nn}(q)=D_{nn}(q)+\frac{V_s(q)^2}{K_s(q)},\quad
 \frac{\delta M(\bm q)}{\delta n(\bm q)}=-\frac{V_s(q)}{K_s(q)\chi_{nn}(q)}.}
\end{equation}
The last formula describes linear response about equilibrium, not an arbitrary constitutive identity for local density outside a source.

At zero wave number the original sharp-cutoff benchmark is independently reproduced: $M=0.693435535520$, $p_F=0.979360586341$, $K_0=0.168741734136$, $V_0=0.068809590036$, $\chi_{nn}(0)=0.147135178153$, and
\begin{equation}\label{eq:qm009-numerics}
\boxed{\left.\frac{\dd M}{\dd n}\right|_{\rm equilibrium,\ sharp}=-2.77146847504.}
\end{equation}
For the same homogeneous source variation, the calculated scalar \emph{operator} vertex is nonzero even though its kinetic derivative is zero:
\begin{equation}\label{eq:qm009-opvertex}
\boxed{\left.\frac{\delta\mathcal D_H}{\delta n}\right|_{q=0,\mathrm{equil.}}=+2.77146847504\,I_4,\qquad\frac{\delta}{\delta n}\left(\partial_{p_\mu}\mathcal D_H\right)=0.}
\end{equation}
For a spinor with \emph{fixed canonical momentum} $p=p_F$ of the reference state, $\zeta=p/(E_p+M)$ satisfies $\partial_M\zeta|_p=-\zeta/E_p$. Therefore
\begin{equation}\label{eq:qm009-zeta}
\boxed{\zeta_0=0.517239994692,\qquad \left.\frac{\partial\zeta}{\partial n}\right|_{p\,\mathrm{fixed}}=1.19459528276.}
\end{equation}
This is a \emph{genuine source-state-to-spinor response}. It is not a derived map from the spinor ratio to the QGD gravitational natural invariant. The fixed-$p$ derivative in Eq.~\eqref{eq:qm009-zeta} must not be confused with the derivative tracking the moving Fermi surface $p_F(n)$.

For a finite-$q$ inhomogeneous calculation, a sharp cutoff independently imposed on both $|\bm p|$ and $|\bm p+\bm q|$ produces a nonanalytic $|q|$ edge artifact. We checked this numerically and do \emph{not} interpret it as collective dynamics. Instead, match the same $M$, $\mu$, and $m_0$ with smooth vacuum form factors
\begin{equation}\label{eq:qm009-regulator}
R_{n}(p)=\exp[-(p/\Lambda)^n],\qquad n=8,16,
\end{equation}
implemented as $\sqrt{R_n(p)R_n(p')}$ across the vacuum loop, and recalibrate $G_4$ by the respective gap equation. The resulting $T\to0$ zero-wave-number comparisons are
\begin{center}\small
\begin{tabular}{@{}lrrrr@{}}\toprule
Vacuum scheme & Matched $G_4$ & $K_s(0)$ & $\chi_{nn}(0)$ & $K_s(q=0.2)$ at $T=0.005$\\\midrule
Original sharp, $q=0$ only &0.56542557&0.16874173&0.14713518& ---\\
Smooth $n=8$ &0.62924767&0.14993818&0.15065405&0.15188956\\
Smooth $n=16$ &0.60370423&0.15737518&0.14916178&0.16062232\\\bottomrule
\end{tabular}
\end{center}
Here the nonzero-$q$ values use small thermal smoothing $T=0.005$ solely to control the Fermi-surface quadrature. Repeating with $T=0.010$ and increased quadrature resolution gives the expected stable infrared trend; the local model is visibly regulator dependent. The computed small-$q$ expansions are $K_s(q)\simeq K_s(0)+\kappa_s q^2$ with $\kappa_s\simeq0.0486$ (smooth $n=8$) and $0.0816$ (smooth $n=16$) in these model units. These are gradient coefficients within the specified schemes, \emph{not} an independently computed long-distance correlation length or an induced gravitational field.

\begin{figure}[htbp]\centering
\includegraphics[width=.86\linewidth]{smooth_scalar_Hessian.png}
\caption{Independently calibrated smooth-vacuum-regulator calculations of the finite-wave-number scalar response. Both have a nonzero stable $K_s(0)$; neither has an infrared static massless scalar pole in the tested phase.}
\end{figure}
\begin{figure}[htbp]\centering
\includegraphics[width=.86\linewidth]{sharp_cutoff_artifact.png}
\caption{The spurious $|q|$ term produced by a naive joint sharp cutoff on two shifted loop momenta, compared with smooth cutoffs. This regulator boundary effect is not a newly derived collective mode.}
\end{figure}

\subsection{Resulting source-to-kinetic verdict and the exact next physical requirement}
For both matched smooth regulators the static scalar Hessian remains positive at sufficiently small wave number and the source vertex is finite. Hence
\begin{equation}\label{eq:qm009-pole}
\boxed{\lim_{q\to0}q^2\frac{V_s(q)^2}{K_s(q)}=0.}
\end{equation}
The position-dependent scalar source changes the local Dirac spinor ratio, as Eq.~\eqref{eq:qm009-zeta} demonstrates, but the microscopic \emph{principal-symbol source-to-kinetic derivative} remains Eq.~\eqref{eq:qm009-zero}. Thus this phase does not independently produce the specific null-sector target
\begin{equation}\label{eq:qm009-target}
\frac{\delta\Gamma^\mu(x)}{\delta\rho_m(y)}\stackrel{\rm QGD\ target}{=}
-\frac{G_N}{c^2}\frac{\ell^\mu\not\!\ell}{|\bm x-\bm y|},
\end{equation}
which, in the Schwarzschild-specialized matching calculation, follows from the \emph{adopted} gravitational source law, not from this contact QFT. The mass density $\rho_m$ and the number-source $\mu$ are distinct perturbations; one must not equate their response kernels without a source/operator conversion and the correct units. Equation~\eqref{eq:qm009-target} remains an \textbf{[O]} microscopic target and is not a universal formula for every sector of the full QGD metric.

\textbf{Next independent test.} At the earliest momentum-dependent self-energy order, compute $\delta[\partial_{p_\mu}\Sigma_R(p;x)]/\delta n(y)$ with the matching two-particle irreducible source vertex, isolate any nonanalytic static $q\to0$ behaviour, then test the \emph{characteristic cone} against spinor-basis changes, the full source stress tensor, light propagation, and the complete coframe/signed-sector/$\cQ$ master metric. Higher-order nonlocality need not be zero; this chapter has not calculated its full finite-$q$ kernel. The effective QGD action route remains valid \emph{conditionally} without that microscopic derivation. The historical $u_B=\mathsf R u_A$, leading WKB Hamilton--Jacobi relation and $r^2j^r=C_J$ are retained throughout.

\textbf{Literature and regulator caution.} Nonrenormalizable NJL-type contact theories have strongly regulator-dependent inhomogeneous channels. For example Pannullo, Wagner and Winstel, \emph{Phys. Rev. D} \textbf{110}, 076006 (2024), explicitly compare regulator-dependent nonuniform phases (DOI:10.1103/PhysRevD.110.076006). These are not imported QGD predictions, but justify why neither a sharp cutoff cusp nor a regulator-specific gradient coefficient should be presented as fundamental gravity.


\section{Microscopic stationary selection: direct scalar-contact test (QM011)}\label{sec:qm011}
\textbf{Research question.} Does the already specified microscopic quantum theory produce a state-dependent stationary functional, and does its physically admissible stationary point select the QGD gravitational natural invariant independently of the adopted gravitational action? We test these questions separately. The scalar model's existence of a stationary collective field is a positive result; identification of that field with QGD's $\sigma_\mu^{(s)}$ is a further physical claim, not a definition.

\subsection{A microscopic Hubbard--Stratonovich functional, not a postulated gravitational energy}
In the zero-electromagnetic-background scalar-contact model of Chapter~1, and $\hbar=c=1$, an algebraic auxiliary scalar $\phi$ can be introduced:
\begin{equation}\label{eq:qm011-hs}
\mathcal L_{\rm HS}=\bar\psi(i\gamma^\mu\partial_\mu-m_0+\phi)\psi-\frac{\phi^2}{4G_4},\qquad M=m_0-\phi.
\end{equation}
Completing the square in $\phi$ recovers the original $+G_4(\bar\psi\psi)^2$ interaction. Integrating the regulated fermions formally gives a state-dependent collective effective action $\Gamma_{\rm micro}[\phi;\mu]$, with a Dirac determinant and a local quadratic auxiliary-field term. Its homogeneous static saddle is precisely the Chapter~1/5 Hartree equation
\begin{equation}\label{eq:qm011-gap}
0=\Omega_M=\frac{M-m_0}{2G_4}-\frac M{\pi^2}\int_{p_F}^{\Lambda}\frac{p^2\dd p}{\sqrt{p^2+M^2}},\quad p_F=\sqrt{\mu^2-M^2}.
\end{equation}
This is an \textbf{[D/C]} microscopic variational-selection functional within the declared regulator/mean-field approximation, with a physically specified collective variable $M$. It has not introduced the QGD metric, $H_{\rm KS}$, $\sigma$, or the Einstein action.

\subsection{The scalar stationary response is calculable and unique in the positive quadratic regime}
Write $q(\bm x)=M(\bm x)-M_\star$ around the stable uniform saddle and apply a small external chemical-potential source $\delta\mu(\bm x)$. The quadratic functional derived by expanding the same fermion determinant is
\begin{equation}\label{eq:qm011-gamma2}
\Gamma_{\rm micro}^{(2)}[q;\delta\mu]=\frac12\int\!\frac{\dd^3k}{(2\pi)^3}\,K_s(k)q(-k)q(k)+\int\!\frac{\dd^3k}{(2\pi)^3}\,V_s(k)q(-k)\delta\mu(k).
\end{equation}
Here $K_s=\Omega_{MM}+\Pi_{MM}(k)-\Pi_{MM}(0)$ and $V_s$ is the consistently computed scalar--density source vertex. Stationarity gives the exact Gaussian relation
\begin{equation}\label{eq:qm011-resp}
\boxed{\frac{\delta\Gamma_{\rm micro}^{(2)}}{\delta q(-k)}=0\quad\Longrightarrow\quad q(k)=-\frac{V_s(k)}{K_s(k)}\delta\mu(k).}
\end{equation}
QM009 determined the small-$k$ static expansion $K_s(k)=K_0+Z_s k^2+O(k^4)$ with two \emph{independently matched smooth vacuum regulators}, keeping the same physical homogeneous mass. Reusing those existing, reproducible inputs yields
\begin{center}\small
\begin{tabular}{@{}lrrrr@{}}\toprule
Scheme & $K_0$ & $Z_s$ & $V_0$ & $\xi_s:=\sqrt{Z_s/K_0}$\\\midrule
Smooth $n=8$&0.14993818&0.04864381&0.06880959&0.56958383\\
Smooth $n=16$&0.15737518&0.08155175&0.06880959&0.71986081\\\bottomrule
\end{tabular}
\end{center}
The $Z_s$ estimates are infrared fits to finite-$k$ loop data with $T=0.005$ quadrature smoothing; they are not exact $T=0$ analytic derivatives. At the displayed truncation order, the stationary equation is
\begin{equation}\label{eq:qm011-elliptic}
\boxed{(-Z_s\nabla^2+K_0)q(\bm x)=-V_0\delta\mu(\bm x).}
\end{equation}
For $K_0,Z_s>0$, the quadratic form is strictly positive in its stated smooth finite-energy field space (with boundary conditions), and the stationary scalar solution is unique there. This does not prove global uniqueness of the full non-Gaussian fermionic effective action.

\subsection{A compact spherical source: independent Green-function and field-equation tests}
Choose the transparent benchmark $\delta\mu(\bm x)=\delta\mu_0\Theta(R-r)$ as a localized \emph{number-source probe in the otherwise uniform finite-density medium}, not as a microscopic model of an isolated star in empty space. Let $m_s^2=K_0/Z_s=\xi_s^{-2}$. For $r>R$, the stationary gradient-truncated solution is
\begin{equation}\label{eq:qm011-exterior}
\boxed{q_{\rm out}(r)=-\frac{V_0\delta\mu_0}{Z_s}\frac{m_sR\cosh(m_sR)-\sinh(m_sR)}{m_s^3}\frac{e^{-m_sr}}{r}.}
\end{equation}
We derived Eq.~\eqref{eq:qm011-exterior} from the Green function $e^{-m_s|\bm x-\bm y|}/(4\pi Z_s|\bm x-\bm y|)$ and independently evaluated the radial three-dimensional Green-function integral. The two agree at 12 tested radii/regulator combinations to relative tolerance $10^{-10}$ or the floating-point absolute floor. Within that illustrative local-gradient model, after normalization to the same field at $r=R$, the ratio of microscopic scalar to a $1/r$ comparator at $r=R+d$ is precisely $e^{-d/\xi_s}$. For $d=1$ it equals $0.17279$ ($n=8$) and $0.24929$ ($n=16$); for $d=2$ it equals $0.02986$ and $0.06214$. The Yukawa profile is a \emph{mathematical diagnostic of the truncated infrared functional}, not a quantitatively controlled far-field prediction: $m_s\simeq1.76$ or $1.39$ is larger than the $q\lesssim0.2$ wave-number window used to fit $Z_s$, and is comparable to or greater than the Fermi momentum $p_F\simeq0.98$. The inferred imaginary Yukawa pole thus lies beyond the verified small-$q$ expansion. The \emph{full} static $T=0$ Fermi-surface response can have additional nonanalyticities at finite momentum (including Friedel-type tails). The reliable infrared conclusion is instead the absence of a $1/q^2$ residue, since $K_s(0)>0$ and $V_s(0)$ is finite.

\begin{figure}[htbp]\centering
\includegraphics[width=.86\linewidth]{selection_test.png}
\caption{Independent variational selection test: the QGD conditional Schwarzschild amplitude $H\propto1/r$ (normalized at $R=2$) versus stationary scalar fields using the previously calculated and refitted microscopic Hessians. The scalar fits have positive $K_0$ and a finite correlation length. These profiles are a comparison, not an identification of the scalar $q$ with $H_{\rm KS}$.}\label{fig:qm011-comparison}
\end{figure}

\subsection{What would have been sufficient, and why it is not yet derived}
For the local-gradient effective action, a static $1/r$ Green-function response of this particular scalar channel would require $K_0=0$ with $Z_s>0$ and a nonvanishing $k\to0$ source vertex, yielding $-Z_s\nabla^2q=-V_0\delta\mu$ and an exterior $q\propto 1/r$. Merely \emph{setting} $K_0=0$ is a hypothetical critical tuning, not a result at the tested Hartree saddle. Even if such a $q$ were selected, one would still have to prove the source conversion from $\delta\mu$ to gravitational mass/stress, establish universal probe coupling, identify the physically admissible null congruence $\ell_\mu$, and derive a legitimate relation $H_{\rm KS}=h(q)$ with $\sigma_\mu=\sqrt{H_{\rm KS}}\,\ell_\mu$ in the declared classical null-sector specialization. The full founding coframe, signed sectors and $\lambda\cQ_{ab}$ remain distinct.

A candidate massless collective field is \emph{not automatically sourced}. For illustration, a shift-symmetric phase mode $\theta$ can couple derivatively as $g\partial_\mu\theta\,j^\mu$. Integrating by parts makes that interaction proportional to $-g\theta\,\partial_\mu j^\mu$ plus a boundary term. For an exactly conserved interior current and trivial boundary contribution it cannot source the static phase field. This is a diagnostic of source coupling, not a new Goldstone mode derived from Chapter~1.

\begin{proposition}[What the microscopic test actually establishes]\label{prop:qm011-verdict}
The specified scalar-contact quantum field theory has a regulated Hubbard--Stratonovich stationary functional. In its matched stable Hartree--Gaussian scalar phase, the independently calculated static Hessian satisfies $K_s(0)>0$, and its scalar-to-number-source vertex is finite. Thus the stationary scalar response is regular at $k=0$ and does \emph{not} select the $1/r$ QGD gravitational amplitude from a localized number-source perturbation by the required infrared mechanism. It does select a physical \emph{scalar mass response} and corresponding Dirac spinor changes. This is a phase/truncation-specific failure of the proposed scalar promotion, not a no-go theorem for a more general physical map $\mathcal F$ or a contradiction of the already closed \emph{conditional} QGD Schwarzschild bridge.
\end{proposition}

\paragraph{Decision for the next microscopic hypothesis.} Keep $\mathcal F_{\rm conditional}$ and Chapter~2's variational Schwarzschild selection closed. For $\mathcal F_{\rm micro}$, the first unresolved property is \emph{the existence of a physically defined state-to-natural-invariant rule} with source and exterior boundary data, not the existence of a standalone graviton. An alternative quantum collective variable must exhibit (i) an independently derived stationary functional, (ii) source/exterior transfer of the required form with justified physical boundary conditions, and (iii) a frame- and probe-independent lift into the full QGD sectors. The original Dirac ratio, WKB principal-symbol equation and $r^2j^r=C_J$ remain independent and preserved; none is silently substituted for the new selection equation.

\paragraph{External comparison.} Standard fermionic Hubbard--Stratonovich/mesonic effective-action calculations likewise identify the inverse collective propagator as $1/(2G)+\Pi(k)$ and distinguish scalar and (in enlarged chiral models) symmetry-related Goldstone channels; see \url{https://doi.org/10.1088/1674-1137/43/9/094103}. That literature supports the method of this test, not the physical identification with QGD's $\sigma$.

% ===== v0.8 research synthesis: inserted without removing original v0.7 derivations =====
\section{Research consolidation: what the programme has actually achieved (QM012--QM014)}\label{sec:v08-small-wins}
\textbf{Editorial status (2 November 2026).} This section records independent research notes QM012, QM013 and QM014 after the v0.7/QM011 edition. It preserves all earlier calculations and equation labels. Its purpose is to identify small but decisive advances and to prevent an externally proposed induced-action construction from being mistaken for the original microscopic QGD closure. The additional calculations are not silently promoted into the founding axioms. The status convention [D/C/A/O] is that of the chapter's opening contract.

\subsection{Two arrows, two logically different scientific questions}\label{sec:v08-two-arrows}
The founding objective is a source-state and macroscopic-Dirac-to-natural-invariant selection law, schematically
\begin{equation}\label{eq:v08-original-direction}
\boxed{\mathcal D_{\rm eff}\Psi=0\ \longrightarrow\ (\Psi,S,\mathsf R,j^\mu;\hat\varrho_{\rm src})
\ \xrightarrow{\ \mathcal F_{\rm micro}\ }\ [\Sigma]_{\rm physical}.}
\end{equation}
This notation is a \emph{research target}, not an equation derived by applying the Dirac operator: a probe spinor does not contain the full source-state information in general. The required map must use independently characterized source-state observables and must not depend on selecting an arbitrary probe velocity, spin polarization, spinor basis or already-known gravitational field. Chapter~1 and Eqs.~\eqref{eq:F-gap}--\eqref{eq:qm010-microstationarity} state the controlled ingredients and the remaining gap.

The alternative suggested in the external Claude compendium begins with a \emph{new premise (Q)}: independent admissible QGD sector variables $\Sigma$, a specified functional measure/saddle prescription, and matter universally coupled to the composite metric $g(\Sigma)$. Its logical direction is
\begin{equation}\label{eq:v08-induced-direction}
\boxed{S_{\rm micro}[\psi;g(\Sigma)]\ \xrightarrow{\int\! D\psi D\bar\psi}\
\Gamma_{\rm ind}[\Sigma]\ \xrightarrow{\delta\Gamma_{\rm ind}=0}\
[\Sigma]_{\rm stationary}\ \longrightarrow\ (\Psi,S,\mathsf R,j^\mu).}
\end{equation}
The induced action and its solutions can be meaningful \textbf{[C]} calculations. They do not invert Eq.~\eqref{eq:v08-original-direction}: the metric and its independent sector variables were supplied at the beginning. On a regular spanning branch $\delta_\Sigma\Gamma_{\rm ind}=(D_\Sigma g)^*\delta_g\Gamma_{\rm ind}$; at a rank-deficient branch its vanishing need not imply $\delta_g\Gamma_{\rm ind}=0$. The proposed premise (Q) was \textbf{not adopted} as a proven rewriting of the flat-reference original microscopic F1 theory.

\subsection{Small-win ledger: positive results, negative results, and their scope}\label{sec:v08-ledger}
The programme contains more than the still-open final arrow. The following claims deliberately distinguish a result established under a defined mathematical/model framework from a claim of fundamental quantum-gravitational emergence.
\begin{longtable}{@{}p{.075\linewidth}p{.36\linewidth}p{.075\linewidth}p{.38\linewidth}@{}}
\toprule
No.&Verified or restricted advance&Status&What it actually buys QGD\\\midrule\endhead
01&Exact positive-energy Dirac upper/lower block relation $u_B=\mathsf R u_A$.&D/C&Controlled microscopic spinor kinematics; not yet a gravitational one-form.\\
02&WKB Hamilton--Jacobi/transport and the stationary spherical $r^2j^r=C_J$ law.&D/C&Preserves the historic wavefunction--momentum/current bridge without inferring a force from flux.\\
03&Chapter~1's separated large-$N$ centre-of-mass $\hbar/N$ limit and correlation counterexample.&D/C&Demonstrates a precise possible classical collective limit; rejects the claim that large $N$ alone yields a unique Dirac collective spinor or gravity.\\
04&Complete coframe/signed-sector/$\mathcal Q_{ab}$ metric and full Euler pullback.&A/D&Provides the exact representation and variational language, without deriving its microscopic ontology.\\
05&Independent $A,B$ Einstein--Hilbert variations select $A^2=B^2=2G_NM/(c^2r)$ on the specified nonzero exterior branch.&C&Schwarzschild natural invariant is \emph{dynamically} selected within adopted dynamics; not algebraically guessed.\\
06&One Newtonian mass calibration, exact static energy/potential consistency and distinct wavefunction vs geometric normalizations.&D/C&Eliminates spurious gravitational constants and unsupported universal $\hbar$-force claims.\\
07&Chapter~5 substitution $H=2g_cq/c^2$ makes the assumed collective and Newtonian $H$ energies exactly equal, with $G_{\rm eff}=g_c^2/(4\pi Z)$.&C&Closes the specified static axiomatic correspondence while leaving premise MW2 physical.\\
08&Source-conditioned stationary $H$ functional has a unique regular decaying solution for a compact source.&C&Fixes $H\sim1/r$, its mass coefficient and $\sigma_{\rm amp}\sim r^{-1/2}$ on the designated spherical null branch.\\
09&Exact null-sector Clifford/principal-symbol matching and radial Hamilton--Jacobi roots for \emph{given} $H$.&D/C&Checks the geometry-to-Dirac return map; does not reverse it to derive $H$.\\
10&Explicit source-to-$H$ and source-to-kinetic $1/|\bm k|^2$ \emph{matching targets}.&D/C&Converts the missing origin into a falsifiable source-response test, not a new microscopic calculation.\\
11&Regulated Hartree stationary scalar and Gaussian Hessian/source vertex constructed from actual F1.&D/C&Proves microscopic collective stationarity is calculable in one specified quantum phase.\\
12&Two matched smooth regulators find $K_s(0)>0$ and finite $V_s(0)$; QM011 Green-function check passes.&D/C&Rules out the desired static massless scalar residue in this tested regime; does not rule out all possible phases.\\
13&Leading Hartree/first local Fock principal-symbol source derivative is zero.&D/C&Exposes a concrete mismatch to the required universal source-to-kinetic response in the tested phase.\\
14&Isotropic scalar--TT mixing selection, collisionless TT continuum, full leading intraband stress matrix and Ward-contact checks.&D/C&Distinguishes real stress dynamics from a relabelled scalar; no autonomous graviton-like pole derived.\\
15&QM012 shows naive Laplace stiffness for $\sigma_{\rm amp}$ has wrong radial exponent; the $H=\sigma_{\rm amp}^2$ stiffness succeeds conditionally.&D/C&Identifies the correct \emph{effective field coordinate} and its nonlinear $\sigma$ stiffness.\\
16&QM012 WKB density fourth-root exponent, spinor-bilinear null construction and statistical selection controls.&D/C&Provides candidates but also explicit failures of source-mass normalization, spin/probe independence or exterior selection.\\
17&Regulated Dirac determinant induces a local curvature coefficient when matter is coupled to an \emph{already defined} metric.&D/C&Demonstrates induced-action consistency, not the original matter-to-$\Sigma$ construction.\\
18&Regular eight-sector Jacobian rank 10, nullity 22 and $O(4,4)$ orbit equality; singular-rank counterexamples.&D&Identifies physical-representation redundancy only on the stated metric-only spanning branch.\\
19&Conditional induced Einsteinian operator has two on-shell TT polarizations and a gauge-dependent local canonical-energy example.&C/D&Checks the adopted low-energy spectrum/gauge accounting, not a full microscopic graviton or Hamiltonian-constraint theorem.\\
20&Chapter~8 regular metric-only no-extra-equation theorem.&C&Explains why a field redefinition alone gives no new cosmological observable beyond the input semiclassical metric action.\\\bottomrule
\end{longtable}
A \textbf{negative test} is a real gain only in its stated domain: neither the failure of a local scalar-contact channel nor the success of a universal representation is a universal theory of gravity. Conversely, source-visible massless static response and an admissible common propagation metric are separate requirements.

\subsection{QM012: why varying the correct effective variable matters}\label{sec:v08-H-not-sigma}
In the \emph{fixed-reference}, $\lambda\mathcal Q=0$, positive null-sector specialization alone, take a fixed null direction $\ell_\mu$ and write
\begin{equation}\label{eq:v08-H-branch}
\boxed{g_{\mu\nu}=\eta_{\mu\nu}+H\ell_\mu\ell_\nu,\quad
\sigma_\mu=\sigma_{\rm amp}\ell_\mu,\quad
H=\sigma_{\rm amp}^2\geq0.}
\end{equation}
This is \emph{not} a replacement for the full founding metric. Consider first the \emph{specific trial} canonical vacuum energy $\mathfrak E_{\rm can}[\sigma_{\rm amp}]=(Z/2)\int |\nabla\sigma_{\rm amp}|^2\dd^3x$. Its exterior spherical stationary equation is $\nabla^2\sigma_{\rm amp}=0$, with decaying solution $\sigma_{\rm amp}=a/r$. Consequently $H=a^2/r^2$, which is not the Schwarzschild $H\propto1/r$. This failure excludes this particular simple stiffness, not all possible microscopic actions for $\sigma$.

The \emph{adopted} static QGD/GR functional, Eq.~\eqref{eq:qm010-energy}, is instead varied in $H$. Its exact Euler derivative is
\begin{equation}\label{eq:v08-H-Euler}
\frac{\delta\mathfrak E_{\rm N}}{\delta H}=-\frac{c^4}{16\pi G_N}\nabla^2 H-\frac{c^2}{2}\rho_m,
\qquad \nabla^2H=-\frac{8\pi G_N}{c^2}\rho_m.
\end{equation}
Regularity through a compact source, $H\to0$ at infinity and the mass integral then fix the exterior solution $H=2G_NM/(c^2r)$, $\sigma_{\rm amp}=\sqrt{2G_NM/(c^2r)}$. Its uniqueness refers to the chosen source, function space and boundary data; vacuum exterior harmonicity alone leaves the coefficient undetermined.

\textbf{A useful new explicit identity [D/C].} Rewriting the \emph{same} energy in the scalar amplitude is not the naive canonical energy:
\begin{equation}\label{eq:v08-nonlinear-stiffness}
\boxed{\mathfrak E_{\rm N}[\sigma_{\rm amp}^2;\rho_m]
=\int\!\dd^3x\left[\frac{c^4}{8\pi G_N}\sigma_{\rm amp}^2|\nabla\sigma_{\rm amp}|^2
-\frac{c^2}{2}\rho_m\sigma_{\rm amp}^2\right].}
\end{equation}
By the chain rule its unrestricted scalar-amplitude Euler derivative is
\begin{equation}\label{eq:v08-zero-degeneracy}
\boxed{\frac{\delta\mathfrak E_{\rm N}[\sigma_{\rm amp}^2]}{\delta\sigma_{\rm amp}}
=2\sigma_{\rm amp}\left[-\frac{c^4}{16\pi G_N}\nabla^2(\sigma_{\rm amp}^2)-\frac{c^2}{2}\rho_m\right].}
\end{equation}
On $\sigma_{\rm amp}\ne0$, this is equivalent to Eq.~\eqref{eq:v08-H-Euler}. At $\sigma_{\rm amp}=0$ it vanishes even when the $H$ Euler equation fails, revealing a \emph{spurious stationary branch of a noninvertible field coordinate}. In the vacuum $r>R$, the correct $\sigma_{\rm amp}=\sqrt{r_s/r}$ has
\begin{equation}\label{eq:v08-sigma-radial-identity}
\nabla^2\sigma_{\rm amp}=-\frac{\sigma_{\rm amp}}{4r^2},\qquad
|\nabla\sigma_{\rm amp}|^2=\frac{\sigma_{\rm amp}^2}{4r^2},\qquad
\nabla^2(\sigma_{\rm amp}^2)=0.
\end{equation}
This resolves the apparent discrepancy between harmonic $H$ and non-harmonic $\sigma_{\rm amp}$. Selecting $H$ improves the conditional stationary coordinate in this branch; it is not a proof that the original flat Dirac QFT generates Eq.~\eqref{eq:v08-nonlinear-stiffness}. In a general multi-sector configuration the individual sector amplitudes and $H$ representatives can change under redundancies, so an operational, gauge/representation-invariant lift remains necessary.

\subsection{QM012 controls: which apparent ``spinor origins'' actually fail}\label{sec:v08-qm012-controls}
The independent comparative test supplies three especially useful controls: (i) a stationary outgoing free WKB wave has $|\Psi|\sim r^{-1}$ and constant flat-space local lower/upper ratio, not $\sigma_{\rm amp}\sim r^{-1/2}$; (ii) a trial fourth-root density $(\Psi^\dagger\Psi)^{1/4}$ has the correct radial \emph{power} but a free outgoing $1/r$ wave is not an isolated normalizable source and its coefficient is a flux rather than source-mass normalization; (iii) pure-spinor Fierz bilinears can furnish null directions algebraically but generally retain spin-polarization dependence, while mixed unpolarized sources need not select a unique null congruence. Maximum entropy selects a state only under imposed constraints, not automatically a universal exterior $H$. These are useful restrictions on $\mathcal F_{\rm micro}$, not assertions that all collective-spinor possibilities fail.

\subsection{QM013: proper-time induced-curvature win and the limits of premise (Q)}\label{sec:v08-qm013}
Assuming metric-coupled Dirac quantum matter with a covariant proper-time prescription, the checked local vacuum coefficient in the specified one-four-component-Dirac convention is
\begin{equation}\label{eq:v08-induced-coeff}
I_1=\int_{\Lambda^{-2}}^\infty\!s^{-2}e^{-sM^2}\dd s,
\quad C_R=-\frac{I_1}{96\pi^2},\quad
M_P^2=-2C_R=\frac{I_1}{48\pi^2},\quad G_{\rm ind}=\frac{6\pi}{I_1}.
\end{equation}
At $M=0.693435535520$ and $\Lambda=4$ in the original dimensionless model, the independent audit reports $I_1=14.10422862355$, $C_R=-0.01488601185938$ and $G_{\rm ind}=1.3364471340$. Claude's alternative expression $12\pi/(\sum_i w_iI_{1,i})$ agrees for one such Dirac species \emph{only if} its unspecified weight is $w_i=2$; setting $w_i=1$ yields a factor-of-two discrepancy. Neither matching a heat-kernel coefficient nor obtaining the textbook $1/|\bm k|^2$ weak-field response \emph{after} admitting a dynamical metric establishes that $\Sigma$ was derived as a composite observable of the original flat QFT. Vacuum-energy tuning, regularity, physical source vertices, regulator matching, the measure and no double counting of a separately retained Einstein term are still necessary. A causal non-equilibrium source response also requires an appropriate in-in rather than merely Euclidean vacuum prescription.

\subsection{QM014: exact sector ranks and a narrow induced-spectrum check}\label{sec:v08-qm014}
For the special fixed-coframe, $\lambda\mathcal Q=0$, eight-sector model, let $J=\operatorname{diag}(I_4,-I_4)$ and $g=\eta+\Sigma^{\mathsf T}J\Sigma$. Exact rational-matrix tests reproduce:
\begin{center}\small
\begin{tabular}{@{}lrrrr@{}}\toprule
Configuration & $\operatorname{rank}D_\Sigma g$ & $\dim\ker$ & $O(4,4)$ orbit & Stabilizer\\\midrule
Flat spanning pair & 10 & 22 & 22 & 6\\
Second spanning Lorentzian pair & 10 & 22 & 22 & 6\\
One rank-one sector & 4 & 28 & 7 & 21\\
All sectors zero & 0 & 32 & 0 & 28\\\bottomrule
\end{tabular}
\end{center}
The 22-dimensional kernel equals the local $O(4,4)$ orbit tangent \emph{on the spanning branch}; singular branches have additional null directions that are not automatically gauge transformations. At a tuned, flat, spanning background and \emph{assuming} the two-derivative induced Einsteinian action, the Fourier metric Hessian has rank 6 at generic non-null momentum (four diffeomorphism null directions) and rank 4 at null momentum (four gauge plus two physical wave solutions). Pullback to 32 sectors gives respective nullities 26 and 28, of which 22 are algebraic sector redundancies. These are reproducible \textbf{[D/C]} operator-rank checks, not a full nonlinear Dirac--Bergmann count for the original microscopic theory.

The same audit supplies an explicit gauge transformation of a TT wave for which the local canonical Fierz--Pauli $t^{00}_{\rm can}$ changes, although its periodic spatial average does not. Therefore $\mathcal G+\mathcal T+\mathcal Q=0$ (a sector Euler balance) cannot be read as three separately physical local energy densities, and the total matter stress should not be called zero inside an ordinary sourced Einstein solution. Most importantly, the Weinberg--Witten obstruction is \emph{not} eliminated just by writing a stationary total-stress equation. A proposed microscopic spin-two construction must identify its true Lorentz-covariant stress operator and show precisely which theorem assumptions fail after the constraints are constructed. The classical sector orbit counts alone cannot do this.

\subsection{Current research decision: the smallest still-missing physical result}\label{sec:v08-decision}
\textbf{Conditional QGD objective achieved [C].} Under a supplied effective QGD/Einstein gravitational action, specified source, boundary data, regular nonzero branch and null-sector matching, Chapter~2 selects Schwarzschild, Chapter~5 exactly matches its static collective functional, and the induced curved Dirac/WKB/current structure is consistent. This is a closed \emph{conditional} statement, independent of whether premise (Q) is later adopted.

\textbf{Physical genesis remains unproved [O].} The minimum new result that would advance the original direction is a \emph{source-theory-derived}, physically defined, probe-independent collective observable $H_{\rm micro}$ (or a suitable tensor/constitutive replacement) and a constrained state-dependent functional $\Gamma_{\rm micro}[H_{\rm micro},\ldots;J]$ whose Hessian and \emph{physical mass/stress} source vertex determine an exterior, source-visible long-range response. On an ordinary regular Gaussian branch the benchmark is
\begin{equation}\label{eq:v08-micro-target}
\boxed{\frac{\delta H_{\rm micro}(\bm k)}{\delta\rho_m(\bm k)}
=-[K_{\rm phys}^{-1}(\bm k)B_{H\rho}(\bm k)]
\stackrel{?}{\sim}\frac{8\pi G_N}{c^2|\bm k|^2}
\qquad(|\bm k|\to0).}
\end{equation}
The displayed $G_N$ is a measured effective matching target, not something QGD must independently predict to count an effective milestone. The formula presumes a specified linear static regime, consistent definitions and fixed Fourier conventions; a genuinely nonlinear/non-Gaussian selection mechanism would need its own sourced exterior solution, not a forced Gaussian pole. Only after this physical response is obtained should one derive the frame/probe-independent lift $H_{\rm micro}\to[\Sigma]$ and verify universal matter/light propagation, source stresses, causal constraints, and the full coframe/signed-sector/$\mathcal Q$ history. Inducing a curvature coefficient from a supplied metric does not substitute for that microscopic step.

\textbf{Optional alternative to audit, not an adopted foundational closure.} Premise (Q) may be explored on a separate track: supply a complete auxiliary-$\Sigma$/Dirac Lagrangian and measure, a consistent vacuum prescription, full nonlinear constraints including spin-connection and fermionic Dirac brackets, diffeomorphism and sector generators, boundary charges and Ward identities; then show how a composite universal metric and its microscopic physical states arise or declare them new axioms. The QM014 rank tests are necessary but not sufficient. A spin-two pole calculation is a diagnostic, not a required first step of the founding $\mathcal F_{\rm micro}$ project.

\textbf{Reproducibility and provenance.} The numerical QM011 scalar comparison was rerun using \texttt{test\_qm011.py}; the QM012 selection diagnostics were rerun using \texttt{tests\_qm012.py}; and QM014 rank/gauge diagnostics were rerun using \texttt{verify\_qm014.py}, each using the preserved research packages. QM013's normalization comparison is independently archived in \texttt{QM013\_Independent\_Audit.md}. These are tests of stated mathematical/model claims, not evidence for a completed microscopic-to-gravity derivation. Relevant primary method literature: Cornwall, Jackiw and Tomboulis, \emph{Phys. Rev. D} \textbf{10}, 2428 (1974), \url{https://doi.org/10.1103/PhysRevD.10.2428}; Visser, \emph{Sakharov's induced gravity}, \url{https://arxiv.org/abs/gr-qc/0204062}; Weinberg and Witten, \emph{Phys. Lett. B} \textbf{96}, 59 (1980), \url{https://doi.org/10.1016/0370-2693(80)90212-9}.
% ===== end v0.8 research synthesis =====

\section{Conclusion and dependency ledger (v0.7 baseline preserved)}
\begin{center}\small
\begin{tabularx}{\linewidth}{@{}>{\raggedright\arraybackslash}p{.31\linewidth}>{\raggedright\arraybackslash}p{.14\linewidth}>{\raggedright\arraybackslash}X@{}}\toprule
Research result&Status&Meaning\\\midrule
Dirac ratio, selected WKB, total current conservation&D/C&Matter-to-transport bridge retained\\
Hartree scalar Hessian and density-source vertex&D/C&Numerically reproduced using the original scalar-contact cutoff model\\
Stable scalar-channel static pole&D/C&Absent in sharp-3D and matched proper-time schemes at the stated saddle\\
TT--scalar two-point mixing&D/C&Zero in an isotropic state by rotational tensor selection\\
Strictly local matter-to-$\Sigma$ exterior rule&D/C&Insufficient to produce a nonconstant source-dependent vacuum exterior\\
TT shear, rank-nine intraband matrix, and kinetic Ward contact&D/C&Independently checked, no isolated TT pole at this order\\
Null-sector Dirac principal symbol and radial WKB&D/C&Exact Clifford/coframe identity; reproduces QGD Kerr--Schild propagation given $H$\\
Conditional source-to-kinetic response&D/C&Calculated $1/|\bm q|^2$ functional target from adopted gravity equation; not derived from QFT\\
Hartree/Fock source-to-kinetic vertex&D/C&Zero at tested leading orders, fails conditional target\\
Full source-visible stress/tensor pole&O&No pole in scalar-Hartree Gaussian TT channel; other interactions require conserving tests\\
Universal optical/matter metric&O&Requires multi-probe kinetic and interaction tests\\
Microscopic Hubbard--Stratonovich stationary scalar&D/C&Existence derived; selected massive scalar response does not match $H_{\rm KS}$ in tested phase\\
Variational weak-field $H$ selection&C&Source plus adopted stationary field energy and boundary conditions uniquely determine the spherical exterior $H$\\
Physical state-to-$[\Sigma]$ selection&O&Requires invariant quantum-state/Dirac-to-invariant rule or microscopic selection functional\\
Induced complete QGD Einsteinian action&O&Currently adopted in Chapters 2--5 and tested conditionally\\\bottomrule
\end{tabularx}
\end{center}
\textbf{Revised v0.7 research decision.} There are two separate closure goals. The \emph{conditional QGD effective correspondence} is closed for the stated Schwarzschild exterior: quantum matter supplies the source, the \emph{adopted} gravitational Euler equations select $H$ and the natural invariant, the compatible curved Dirac operator has the exact required principal symbol, and WKB/current conservation match without discarding the original bridge. This does \emph{not} establish the microscopic origin of gravitational geometry. QM009 adds the source-dependent spinor response but proves the Hartree--Fock kinetic vertex remains zero for the actual local scalar-contact model; a momentum-dependent self-energy and its conserving source vertex remain to be tested.

QM011 has now tested the microscopic stationarity hypothesis itself: a regulated matter-state-dependent functional exists and selects a unique scalar field at Gaussian order, but its measured positive scalar Hessian and finite source vertex do not select QGD's $1/r$ exterior. This negative result does not reopen the conditional closure and does not impose a graviton-first requirement.\par\noindent The \emph{independent origin} remains open at a now precise point: find a physically defined, observer- and probe-independent map from the solved macroscopic Dirac/quantum source state to $[\Sigma]$, or derive a microscopic constrained selection functional with that map as its stationary solution. The source-to-kinetic vertex of Eq.~\eqref{eq:actual-target} is one testable \emph{candidate route}, not an obligatory definition of $\mathcal F$; if that route is chosen, its result must match Eq.~\eqref{eq:source-vertex} without inserting the gravitational Green function. In the tested scalar-Hartree/first-Fock phase that vertex vanishes, whereas the stipulated QGD target is nonzero and nonlocal. The previously computed tensor/stress response remains a valuable parallel investigation, not the automatic prerequisite for the original spinor-to-invariant matching. Chapters 7 and 8 may be studied conditionally on the declared QGD effective action while the stronger origin question is pursued.


\section{Final v0.8 research decision and non-circular acceptance test}\label{sec:v08-final}
\textbf{Result of the consolidation.} The research has \emph{already} established a mathematically controlled geometry/Dirac correspondence and an exactly matched \emph{conditional} static collective energy. Its shortfall is not a missing Schwarzschild calculation, a missing radial exponent or a missing stress-tensor notation. It is the absence of a quantum-source-derived physical selection law $\mathcal F_{\rm micro}$ which fixes the same long-range field and universal coupling \emph{without starting from a supplied metric, Einstein equation or gravitational energy}. The massless collective postulate MW2 is neither contradicted by the tested gapped Hartree phase nor derived from it. The separate premise (Q) remains a distinct, explicitly unadopted induced-action route.

\textbf{Operational closure criterion.} To claim original-direction microscopic closure, a candidate calculation must (i) specify an actual quantum state/phase and observable $\mathcal O_{\rm micro}$, with its non-arbitrary definition, (ii) derive an effective constrained functional and \emph{conserving physical source vertex} from the same microscopic interaction, (iii) obtain the sourced exterior relation and normalize it to source mass without using the target Green function as input, (iv) show a physically justified common kinetic/light-cone response and a nondegenerate map to a QGD $[\Sigma]$, and (v) pass the full-master-metric, sector-redundancy, stability, conservation and boundary tests. Matching a coefficient or radial power in isolation is a useful diagnostic but is not sufficient.

\textbf{The next concrete research calculation.} Construct a specified composite/2PI observable family containing source-state stress, kinetic and Dirac-bilinear data; calculate its regulated $\Gamma_{\rm micro}[\varphi;J]$, its reduced inverse Hessian after constraints, and the static mass-source response. First decide whether that response is finite, critical, nonanalytic or non-Gaussian. \emph{Only if} a physically sourced long-range mode survives should one attempt the QGD $H\to[\Sigma]$ lift and universal test. As a parallel, independent research programme, an adopted premise (Q) requires a complete nonlinear Dirac--sector constraint algebra rather than a claim of automatic Weinberg--Witten evasion. Neither route invalidates the earlier closed conditional results.

\textbf{Editorial seal.} The original v0.7 core, source calculations and historic proof-status ledger remain intact; the present v0.8 text appends a narrowly scoped QM012--QM014 synthesis and this updated stopping rule. It does not promote any open [O] claim to [D].



% ===== BEGIN v0.9 QM015 extension; original v0.7 and v0.8 sections preserved =====
\section{QM015: source-overlap criterion and coupled-channel micro-origin programme}\label{sec:v09-qm015}
\textbf{Purpose and proof scope.} The previous sections isolate the physical missing arrow from quantum matter to a sourced natural invariant. We now sharpen it using an exact spectral-response criterion, two independent false-positive controls, and a solvable mixed-channel stationary model. In this extension \textbf{[D]} is a mathematical conclusion under explicitly stated hypotheses, \textbf{[C]} is conditional on the specified toy functional or a named approximation, \textbf{[A]} is an independently supplied kinetic/source ingredient, and \textbf{[O]} remains an unproved QGD microscopic mechanism. These calculations do not modify the founding master metric, existing Chapters 1--5 or the alternatives catalogued as premise (Q).

\subsection{The spectral source-overlap theorem [D]}\label{sec:v09-source-overlap}
Work at zero frequency in an independently chosen translation-invariant microscopic source state. Let $\varphi_A$ denote renormalized, physically specified composite observables, which may include scalar bilinears, conserved stress channels, a momentum-dependent Dirac self-energy, or constrained order parameters. Their stationary quadratic response must follow from a \emph{single} quantum-matter functional, not from an imported gravitational propagator. In a real positive Euclidean static subspace, write its quadratic change under an operational physical source $\rho$ as
\begin{equation}\label{eq:v09-general-functional}
\delta^2\Gamma_{\rm micro}=\frac12\int_{\bm k}\varphi_A(-\bm k)K_{AB}(\bm k)\varphi_B(\bm k)
-\int_{\bm k}\varphi_A(-\bm k)B_A(\bm k)\rho(\bm k),\qquad
\int_{\bm k}\equiv\int\frac{\dd^3k}{(2\pi)^3}.
\end{equation}
The source $\rho$ is an operational long-wavelength mass/energy perturbation, to be defined from a second microscopic matter species or a conserving stress insertion. A trial linear observable $H_{\rm micro}=A_A\varphi_A$ has
\begin{equation}\label{eq:v09-susceptibility}
\boxed{\chi_{H\rho}(\bm k)\equiv\frac{\delta H_{\rm micro}}{\delta\rho}
=A_A(\bm k)[K^{-1}(\bm k)]_{AB}B_B(\bm k).}
\end{equation}
The sign agrees with the displayed $-\varphi\cdot B\rho$ convention; an opposite microscopic source convention reverses it. The formula is valid on the constrained \emph{physical} subspace after zero-gauge directions have been appropriately quotiented; it is not a license to invert an un-fixed gauge Hessian.

\begin{proposition}[Finite-Hessian obstruction and source-overlap criterion]\label{prop:v09-overlap}
Assume finitely many physical static channels, $K(\bm k)$ real symmetric, analytic near $\bm k=0$, and with $K(0)$ uniformly positive definite. If $A(\bm k)$ and $B(\bm k)$ are also analytic there, then $\chi_{H\rho}(\bm k)$ is bounded and analytic near the origin; it cannot contain the required static $1/|\bm k|^2$ pole. Conversely, suppose a normalized isolated eigenmode $u_0$ has $\lambda_0(k)=\mu^2+Zk^2+O(k^4)$ with $Z>0$, while all other relevant modes are gapped and analytic. Its leading response is
\begin{equation}\label{eq:v09-residue}
\chi_{H\rho}^{(0)}(k)=\frac{(A\cdot u_0)(u_0\cdot B)}{\mu^2+Zk^2}+\text{terms regular at }k=0.
\end{equation}
A $1/k^2$ mass-source response therefore requires $\mu=0$, \emph{nonzero overlap on both sides}, and a nonvanishing net residue. The conditions are necessary and sufficient within the stated isolated Gaussian $k^2$ mode class, not across all non-Gaussian or nonlocal selection mechanisms.
\end{proposition}
\textit{Proof.} If $K(0)>0$, $K^{-1}(k)$ has a convergent analytic expansion at the origin; multiplication by regular $A$ and $B$ preserves this property. For the second claim, insert the orthonormal eigenmode decomposition of $K^{-1}$ and isolate the eigenvalue $\lambda_0$. If either projection vanishes, the pole is invisible to that source-observable pair. Multiple critical modes may cancel, and a symmetry can force $u_0\cdot B=0$ even when the mode is genuinely gapless.\hfill$\square$

This is the operative upgrade to the Chapter 5 premise MW2: instead of asking only whether a mode has $Zk^2$, we must ask whether the \emph{same microscopic calculation} yields a nonzero mass/stress vertex $B$, a physical observable projection $A$, and the positive long-range residue. A density pole, a time-dependent sound pole and a gravitational static response are not interchangeable.

\subsection{A concrete negative test: a massless Goldstone mode can be source-invisible [D/C]}\label{sec:v09-parity}
In a homogeneous, isotropic, parity-preserving state of Dirac quantum matter define the parity-even scalar $S=\bar\psi\psi$ and parity-odd pseudoscalar $P=\bar\psi i\gamma_5\psi$. At zero frequency, their mixed response obeys
\begin{equation}\label{eq:v09-parity}
\begin{aligned}
\chi_{PS}(0,\bm k)&=-\chi_{PS}(0,-\bm k) &&\text{(parity)},\\
\chi_{PS}(0,\bm k)&=+\chi_{PS}(0,-\bm k) &&\text{(isotropy)},\\
&\therefore\quad\boxed{\chi_{PS}(0,\bm k)=0}.&&
\end{aligned}
\end{equation}
If the microscopic mass insertion is parity even, an exactly gapless parity-odd Goldstone mode need not respond to it linearly. Chiral fermion theories provide a familiar controlled \emph{candidate category}: pseudoscalar modes may be symmetry-protected while scalar amplitude modes are massive. This excludes the inference ``fermionic Goldstone pole $\Rightarrow$ static gravity''. Equation~\eqref{eq:v09-parity} is exact only under the explicitly stated symmetry and source assumptions; parity-breaking states, anisotropic media or different source vertices must be analyzed separately.

A second independent control is a zero-temperature idealized compressible superfluid. At quadratic hydrodynamic order a density fluctuation $\delta n$ and phase $\theta$ have the static functional
\begin{equation}\label{eq:v09-hydro}
\mathfrak H[\delta n,\theta;J]=\int\dd^3x\left[\frac{(\delta n)^2}{2\chi_n}+\frac{\rho_s}{2}|\bm\nabla\theta|^2-J\delta n\right],\quad \chi_n,\rho_s>0.
\end{equation}
The phase has a gapless dynamical phonon after adding the canonical $\delta n\dot\theta$ term, but static minimization gives $\delta n=\chi_n J$ and $\nabla^2\theta=0$. More explicitly, with $c_s^2=\rho_s/\chi_n$ the linear density response has the idealized form
\begin{equation}\label{eq:v09-superfluid}
\chi_{nn}^{R}(\omega,k)=\frac{\chi_n c_s^2k^2}{c_s^2k^2-(\omega+i0)^2},\qquad
\boxed{\chi_{nn}^{R}(0,k)=\chi_n,}
\end{equation}
not $1/k^2$. Its propagating pole is not an unscreened static gravitational mass-source pole. The relevant limit is \emph{first} $\omega=0$ for a static source and \emph{then} $k\rightarrow0$; interchanging limits can change the answer.

\subsection{Mixed-channel existence theorem: the source need not couple directly to the gapless field [D/C]}\label{sec:v09-mixing}
The selection theorem suggests a route less restrictive than demanding that the long-range collective coordinate itself directly overlap with a microscopic density bilinear. Consider two \emph{hypothetical} collective coordinates, a gapless $q$ and a massive source-facing $s$, with positive constants $Z,Z_s,M_*,\lambda,g_s$ and stationary energy
\begin{equation}\label{eq:v09-two-field-energy}
\boxed{\mathfrak E[q,s;\rho_m]=\int\dd^3x\left[\frac Z2|\bm\nabla q|^2+\frac{Z_s}{2}|\bm\nabla s|^2+\frac{M_*^2}{2}(s-\lambda q)^2-g_s\rho_m s\right].}
\end{equation}
\textbf{[A]} The full functional and the joint shift symmetry $(q,s)\mapsto(q+\epsilon,s+\lambda\epsilon)$ of its \emph{source-free} part are supplied for this toy test, not yet generated from Chapter 1's four-fermion action. In particular, no externally assumed QGD or Einstein metric is used in solving the model. The source acts \emph{only} on the otherwise massive $s$ channel. The Fourier Hessian, source vector and its determinant are
\begin{equation}\label{eq:v09-two-field-K}
K(k)=\begin{pmatrix}Zk^2+\lambda^2M_*^2&-\lambda M_*^2\\ -\lambda M_*^2&M_*^2+Z_s k^2\end{pmatrix},\quad B=\begin{pmatrix}0\\ g_s\end{pmatrix},\quad
\boxed{\det K=k^2\left[M_*^2(Z+\lambda^2 Z_s)+ZZ_s k^2\right].}
\end{equation}
An exact inversion gives
\begin{equation}\label{eq:v09-mixed-response}
\boxed{\frac{\delta q(k)}{\delta\rho_m(k)}=\frac{\lambda g_s}{Z_{\rm eff}\,k^2(1+\ell^2k^2)},\qquad
Z_{\rm eff}=Z+\lambda^2Z_s,\qquad
\ell^2=\frac{ZZ_s}{M_*^2Z_{\rm eff}}.}
\end{equation}
\emph{The mass source directly excites the $s$ channel; mode mixing transfers a nonzero residue into the joint gapless mode.} The full static $K(k)$ is positive for every $k\neq0$ when the displayed coefficients are positive; the only long-wavelength flat direction is the deliberately supplied joint shift mode. Thus the calculation shows that a gapped observable may transmit source information to a long-range mode, not that a massless mode can emerge from a uniformly positive Hessian.

For an idealized point mass $\rho_m(\bm x)=M_{\rm src}\delta^{(3)}(\bm x)$, the \emph{independently Fourier-inverted} result is
\begin{equation}\label{eq:v09-pointfield}
\boxed{q(r)=\frac{\lambda g_s M_{\rm src}}{4\pi Z_{\rm eff}r}\left(1-e^{-r/\ell}\right).}
\end{equation}
If the effective null-sector observable is \emph{conditionally} identified as $H=2g_{\rm eff}q/c^2$, with $g_{\rm eff}=\lambda g_s$, then
\begin{equation}\label{eq:v09-G-and-H}
\boxed{G_{\rm eff}=\frac{g_{\rm eff}^2}{4\pi Z_{\rm eff}},\qquad
H(r)=\frac{2G_{\rm eff}M_{\rm src}}{c^2r}\left(1-e^{-r/\ell}\right).}
\end{equation}
In the controlled exterior regime $r\gg\ell$ this reproduces the Chapter 5 conditional $H\sim r^{-1}$ and the positive null-sector amplitude $\sigma_{\rm amp}\sim r^{-1/2}$, while giving an explicit exponentially short-range deviation at finite $r$. The latter should not be fitted to observations until $\ell$ is derived and constrained from an actual matter theory.

\begin{figure}[htbp]\centering
\includegraphics[width=.83\linewidth]{qm015_mixed_mode_response.png}
\caption{Exact QM015 toy-model radial ratio to the asymptotic QGD/Newtonian null-sector field. The finite-range correction follows from the two-field Hessian, not an empirical fit. There is no claim that the original QGD microscopic theory produces these parameters.}\label{fig:v09-mixed}
\end{figure}

\textbf{Independent check.} In dimensionless units $Z=2$, $Z_s=3$, $M_*=4$, $\lambda=1/2$, $g_s=1.2$, $c=1$, we obtain $Z_{\rm eff}=2.75$, $\ell=0.369274472938$, $g_{\rm eff}=0.6$, $G_{\rm eff}=0.0104174144569$. The same $q(r)$ is obtained to better than $2\cdot10^{-7}$ absolute accuracy by direct numerical three-dimensional Fourier inversion for $M_{\rm src}=3$ and $r=0.1,0.5,1,3,8$, independently of the analytic Yukawa decomposition. The exact determinant and susceptibility identity are also checked symbolically; full verifier and data are packaged with this edition.

\subsection{Schur-complement version of the mechanism [D]}\label{sec:v09-schur}
For a general coupled physical block $\varphi=(q,s)$ with $K_{ss}$ invertible at the chosen momentum, eliminating $s$ yields
\begin{equation}\label{eq:v09-schur}
K_{q,\mathrm{red}}=K_{qq}-K_{qs}K_{ss}^{-1}K_{sq},\qquad
B_{q,\mathrm{red}}=B_q-K_{qs}K_{ss}^{-1}B_s.
\end{equation}
These formulae expose an important difference between two questions: whether a channel is itself gapless, and whether it becomes \emph{source-visible by mixing}. Even when $B_q=0$ at the bare level, the reduced vertex can be nonzero if $K_{qs}B_s\neq0$. In the specific model above $K_{qs}=-\lambda M_*^2$, so $B_{q,\mathrm{red}}(0)=\lambda g_s$. The effective massless stiffness is $Z_{\rm eff}$. The route to QGD would therefore be a \emph{microscopically derived Schur complement} whose eigenmode and source overlap satisfy the static QGD target. It is not necessary to postulate a direct $-g_c q\rho$ term in an original fermionic Lagrangian, but the interaction that generates the mixing must still be derived.

\subsection{A controlled candidate microscopic experiment: exact Hubbard--Stratonovich identity, then 2PI selection [O]}\label{sec:v09-2pi-test}
A viable research path retains the actual matter degrees of freedom instead of inserting $g(\Sigma)$ as their initial background. For illustration, an $N$-species flat-background Euclidean fermion theory with a scalar contact operator $\mathcal O=\sum_{i=1}^N\bar\psi_i\psi_i$ has the \emph{algebraic} Hubbard--Stratonovich identity
\begin{align}\label{eq:v09-hs-identity}
\exp\left[+\frac{G_0}{2N}\int\mathcal O^2\right]
&\propto\int\mathcal Dq\,\exp\left[-\int\left(\frac{Nq^2}{2G_0}-q\mathcal O\right)\right],\\
\Gamma_{\rm HS}[q]&=\frac{N}{2G_0}\int q^2-N\,\Tr\log(D_0-q)+\text{counterterms},
\end{align}
where the second line assumes identical species and a regulator. Here the $q$ integration is initially an \emph{exact rewriting of the chosen contact interaction}, not premise (Q), and is not yet a universal metric variable. Expanding about a true self-consistent saddle produces $K_q(k)=K_q(0)+Z_qk^2+\cdots$ with the real mass-source vertex computed by inserting a \emph{separate sourced microscopic matter species} or the appropriate conserved energy-stress operator. The already performed QM011 model yields $K_q(0)>0$ under the two tested regulator schemes; nothing in the HS identity forces a pole.

The next numerical experiment should therefore calculate, \emph{from one consistent regulated fermion theory}, the combined basis $\{\bar\psi\psi,\,\bar\psi i\gamma_5\psi,\,T^{00},\,T^{0i},\,T^{ij},\,\partial_{p_\mu}\Sigma_{\rm Dirac},\ldots\}$, renormalize its composite mixing, and compute its connected source response through a conserving 2PI/Bethe--Salpeter differentiation. A finite truncation can preserve specified global conservation laws without automatically satisfying every gauge Ward identity; the exact identities and truncation dependence must be separately checked. The trial QGD observable $H_{\rm micro}$ should be defined operationally, not by already inserting $g_{\mu\nu}(\Sigma)$ into the fermion kinetic term.

\textbf{Concrete decision procedure.} (i) Solve the source-free homogeneous phase with the matched Chapter 5/QM011 regulator and record all stationary branches and thermodynamic stability. (ii) Add a second, spatially localized matter source with independently conserved stress; obtain $B_A$ by differentiating the same functional. (iii) Compute $K_{AB}(\omega,\bm k)$, its Ward-completed contacts and its constrained static Schur complement with $\omega=0$ before $k\to0$. (iv) Determine whether an eigenvalue is genuinely $Zk^2$, whether $A\cdot u$ and $u\cdot B$ are nonzero, and whether the residue is positive and stable under regulator/volume change. (v) Only then infer $H_{\rm micro}$ and compare the independently computed $\partial_{\rho} \partial_{p_\mu}\Sigma_{\rm Dirac}$ with the common kinetic deformation required by the QGD null-sector Dirac operator. (vi) Test separate massive, massless, polarized, antiparticle and source-composition probes and ultimately the full coframe/signed-sector/$\mathcal Q$ metric rather than a single scalar fifth force.

\subsection{Why even a successful scalar pole would not yet be gravity [D/C/O]}\label{sec:v09-universality}
A scalar $1/k^2$ pole suffices to generate a Newtonian-looking static force but is not by itself a proof of universal relativistic gravitation. In particular, a mere scalar mass self-energy shifts the Dirac mass term but need not change the principal symbol $\gamma^\mu p_\mu$; in the original QM009 tested Hartree--Fock channel the source-to-kinetic vertex is zero. A purely conformal perturbation of a four-dimensional reference metric also preserves unparameterized null geodesics, so its lensing cannot be assumed equivalent to the nonconformal QGD Kerr--Schild lift. The correct target must include the microscopic source-dependent \emph{derivative} vertex and propagation of more than one probe species.

For the QGD positive null-sector specialization, define the classical geometric representation only \emph{after} a physical response is obtained,
\begin{equation}\label{eq:v09-lift}
H_{\rm micro}\quad\xrightarrow[\text{additional physical selection}]{\text{nondegenerate sector/orbit lift}}\quad
[\Sigma],\qquad g_{\mu\nu}=\eta_{\mu\nu}+H\ell_\mu\ell_\nu,\quad \sigma_\mu=\sqrt{H}\ell_\mu.
\end{equation}
The lift and the vector $\ell_\mu$ cannot be defined by the motion of the very test particle used to validate the field. Moreover the founding complete metric remains
\begin{equation}\label{eq:v09-full-metric}
\boxed{g_{\mu\nu}=T^a{}_{\mu}T^b{}_{\nu}\left[\eta_{ab}-\sum_s\varepsilon_s\sigma_a^{(s)}\sigma_b^{(s)}-\lambda\mathcal Q_{ab}\right],\qquad s_s=-\varepsilon_s.}
\end{equation}
The $\lambda\mathcal Q=0$, fixed-reference Kerr--Schild relation is a declared specialization, not the general microscopic QGD theory.

\section{Final v0.9 proof ledger, options, and next calculation}\label{sec:v09-decision}
\begin{center}\small
\begin{tabularx}{\linewidth}{@{}>{\raggedright\arraybackslash}p{.30\linewidth}>{\raggedright\arraybackslash}p{.11\linewidth}>{\raggedright\arraybackslash}X@{}}\toprule
Result or candidate&Status&Scientific meaning\\\midrule
Analytic positive physical static Hessian forbids a $1/k^2$ pole&[D]&No gapless long-range Gaussian response in that fixed phase\\
Source-overlap residue $(A\cdot u)(u\cdot B)/Z$&[D]&Necessary observable and source couplings, in addition to a pole\\
Parity-even scalar source versus parity-odd massless mode&[D/C]&Gapless collective modes may be physically invisible to mass\\
Hydrodynamic phonon versus static density response&[D/C]&A propagating gapless phonon does not imply divergent static density susceptibility\\
Two-field source-visible massless mode by mixing&[D/C]&Explicit stationary existence model with an exact $1/r$ exterior and predicted finite-$r$ correction\\
Microscopic origin of the two-field shift symmetry and mixing&[O]&Coefficients/fields remain postulated in QM015 toy; not a Dirac derivation\\
Exact HS transformation of a chosen fermion contact action&[D/C]&Does not itself add premise (Q); dynamically generated pole still open\\
QGD probe-independent $H\to[\Sigma]$ and universal kinetic coupling&[O]&Requires independent dynamical/operator proof\\
Claude's independent induced-$g(\Sigma)$ route (Q)&[A/O]&Separate unadopted additional microscopic premise with open constraints\\\bottomrule
\end{tabularx}
\end{center}
\textbf{Recommended next concrete calculation.} Re-use the original Chapter 1 microscopic fermion contact action and QM011 matched regulators, enlarge the HS/2PI collective basis to the minimal scalar--stress--kinetic channels, and \emph{measure the complete static source-overlap residue} at its self-consistent stable phase(s). Vary the contact coupling through any legitimate critical phase; do not set $K(0)=0$ by hand. If a protected gapless sector is found, explicitly test whether the physical source vertex vanishes by symmetry or kinematics, and whether a mixed mode supplies a nonzero finite overlap. Verify that the same micro theory yields the source-dependent Dirac principal symbol. Negative answers would be publishable restrictions on this mechanism even if the target pole is absent.

\textbf{Precise research outcome.} QM015 contributes a source-overlap theorem, exact mixed-mode transfer mechanism and two decisive false-positive controls. It \emph{does not} derive a long-range gravitational mode or any new fundamental field from the existing microscopic QGD interaction. This preserves the Chapter 5 conditional closure and makes Chapter 6's strongest unresolved task more concrete: derive a source-visible, stable physical collective operator and its universal propagation effects directly from quantum matter before promoting it to the natural-invariant sector. All prior v0.7/v0.8 labels, calculations and proof statuses are retained.

\textbf{Literature supporting methodological comparators.} G. Baym and L. P. Kadanoff, \emph{Phys. Rev.} \textbf{124}, 287 (1961), \url{https://doi.org/10.1103/PhysRev.124.287}; J. M. Cornwall, R. Jackiw and E. Tomboulis, \emph{Phys. Rev. D} \textbf{10}, 2428 (1974), \url{https://doi.org/10.1103/PhysRevD.10.2428}; overview of NJL scalar/pseudoscalar collective channels, \url{https://doi.org/10.1088/1674-1137/43/9/094103}; Weinberg soft-graviton universality (1965), \url{https://doi.org/10.1103/PhysRev.140.B516}; J. Bonifacio \emph{et al.}, \emph{JHEP} \textbf{07}, 056 (2020), \url{https://doi.org/10.1007/JHEP07(2020)056}. The cited literature provides established comparators, not an independent endorsement of QGD's unproved microscopic origin.
% ===== END v0.9 QM015 extension =====

% ===== BEGIN v1.0 QM016 extension; v0.9 preserved =====
\section{QM016: a genuinely microscopic source--scalar--density block and a strict Hartree obstruction}\label{sec:v10-qm016}
\textbf{Scientific scope.} This is a new calculation within the \emph{original Chapter~1} scalar-contact Dirac model of Eq.~\eqref{eq:micro}, not the freely chosen two-field energy of QM015. We analyze the self-consistent $T=0$ homogeneous Hartree phase at fixed $\mu$, $m_0$ and sharp spatial cutoff $\Lambda$, then test the two previously matched smooth finite-wave-number regulators at their benchmark point. The independent fields being varied are the scalar Hubbard--Stratonovich mass coordinate $M$ and its conjugate number-density source; \emph{neither} is identified with the QGD gravitational $H$ or the complete tensor stress of a separate source. Proof status [D/C]; the universal microscopic-to-natural-invariant map stays [O].

\subsection{A positive-curvature identity that excludes tuning a scalar massless pole in a controlled parameter domain [D/C]}\label{sec:v10-positive}
Recall the stationary gap and static scalar Hessian from Eqs.~\eqref{eq:gap} and \eqref{eq:K0}, with $E_p=\sqrt{p^2+M^2}$, $M>0$, $G_4>0$, $m_0>0$, $\Lambda>\mu>0$ and $p_F=\sqrt{\max(\mu^2-M^2,0)}$. The exact algebraic identity
\begin{equation}\label{eq:v10-gap-identity}
\frac{p^2}{E_p}-\frac{p^4}{E_p^3}=\frac{M^2p^2}{E_p^3}
\end{equation}
permits elimination of $1/(2G_4)$ using the \emph{same self-consistent} gap equation, not a gravitational matching or a fitted kernel:
\begin{equation}\label{eq:v10-Kgap}
\boxed{K_0=\frac{m_0}{2G_4M}+\frac{M^2}{\pi^2}\left[\int_{p_F}^{\Lambda}\frac{p^2\,\dd p}{(p^2+M^2)^{3/2}}-\frac{p_F}{\mu}\right].}
\end{equation}
For an occupied Fermi sea, $0<M<\mu$, its bracket is the explicit function
\begin{align}\label{eq:v10-F-def}
F(M;\Lambda,\mu)&=\log\!\left(\frac{\Lambda+\sqrt{\Lambda^2+M^2}}{\mu+\sqrt{\mu^2-M^2}}\right)-\frac{\Lambda}{\sqrt{\Lambda^2+M^2}},\\
\frac{\partial F}{\partial M}
&=\frac{M}{E_\Lambda(\Lambda+E_\Lambda)}+\frac{M}{p_F(\mu+p_F)}+\frac{\Lambda M}{E_\Lambda^3}>0,
\quad E_\Lambda=\sqrt{\Lambda^2+M^2},\label{eq:v10-F-prime}\\
\lim_{M\to0^+}F(M;\Lambda,\mu)&=\log(\Lambda/\mu)-1.\label{eq:v10-F-limit}
\end{align}
Here $\partial_M F>0$ is evaluated at fixed $\Lambda,\mu$ and $0<M<\mu$. The derivative was also checked by exact symbolic differentiation. For an unoccupied Fermi sea, $M\ge\mu$, $p_F=0$ and the square bracket in Eq.~\eqref{eq:v10-Kgap} is manifestly positive.
\begin{proposition}[Scalar Hartree criticality exclusion in the stated cutoff domain]\label{prop:v10-positivity}
For the homogeneous positive-mass solutions of the sharp-cutoff Hartree gap equation with $m_0>0$, $G_4>0$ and $\Lambda/\mu\ge e$ (and $\Lambda>\mu>0$), the zero-wave-number scalar curvature is \emph{strictly positive}, $K_0>0$, for any self-consistent value of $G_4$. This excludes a zero-momentum massless scalar \emph{within this model, regulator and approximation}; it makes no claim about nonlocal self-energy, stress-tensor eigenmodes, non-Gaussian criticality, different cutoff ratios, or other microscopic interactions.
\end{proposition}
\begin{proof}
For $M<\mu$, Eqs.~\eqref{eq:v10-F-prime} and \eqref{eq:v10-F-limit} give $F(M)\ge\log(\Lambda/\mu)-1\ge0$. The separate term $m_0/(2G_4M)$ is strictly positive. For $M\ge\mu$, the integral in Eq.~\eqref{eq:v10-Kgap} is positive, as is the first term. Thus $K_0>0$ for each such self-consistent solution. No local-minimum assumption is needed for the algebraic conclusion, although thermodynamic global selection is a separate question.
\end{proof}
This result is stronger than sweeping $G_4$ numerically; unlike a generic ``no-go'' claim, it has an explicit parameter-domain proof. At our $\Lambda=4$, $\mu=1.2$ benchmark, $\Lambda/\mu=10/3>e$ and the positive lower bound on $F$ is $\log(10/3)-1\simeq0.203973$. The original $m_0=0.12$ is held fixed. As emphasized in QM009--QM011, a hard spatial cutoff is \emph{not} accepted as a reliable finite-$q$ regulator; this theorem concerns only the self-consistent $q=0$ Hartree thermodynamics.

\subsection{The genuine two-coordinate scalar--density static Hessian [D/C]}\label{sec:v10-legendre}
A physical number-density response cannot be checked by calling a chemical-potential source $\delta\mu$ an independent propagating boson. Introduce the Legendre-transformed static Helmholtz density $\mathcal F(M,n)=\Omega(M,\mu)+\mu n$ with $n=-\Omega_\mu$ in the occupied phase. Using the \emph{same} fermionic grand-potential derivatives
\begin{equation}\label{eq:v10-thermo-derivs}
K=\Omega_{MM},\qquad V=\Omega_{M\mu}=\frac{Mp_F}{\pi^2},\qquad
D=-\Omega_{\mu\mu}=\frac{\mu p_F}{\pi^2}>0,
\end{equation}
and eliminating $\delta\mu$ at fixed $\delta n$ yields the symmetric two-coordinate physical static Hessian
\begin{equation}\label{eq:v10-hessian}
\boxed{\mathsf H_{(M,n)}=
\begin{pmatrix}K+V^2/D&V/D\\ V/D&1/D\end{pmatrix},\quad
\det\mathsf H_{(M,n)}=\frac KD>0,\quad
\chi_{nn}=D+\frac{V^2}{K}.}
\end{equation}
Its upper-left principal minor is positive and its determinant is positive; hence both thermodynamic eigenvalues are positive where $D>0$. Number density is the Legendre conjugate of $\mu$, not a separately proven propagating gapless field. In the empty-Fermi-sea branch $D=V=0$ at $T=0$, the Legendre matrix in Eq.~\eqref{eq:v10-hessian} is singular as a \emph{coordinate choice} and must not be interpreted as a new massless collective mode; the actual number-source response vanishes.

The static source response therefore remains finite for the tested finite-density phase. If the finite-wave-number continuation has the regularity independently established for the QM009 matched scalar channel, the proposed scalar source-visible pole residue obeys
\begin{equation}\label{eq:v10-residue}
\boxed{\lim_{|\bm k|\to0}|\bm k|^2\left[D(\bm k)+\frac{V(\bm k)^2}{K(\bm k)}\right]=0.}
\end{equation}
The constant-density susceptibility and its $k\to0$ order are essential: a conserved density can have a hydrodynamic sound excitation at nonzero frequency while its \emph{static} compressibility remains finite. The compressibility sum-rule control is standard in self-consistent NJL response theory; see \url{https://doi.org/10.1088/1674-1137/43/9/094103}.

\subsection{New self-consistent coupling sweep and independent thermodynamic sum-rule check [D/C]}\label{sec:v10-numerical}
For $m_0=0.12$, $\mu=1.2$, $\Lambda=4$ and the original sharp zero-momentum normalization, we re-solved the \emph{nonlinear} gap equation at 151 distinct couplings $G_4\in[0.25,0.95]$ rather than inserting the precomputed mass. Only one positive-mass root was found at each scanned coupling using the stated root-search grid. The populated branch's minimum scalar curvature was $K_{\min}\simeq0.14633658$ near $G_4\simeq0.60048147$. The Fermi surface closes at
\begin{equation}\label{eq:v10-Gstar}
\boxed{G_4^*=\frac{\pi^2(\mu-m_0)}{2\mu\int_0^\Lambda p^2\,\dd p/\sqrt{p^2+\mu^2}}=0.637146968076\ldots.}
\end{equation}
When $G_4>G_4^*$ on the tested homogeneous branch, $M>\mu$, so $p_F=0$ and $V=0$; the charge/number source then does not excite the scalar at zero temperature. Thus increased scalar attraction does not open a long-range \emph{number-source-visible} channel here: throughout the occupied branch $K_0>0$, and beyond its endpoint the source overlap vanishes. The analytical positivity proposition, not the finite scan, is the warrant for the stated sharp-regulator domain.

At the original benchmark $G_4=0.5654255716$, the recomputed values are $M=0.693435535520$, $K_0=0.168741734136$, $V_0=0.068809590036$, $D_0=0.119075968585$, and $\chi_{nn}=0.147135178153$. Independently re-solving the gap equation at $\mu\pm10^{-4}$ and differentiating $n=p_F^3/(3\pi^2)$ gives $\chi_{nn}=0.147135178036$, agreeing within $1.2\times10^{-10}$. The algebraic $K_0$ identity also matches the direct fluctuation-integral Hessian along the scanned branch to better than $2\times10^{-11}$ absolute numerical tolerance.

The two \emph{previously independently matched} QM009 smooth static infrared regulator benchmarks provide an additional, different check of the same thermodynamic source block, not a new smooth-regulator coupling sweep:
\begin{center}\small
\begin{tabular}{@{}lrrrr@{}}\toprule
Regulator & $K_s(0)$ & $D_0$ & $\chi_{nn}$ & $\lambda_{\min}(\mathsf H_{(M,n)})$\\\midrule
Smooth $n=8$ &0.14993818&0.11907597&0.15065405&0.14921889\\
Smooth $n=16$ &0.15737518&0.11907597&0.14916178&0.15661953\\\bottomrule
\end{tabular}
\end{center}
Both use $V_0=0.06880959$ and have positive $\det\mathsf H=K_s/D_0$. They therefore reinforce the positive scalar--density static conclusion at the tested reference state. Their finite-$q$ stiffnesses are still those of QM009; an independent full matrix $K_{AB}(\omega,\bm k)$ with stress and kinetic vertices was \emph{not} computed here.
\begin{figure}[htbp]\centering
\includegraphics[width=.82\linewidth]{qm016_scalar_density_sweep.png}
\caption{QM016 nonlinear gap-equation coupling sweep for the Chapter 1 sharp $q=0$ contact model. The scalar curvature stays positive across the tested interval. When the Fermi surface closes at $G_4^*$, the number-source scalar overlap goes to zero. The analytic positivity theorem applies to the specified sharp-cutoff domain even beyond the finite sweep; the plot is a numerical cross-check, not a demonstration of universal quantum gravity.}\label{fig:v10-sweep}
\end{figure}

\subsection{Why this obstructs the QM015 mixed-mode mechanism in the tested truncation [D/C/O]}\label{sec:v10-missing}
QM015's exact long-range \emph{toy} mechanism had a second, shift-symmetric collective coordinate and an off-diagonal mixing which transfers a number-source response into its zero eigenvalue. The original scalar-contact Hartree functional \emph{does not} dynamically generate that extra coordinate: the scalar stiffness remains nonzero, and the independent density Legendre coordinate has positive curvature whenever occupied. The $(M,n)$ two-coordinate response is therefore \emph{not} a realization of the QM015 massless-mixing theorem. No Ward symmetry of this microscopic interaction protects the required $1/k^2$ scalar pole. In the occupied, isotropic phase the scalar-to-helicity-two stress mixing also vanishes by the earlier QM007 rotational selection rule; introducing the density coordinate does not manufacture the missing TT excitation.

Within the local Hartree--Fock truncation, $\partial_{p_\mu}\mathcal D_{\rm HF}=\gamma^\mu$ under a source-state perturbation even when $M(x)$ changes; Eq.~\eqref{eq:qm009-zero} still gives a vanishing source-to-kinetic principal-symbol deformation. This is a \emph{separate and indispensable} obstruction to universal QGD-like propagation, not something repaired merely by increasing the static scalar susceptibility. These statements do not extend to the full nonlocal $\Sigma_{\rm Dirac}(p,x)$ at higher skeleton/2PI order.

\textbf{Physical-source caveat.} The perturbation $\delta\mu$ is a conserved-number external source \emph{inside the chosen finite-density fermion medium}. It is not a compact object's full independent stress tensor, and the homogeneous medium is not an isolated asymptotically empty exterior. A second matter species $\chi$ which factorizes exactly from the original $\psi$ action has zero connected $\chi$--$\psi$ cross-response, so an independent-source coupling must be generated or specified in the microscopic theory; the $\psi$ density susceptibility cannot silently be renamed the universal mass-source vertex. This distinction is necessary for an honest $\rho_m\to[\Sigma]$ selection test.

\subsection{QM016 decision and the next actually different microscopic calculation [O]}\label{sec:v10-decision}
\begin{enumerate}[label=\textbf{QM017-\arabic*.},leftmargin=18mm,itemsep=3pt]
\item Construct a \emph{single consistent} 2PI $\Phi[G]$ or Bethe--Salpeter kernel at the next genuinely momentum-dependent order for the Chapter 1 interaction. The new physics must be a frequency- and momentum-dependent Dirac self-energy/vertex, not merely re-diagonalization of the already-positive $(M,n)$ static block.
\item Couple independent conserved stress sources as \emph{bookkeeping probes of the original flat-space QFT}, including the needed stress-contact/seagull insertions; do not promote their metric source into an independently integrated field by Claude's unadopted premise (Q). Evaluate the retarded constrained $(T^{00},T^{0i},T^{ij},S,J^\mu,\partial_p\Sigma)$ response and the helicity decompositions.
\item Verify the static density compressibility Ward/sum-rule and the energy-momentum Ward identities from the same functional approximation. A $\Phi$-derivable approximation enforces specified conservation laws but does not automatically satisfy all Ward--Takahashi identities for truncated internal vertices; document the external conserving vertex separately. See Baym--Kadanoff, \url{https://doi.org/10.1103/PhysRev.124.287}, and van Hees--Knoll, \url{https://doi.org/10.1103/PhysRevD.66.025028}.
\item Test an isolated, positive-residue, source-visible infrared pole \emph{and} a nonzero universal source-to-Dirac-principal-symbol response. Only after both checks should we attempt the source/probe-independent natural-invariant lift and the entire complete coframe/$\mathcal Q$ QGD metric.
\end{enumerate}
\textbf{Conclusion of QM016.} A newly derived analytic Hartree curvature identity excludes tuning the scalar to a massless source-visible pole for $\Lambda/\mu\ge e$, $m_0>0$ and the other stated conditions. The full static two-coordinate number--scalar thermodynamic Hessian is positive on the occupied branch, while the source overlap vanishes as the Fermi sea empties. The original Chapter~1 contact theory has therefore \emph{not} supplied the gapless source-visible mixed collective mode assumed only in QM015; the earlier conditional Schwarzschild theorem remains intact. The next relevant experiment must introduce \emph{new dynamics computed from the same microscopic theory}, not a relabeling of the finite Hartree susceptibility as gravity.
% ===== END v1.0 QM016 extension =====

% ===== BEGIN QM017--QM019: sustained microscopic-to-universal-geometry audit =====
\newpage
\section{QM017--QM019: Does microscopic matter produce a universal kinetic geometry?}\label{sec:v11-qm017-019}
\noindent\textbf{Date: 2 November 2026. Research status: mixed constructive result and a decisive failure of the declared minimal mechanism; not a derivation of gravitational emergence.} This section follows QM016 through multiple tests rather than stopping at a positive self-energy coefficient. We maintain the historical target
\begin{equation}\label{eq:v11-original}
\mathcal D_{\rm eff}\Psi=0\longrightarrow(\Psi,S,\mathsf R,j^\mu)
\overset{\mathcal F_{\rm micro}}{\longrightarrow}[\Sigma],
\end{equation}
and use the operational source formulation
\begin{equation}\label{eq:v11-operational}
\boxed{(S_{\rm micro},\hat\varrho_{\rm source},{\rm BC})
\longrightarrow g_{\mu\nu}^{\rm physical}\longrightarrow[\Sigma]_{\rm physical}.}
\end{equation}
The second arrow is an \emph{algebraic} local equivalence on a regular, fixed-coframe, $\lambda\mathcal Q=0$, spanning signed-sector branch. The first arrow requires independent physical dynamics, a source-supported exterior and cross-species universality.

\subsection{QM017A: a concrete nonlocal microscopic kinetic response, and its scope [D/C]}\label{sec:v11-exchange}
The Chapter~1 benchmark at $M=0.693435535520$, $\mu=1.2$ has $p_F=\sqrt{\mu^2-M^2}=0.979360586341$. QM011 independently obtained the \emph{infrared} static HS scalar Hessian $K_s(k)=K_0+Zk^2+O(k^4)$ for matched smooth regulators. As an explicitly defined screened-exchange \emph{diagnostic}, continue this quadratic expression into the internal momenta of a positive-energy Fermi sea:
\begin{equation}\label{eq:v11-screen}
 D_s(\bm q)=\frac1{K_0+Z\bm q^2},\qquad
 \xi=\sqrt{Z/K_0},\quad K_0,Z>0.
\end{equation}
The use of this $D_s$ at all $|\bm q|\le2p_F$ is a \emph{model extension}, not an independently computed full-$q$ Dirac/HS propagator. In particular, QM011's low-$k$ fit does not fix the complete exchange integral. The following coefficient is therefore a reproducible constructive proxy, \emph{not} a self-consistent 2PI prediction of the original contact action.

For a scalar Yukawa normalization $g_*^2$ (set to one only for the numerical table), take the static positive-energy exchange self-energy, in an explicitly specified rest-frame Dirac convention,
\begin{equation}\label{eq:v11-Fock}
\Sigma_F(\bm p)=g_*^2\int_{|\bm k|<p_F}\frac{d^3k}{(2\pi)^3}
 \frac{\gamma^0 E_k-\bm\gamma\!\cdot\!\bm k+M}{2E_k}
 \frac1{K_0+Z|\bm p-\bm k|^2},\qquad E_k=\sqrt{k^2+M^2}.
\end{equation}
The overall Fock sign and the normalization of the effective auxiliary Yukawa vertex are fixed \emph{here by definition}; they cannot be imported as an uncomputed sign or coupling of a covariant 2PI skeleton. Different microscopic conventions can reverse the displayed sign of the kinetic correction, not its nonzero magnitude.

\begin{proposition}[Nonzero spatial kinetic derivative in the screened-exchange diagnostic [D/C]]\label{prop:v11-Cs}
The odd-in-$\bm p$ spatial part of Eq.~\eqref{eq:v11-Fock} has
\begin{equation}\label{eq:v11-Cs}
\boxed{\Sigma_F(\bm p)\big|_{\bm\gamma\cdot\bm p}=-C_s\,\bm\gamma\cdot\bm p+O(|\bm p|^3),\quad
 C_s=\frac{g_*^2Z}{6\pi^2}\int_0^{p_F}\!dk\,
 \frac{k^4}{E_k(K_0+Zk^2)^2}>0.}
\end{equation}
At fixed $M,K_0,Z$, its \emph{partial}, not fully self-consistent, density response is
\begin{equation}\label{eq:v11-partialsource}
\boxed{\left.\frac{\partial C_s}{\partial\mu}\right|_{M,K_0,Z}
 =\frac{g_*^2 Zp_F^3}{6\pi^2(K_0+Zp_F^2)^2}>0,\quad
 \left.\frac{\partial C_s}{\partial n}\right|_{M,K_0,Z}
 =\frac{g_*^2 Zp_F^2}{6\mu(K_0+Zp_F^2)^2},\;n=\frac{p_F^3}{3\pi^2}.}
\end{equation}
\end{proposition}
\begin{proof}
At $\bm p=0$, $\partial D_s(|\bm p-\bm k|)/\partial p_j=2Zk_j/(K_0+Zk^2)^2$. The time-like and scalar numerators integrate to zero after multiplication by the odd function $k_j$. The spatial numerator $-\gamma^i k_i/(2E_k)$ gives $-g_*^2Z\gamma^i p_j\int d^3k\,k_i k_j/[ (2\pi)^3 E_k(K_0+Zk^2)^2]$. Rotational integration substitutes $k_i k_j\mapsto\delta_{ij}k^2/3$, proving Eq.~\eqref{eq:v11-Cs}. Differentiate the moving upper limit $p_F=\sqrt{\mu^2-M^2}$ while holding the other parameters fixed to prove Eq.~\eqref{eq:v11-partialsource}. The fully physical derivative also includes $M'(\mu)$, $K_0'(\mu)$, $Z'(\mu)$ and the appropriate vertex corrections, which are \emph{not} included here.
\end{proof}

\begin{center}\small
\begin{tabular}{@{}lrrrr@{}}\toprule
QM011 scheme&$K_0$&$Z$&$\xi$&$C_s$ at $\kappa=p_F$\\\midrule
Smooth $n=8$&0.14993818&0.04864381&0.56958383&0.00422660\\
Smooth $n=16$&0.15737518&0.08155175&0.71986081&0.00532135\\\bottomrule
\end{tabular}
\end{center}
At $\kappa=0.2p_F$, the same coefficients are $2.9039\times10^{-6}$ and $4.3733\times10^{-6}$, respectively, so an actual nonzero contribution is already present in the conservative small-internal-momentum part. At the full $p_F$ the quoted values depend on continuing the $O(k^2)$ HS kernel beyond the regime where QM011 controlled the fit. Direct Gauss--Legendre angular integration of Eq.~\eqref{eq:v11-Fock}, followed by a symmetric difference in the external $p_z$, agrees with Eq.~\eqref{eq:v11-Cs}. The separately differentiated upper-limit response agrees with a numerical $\mu$ derivative at fixed $M,K_0,Z$. The verifier supplies every test and tolerance.
\begin{figure}[htbp]\centering
\includegraphics[width=.86\linewidth]{qm017_kinetic_exchange.png}
\caption{QM017 screened exchange: a nonzero medium kinetic coefficient appears when an already screened HS collective channel is permitted to carry momentum. This is a controlled diagnostic of a \emph{specified} low-momentum extension, not a full 2PI computation or a gravitational source measurement.}\label{fig:v11-exchange}
\end{figure}

\subsection{QM017B: continuation to the actual exterior response fails [D/C]}\label{sec:v11-range}
The same positive HS Hessian has the exact three-dimensional \emph{gradient-model} Green function
\begin{equation}\label{eq:v11-yukawa}
\boxed{D_s(r)=\frac{e^{-r/\xi}}{4\pi Zr},\qquad
 \lim_{k\to0}k^2D_s(k)=0.}
\end{equation}
For the two QM011 matched schemes, $\xi=0.56958$ and $0.71986$ in model inverse-mass units. Hence this propagator does \emph{not} deliver the $1/k^2$ source-response residue or asymptotic $1/r$ exterior required by the Chapter~5 static QGD target. It can create finite state-dependent kinetic dressing inside a medium without selecting a distant gravitational field. A real gapless Fermi sea may also exhibit nonanalytic dynamic particle--hole and $2p_F$ effects beyond this gradient approximation; they are not automatically an isolated gravitational pole.
\begin{figure}[htbp]\centering
\includegraphics[width=.80\linewidth]{qm017_range_test.png}
\caption{Distance test: screened QM011 scalar-gradient model versus the conditional QGD $1/r$ response, with the profiles normalized at $r=1$. The scalar does not become gravitational merely because its exchange also perturbs a Dirac dispersion.}\label{fig:v11-range}
\end{figure}

\paragraph{A necessary technical distinction.} The derivative $\partial\Sigma/\partial p_i$ in Eq.~\eqref{eq:v11-Cs} modifies a \emph{local low-energy expansion around a chosen momentum}. This does not automatically replace the exact ultraviolet/pseudodifferential principal symbol, and a momentum-dependent nonlocal self-energy may not be describable by one local metric. A frequency-dependent time-like coefficient and consistent retarded vertex are missing. Thus the result disproves a too-strong assertion that matter interactions can \emph{never} alter kinetic coefficients, but it does not meet QGD's universal local-geometry target.

\subsection{QM018: exact all-orders obstruction to universal response of a decoupled spectator [D]}\label{sec:v11-factorization}
Let $\psi$ describe the source and let $\chi$ be an independent neutral probe, with factorizing \emph{dynamics}
\begin{equation}\label{eq:v11-factorHam}
 H=H_\psi\otimes I_\chi+I_\psi\otimes H_\chi,
 \quad O_{\rm src}=O_\psi\otimes I_\chi,\quad
 O_{\rm pr}=I_\psi\otimes O_\chi.
\end{equation}
This model is an explicit extension of the minimal source-only microscopic action by a \emph{genuinely decoupled} spectator. Its purpose is to test the universality claim rather than to assert that all particles in an unrestricted microscopic QFT must be decoupled.
\begin{proposition}[Cross-sector retarded response vanishes for every initial state [D]]\label{prop:v11-noCross}
For any initial density operator $\hat\varrho$ on the tensor-product Hilbert space, \emph{including an entangled state}, the exact retarded response obeys
\begin{equation}\label{eq:v11-Kubo}
\boxed{\chi^R_{{\rm pr},{\rm src}}(t)
 =-i\theta(t)\operatorname{Tr}\bigl(\hat\varrho
 [O_{\rm pr}(t),O_{\rm src}(0)]\bigr)=0.}
\end{equation}
No resummation or higher loop order internal to the $\psi$ sector can make a causally independent $\chi$ particle acquire the source's gravitational propagation. If QGD claims universal influence on $\chi$, its complete microscopic theory must generate an actual interaction, common constraint or other physically specified connection.
\end{proposition}
\begin{proof}
The time-evolution operator factorizes into $U_\psi(t)\otimes U_\chi(t)$ under Eq.~\eqref{eq:v11-factorHam}. Thus $O_{\rm pr}(t)=I_\psi\otimes O_\chi(t)$, which commutes identically with $O_{\rm src}(0)$. The operator commutator vanishes independently of $\hat\varrho$. Equivalently, with independent external source insertions and factorizing path-integral measure, $W[J_\psi,J_\chi]=W_\psi[J_\psi]+W_\chi[J_\chi]$, so $\delta^2W/\delta J_\psi\delta J_\chi=0$. The latter path-integral argument assumes a product preparation; the operator argument does not. The verifier constructs an entangled Bell-state example with connected equal-time correlation equal to one, while its retarded cross-response is zero at every tested time.
\end{proof}
\noindent\textbf{Logical boundary.} The original F1 also permits coupling to gauge fields, and an extended microscopic theory could contain other common interactions. The proposition does \emph{not} prove an all-theories no-go for emergent gravity or that a photon with electromagnetically charged matter is exactly decoupled. It does establish that \emph{the minimal isolated scalar-contact source sector alone} cannot universally influence a separate causally uncoupled probe. A physical QGD law must pass this cross-species test rather than silently assign its candidate scalar response to all probes.

\subsection{QM018B: scalar-resummed TT sector remains unexcited in the isotropic phase [D/C]}\label{sec:v11-TT}
In a parity-even isotropic state the scalar--spatial-stress vertex at a static wavevector must have the tensor structure $\Pi_{ij,S}=a(k^2)\delta_{ij}+b(k^2)k_i k_j$. A transverse-traceless polarization satisfies $e^{ij}_{\rm TT}\delta_{ij}=e^{ij}_{\rm TT}k_i=0$, hence
\begin{equation}\label{eq:v11-ttzero}
\boxed{e^{ij}_{\rm TT}\Pi_{ij,S}=0,\qquad
 e^{ij}_{\rm TT}\left(\Pi_{ij,S}D_s\Pi_{S,mn}\right)e^{mn}_{\rm TT}=0.}
\end{equation}
The zero remains exact for \emph{any} scalar-only resummation $D_s$ preserving the stated rotational structure, including an artificially tuned scalar critical point. It is not an assertion that all possible tensor interactions or genuinely dynamical spin-two composites are absent; it closes the proposed route from a \emph{single isotropic scalar HS mode} to an isolated TT gravitational pole at this level.

\subsection{QM019: an exact species-specific Clifford geometry is still not gravity [D/C]}\label{sec:v11-tomography}
Suppose a low-energy first-order truncation of one fermion's inverse operator has kinetic matrices $\Gamma_a^\mu$. Define
\begin{equation}\label{eq:v11-cliff-test}
 g_a^{\mu\nu}:=-\frac1{2d_a}\operatorname{tr}\{\Gamma_a^\mu,\Gamma_a^\nu\},\quad
 \mathcal R_a^{\mu\nu}:=\{\Gamma_a^\mu,\Gamma_a^\nu\}+2g_a^{\mu\nu}I_{d_a}.
\end{equation}
A single effective Clifford metric exists in the stated convention only when every $\mathcal R_a^{\mu\nu}=0$ with a nondegenerate Lorentzian $g_a$ (plus the necessary spin-connection, locality, physical normalization and causal conditions). Different probes must agree on a \emph{common} metric to justify QGD universality, modulo physically appropriate conformal freedom for strictly massless characteristics.

For the \emph{illustrative} screened first-order approximation $\Gamma_\psi^0=\gamma^0$ and $\Gamma_\psi^i=(1+C_s)\gamma^i$ (overall self-energy sign convention from Eq.~\eqref{eq:v11-Fock}), the Clifford test passes exactly, yielding
\begin{equation}\label{eq:v11-optical}
\boxed{g_\psi^{\mu\nu}=\operatorname{diag}(-1,(1+C_s)^2,(1+C_s)^2,(1+C_s)^2).}
\end{equation}
The decoupled spectator instead has $g_\chi^{\mu\nu}=\eta^{\mu\nu}$; for the $n=8$ benchmark $\|g_\psi-g_\chi\|_F\simeq0.01467233\ne0$. This effective \emph{species-specific medium geometry} is not a unique universal spacetime metric. Moreover, Eq.~\eqref{eq:v11-optical} is a restricted static first-order IR model, not a demonstrated full causal Dirac operator.

The complete QGD founding metric is retained without identifying this medium tensor with a derived gravitational metric:
\begin{equation}\label{eq:v11-master}
\boxed{g_{\mu\nu}=T^a{}_\mu T^b{}_\nu\left[\eta_{ab}
 -\sum_s\varepsilon_s\sigma_a^{(s)}\sigma_b^{(s)}-\lambda\mathcal Q_{ab}\right],\qquad s_s=-\varepsilon_s.}
\end{equation}
In the fixed-reference, $\lambda\mathcal Q=0$ eight-sector specialization, \emph{any} sufficiently regular metric perturbation $h=g-\eta$ can be algebraically represented by $P=bI+h/(4b)$ and $N=bI-h/(4b)$, for $b\ne0$: $P^TP-N^TN=h$. The QM019 verifier independently checks this exact signed lift of both illustrative $g_\psi$. \emph{Therefore the mere ability to assign signed natural invariants to a medium's optical metric is not proof that a unique, source-derived universal gravitational $[\Sigma]$ has emerged.} This is a constructive counterexample to a superficially persuasive but logically incomplete micro-to-$\Sigma$ claim.

\subsection{Decisive research verdict and the restricted path forward}\label{sec:v11-verdict}
\begin{center}\small
\begin{tabularx}{\linewidth}{@{}p{.38\linewidth}X@{}}\toprule
Test&Result and exact scope\\\midrule
Kinetic dressing appears when momentum-dependent screened exchange is included&\textbf{YES [D/C]} within the declared low-$k$-extended, static, frozen-HS diagnostic; not completed covariant 2PI\\
Asymptotic inverse-Laplacian pole of the same scalar mode&\textbf{NO [D/C]}: $K_0>0$, $k^2D_s(k)\to0$, Yukawa decay in its gradient truncation\\
Universal response for a completely independent spectator&\textbf{NO [D]}: exact zero retarded Kubo cross-response for arbitrary initial state\\
TT graviton from isotropic single scalar exchange&\textbf{NO [D/C]} at the stated scalar-resummed RPA level by rotational representation theory\\
Clifford metric of one dressed Dirac species&\textbf{YES [D/C]} algebraically, but it differs from the independent spectator's flat metric\\
Full microscopic source $\to$ common metric $\to[\Sigma]$&\textbf{NOT ESTABLISHED [O]}; cannot be inferred from any of the preceding five observations\\\bottomrule
\end{tabularx}
\end{center}
\textbf{Decision.} The minimal scalar-contact/HS mechanism tested through these approximations \emph{fails} the combined long-range and universality criteria. This is a precise failure of the declared candidate, not of QGD's effective gravitational correspondence or every possible underlying QFT. Increasing exchange order in a decoupled source sector cannot repair the exact cross-species causal factorization. Adding a genuine universal microscopic interaction/constraint, or deriving a nontrivial shared emergent collective mediator, is a separate physical hypothesis that requires its own action and Ward-completed proof. Claude's independent auxiliary-$\Sigma$ premise (Q) stays a clearly labeled alternative, not an adopted derivation.

\paragraph{What would count as the next real breakthrough?}
Derive a single state-dependent, conserved, causal, source-visible \emph{tensorial} collective response whose physical pole or constraint-mediated field is shared by two otherwise distinct probe species. Compute from one consistent regulated microscopic action (i) its complete physical $T^{\mu\nu}$-source vertex and full nonlocal dynamic Hessian, (ii) the Ward-completed fermion and electromagnetic probe kinetic responses, (iii) the long-distance static source residue and time-dependent homogeneous sector, and (iv) the common Clifford/null characteristics. Only then use the exact QGD signed map to construct $[\Sigma]$; do not use that algebraic map to manufacture the input response. For a massless universally coupled helicity-two composite inside an otherwise ordinary flat Lorentz-covariant stress theory, the Weinberg--Witten hypotheses and their possible violations must also be treated at the microscopic level.

\paragraph{External methodological comparators.} Ward-completed 2PI approximations: \url{https://doi.org/10.1103/PhysRevD.66.025028}. Stress-tensor measures and identities: \url{https://doi.org/10.1103/PhysRevD.23.2262}. Perturbative nonlocal Fermi-liquid kinetic self-energy is a standard physical possibility, e.g. the Galitskii low-density $m^*/m=1+[8(7\ln2-1)/(15\pi^2)](k_Fa)^2+\cdots$ coefficient (this is a different nonrelativistic, renormalized model and is a \emph{comparator}, not a computation of the relativistic QGD contact benchmark). Universal soft-graviton coupling is conditional on a real massless spin-two pole and Lorentz-invariant scattering: \url{https://doi.org/10.1103/PhysRev.135.B1049}. Weinberg--Witten theorem: \url{https://doi.org/10.1016/0370-2693(80)90212-9}.

\paragraph{Reproducibility.} Run \texttt{verify\_qm017.py}, \texttt{verify\_qm018.py}, and \texttt{verify\_qm019.py}; inspect the generated JSON, figures and the explicit approximation scopes. No existing v1.0 result has been reclassified as a full microscopic gravitational derivation.
% ===== END QM017--QM019 =====

% ===== QM022: Projective spinor optical bridge -- research edition v1.2 =====
\section{QM022: A spinor-optical construction of a Kerr--Schild weighted basis}\label{sec:v12-qm022}
\textbf{Research scope.} The objective is the founding physical map from a quantum source state to the \emph{physical equivalence class} of natural invariants, not a derivation of gravity from a separate scalar force. We develop exact algebraic/kinematic results and test a restricted constitutive candidate. The projective spinor field below is \emph{constructed from the Kerr reference congruence}; its origin as a solution or observable of the original microscopic Dirac/contact theory is \textbf{open}. Previous QM017--QM019 obstructions and all Chapters 1--5 assumptions remain unchanged.

\subsection{The complete master metric and its restricted null sector}
The complete founding equation remains
\begin{equation}\label{eq:v12-master}
\boxed{g_{\mu\nu}=T^a{}_{\mu}T^b{}_{\nu}
\big[\eta_{ab}-\sum_s\varepsilon_s\sigma_a^{(s)}\sigma_b^{(s)}-\lambda\mathcal Q_{ab}\big],\quad s_s=-\varepsilon_s.}
\end{equation}
Our exact optical construction concerns \emph{only} a fixed flat reference with $\lambda\mathcal Q=0$, one positive signed null update $g=\bar g+Hk\otimes k$, and $\sigma=\sqrt H\,k$ where $H\ge0$. A general signed eight-sector realization remains available but is not the microscopic identity of eight independent quantum spinors.

\subsection{Pure chiral spinors supply exact null directions, not unique gravity}
For a nonzero two-spinor $\chi=(u,v)^T$, let $N=\chi^\dagger\chi$ and $n^i=\chi^\dagger\tau^i\chi/N$. The Pauli matrix identity gives
\begin{equation}\label{eq:v12-Pauli}
(\chi^\dagger\boldsymbol\tau\chi)^2=(\chi^\dagger\chi)^2,
\quad |\boldsymbol n|=1,\quad
j_\chi^\mu=N(1,\boldsymbol n),\quad \eta_{\mu\nu}j_\chi^\mu j_\chi^\nu=0.
\end{equation}
For a regular four-spinor, its two chiral currents are separately null, $j^\mu\pm j_5^\mu=2j^\mu_{R/L}$, up to the convention for $\gamma^5$. This is an exact spinor-to-null-\emph{direction} bridge, not a physically selected exterior gravitational covector. In particular a general incoherent ensemble does not satisfy the pure-spinor Fierz identities for its averaged bilinears.

For a positive Hermitian chiral coherence matrix $C_{A\dot B}=\langle\chi_A\chi_{\dot B}^\dagger\rangle$, write $J^0=\operatorname{tr}C$ and $J^i=\operatorname{tr}(C\tau^i)$. Then
\begin{equation}\label{eq:v12-coherence}
\boxed{(J^0)^2-|\boldsymbol J|^2=4\det C\ge0.}
\end{equation}
A nonzero averaged chiral current is null exactly when $C$ has rank one; a generic full-rank mixture is timelike. For the unpolarized $C=I_2/2$, $J^\mu=(1,0,0,0)$ and both eigendirections are degenerate. Thus a locally isotropic many-body state does not select a unique projective null direction by taking a mean current. This is an exact obstruction to a \emph{naive} averaged-spinor rule, not to nonlocal or dynamically selected congruences.

\subsection{An explicit normalized Kerr projective spinor and its optical transport}
Use $G=c=1$, stationary flat oblate coordinates
\begin{equation}\label{eq:v12-oblate}
 x=\sqrt{r^2+a^2}\sin\theta\cos\phi,\quad
 y=\sqrt{r^2+a^2}\sin\theta\sin\phi,\quad z=r\cos\theta,\quad
 \Sigma=r^2+a^2\cos^2\theta.
\end{equation}
The following normalized spinor defines a projective Pauli direction:
\begin{equation}\label{eq:v12-spinor}
\boxed{\xi=\begin{pmatrix}\cos(\theta/2)\\ e^{i\beta}\sin(\theta/2)\end{pmatrix},\qquad
\beta=\phi-\arctan(a/r),\quad Y=\xi_2/\xi_1=e^{i\beta}\tan(\theta/2).}
\end{equation}
In terms of the flat Cartesian position,
\begin{equation}\label{eq:v12-direction}
\boxed{n_x=\frac{rx+ay}{r^2+a^2},\quad
n_y=\frac{ry-ax}{r^2+a^2},\quad n_z=\frac zr,\quad \boldsymbol n^2=1.}
\end{equation}
Direct differentiation of the oblate coordinate embedding yields the \emph{exact} Kerr--Schild null one-form
\begin{equation}\label{eq:v12-kerr-k}
\boxed{k=dt+\boldsymbol n\cdot d\boldsymbol x
=dt+\frac\Sigma{r^2+a^2}dr-a\sin^2\theta\,d\phi.}
\end{equation}
No Kerr metric or Einstein field equation enters Eqs.~\eqref{eq:v12-spinor}--\eqref{eq:v12-kerr-k}; the physical congruence, parameter $a$, reference slicing and normalization \emph{are supplied as the ansatz}. The projective spinor necessarily needs a second chart at the south polar axis. The signature convention matters: $k_\mu=(1,\boldsymbol n)$ has raised vector $k^\mu=(-1,\boldsymbol n)$ in $(-+++)$. The \emph{future-directed chiral current} aligned with the same null lines is supplied by the antipodal spinor
\begin{equation}\label{eq:v12-helicity}
\zeta=i\tau_y\xi^*=\begin{pmatrix}e^{-i\beta}\sin(\theta/2)\\-\cos(\theta/2)\end{pmatrix},\quad
\zeta^\dagger\boldsymbol\tau\zeta=-\boldsymbol n,\quad
j^\mu_\zeta=\mathcal N(1,-\boldsymbol n)=-\mathcal N k^\mu.
\end{equation}
Thus \emph{one helicity/convention generates the Kerr spatial covector while the antipodal one produces the future-directed raised current} in this signature. Conflating covector components with current-vector components would introduce a sign error; the two are related by the metric and chirality choice, not by an unqualified identity.

Define the flat-reference null-ray transport $\mathcal L=\partial_r-a(r^2+a^2)^{-1}\partial_\phi$ for time-independent fields (the full null tangent is $-\partial_t+\mathcal L$). The following identities hold \emph{exactly}:
\begin{equation}\label{eq:v12-transport}
\boxed{\mathcal L\beta=0,\quad \mathcal L\xi=0,\quad
\mathcal L\boldsymbol n=0,\quad
\mathcal L\boldsymbol x=\boldsymbol n.}
\end{equation}
Thus each flat-reference ray is straight and the projective spinor orientation is parallel-transported along that ray; the congruence as a whole may still have nonzero twist. This transport identity is \emph{not} the four-component curved Dirac equation nor the complete free Weyl equation.

\subsection{Exact optical scalars, Kerr twist and a limited weight rule}
With $\sqrt{\bar\gamma}=\Sigma\sin\theta$ the spatial flat-reference volume density and $n^r=1$, $n^\theta=0$, $n^\phi=-a/(r^2+a^2)$, direct calculation gives
\begin{equation}\label{eq:v12-expansion}
\boxed{\Theta:=\nabla_{\!\bar\gamma}\!\cdot\boldsymbol n=\frac{2r}{\Sigma}.}
\end{equation}
In a screen orthogonal to $\boldsymbol n$ the optical derivative is, up to screen orientation,
\begin{equation}\label{eq:v12-optical}
\boxed{B_{AB}=\frac1\Sigma\begin{pmatrix}r&a\cos\theta\\-a\cos\theta&r\end{pmatrix},\quad
\operatorname{shear}(B)=0.}
\end{equation}
The spacetime-null covector is \emph{not} generally a rescaled real eikonal gradient:
\begin{equation}\label{eq:v12-twist}
[k\wedge dk]_{t\theta\phi}=-2a\sin\theta\cos\theta,
\quad \boldsymbol n\cdot(\nabla\times\boldsymbol n)=-\frac{2a\cos\theta}{\Sigma}.
\end{equation}
A single smooth $S$ cannot supply $k\propto dS$ on a twisting open set. This is fully compatible with Chapter 1's independent phase $S$, spinor orientation $\mathsf R$ and current $j^\mu$: none should be equated to the gravitational natural invariant without a microscopic law.

\textbf{Restricted constitutive identity [D/C].} For the Kerr and Schwarzschild \emph{vacuum} family, the exact gravitational Kerr--Schild weight is
\begin{equation}\label{eq:v12-H-expansion}
\boxed{H_{\mathrm{Kerr}}=\frac{2Mr}{\Sigma}=M\Theta=M\,\nabla_{\!\bar\gamma}\!\cdot\boldsymbol n.}
\end{equation}
This combines one charge $M$ with a first derivative of the projective-spinor direction. It is an identity for this normalized stationary Kerr family, \emph{not} an experimentally established or microscopic universal selection law. The normalization $k_t=1$ and stationary flat reference fix its meaning. In ordinary units the coefficient is $G_NM/c^2$ and $a=J/(Mc)$.

A more useful variable is the \emph{mass aspect} $\mathcal M:=H/\Theta$ (defined where $\Theta\ne0$). For the chosen Kerr family, $\mathcal M=M$; for aligned Kerr--Newman electrovac,
\begin{equation}\label{eq:v12-KN-aspect}
\boxed{\mathcal M_{\mathrm{KN}}=M-\frac{Q^2}{2r},\qquad
H_{\mathrm{KN}}=\Theta\left(M-\frac{Q^2}{2r}\right)}
\end{equation}
with $Q$ geometrized. Write $H=2f(r)/\Sigma$, $f=r\mathcal M(r)$. The independent equatorial angular vacuum Einstein equation from QM020 reads
\begin{equation}\label{eq:v12-angular}
\boxed{R_{\theta\theta}|_{\theta=\pi/2}=\frac{2(rf'-f)}{r^2}=2\mathcal M'(r).}
\end{equation}
It selects $\mathcal M'(r)=0$ in vacuum within the ansatz. \emph{This equation follows from the adopted macroscopic Einstein dynamics, not the microscopic Dirac theory.} The identity \eqref{eq:v12-H-expansion} does not secretly derive it from fermions.

\subsection{Why the conserved spinor current does not select the Kerr weight}
A stationary conserved null current along this reference congruence can have the form $j^\mu=\mathcal N\ell^\mu$, with the future tangent $\ell=\partial_t-\mathcal L=-k^\sharp$ and
\begin{equation}\label{eq:v12-conserve}
\boxed{\nabla^{\rm flat}_\mu j^\mu=0\quad\Longrightarrow\quad
\mathcal N(r,\theta,\phi)=\frac{C(\theta,\beta)}{\Sigma}}
\end{equation}
along the rays, on an open region where the congruence is defined and with stationary boundary data. The ray-constant $C$ is not fixed by continuity. Even if a microscopic chiral spinor is aligned with this congruence, its conserved number-current density is proportional to $\Sigma^{-1}$, whereas the gravitational Kerr weight is proportional to $r\Sigma^{-1}$. Therefore current conservation plus a correct twisting direction does \emph{not} produce the physical gravitational weight. The scalar $\Theta$ in Eq.~\eqref{eq:v12-H-expansion} contains the extra radial factor; the gravitational identity is a \emph{separate constitutive assumption} unless a microscopic action supplies it.

\subsection{Current conservation is weaker than microscopic Weyl dynamics}
A further physical test must not be skipped: the exact null-ray transport $\mathcal L\xi=0$ and the conserved current $\mathcal N\propto\Sigma^{-1}$ do not imply that the spinor solves a microscopic wave equation. For the simplest real current-normalized \emph{opposite-helicity} choice $\psi_\mathrm{trial}=\Sigma^{-1/2}\zeta$ aligned with the future tangent of Eq.~\eqref{eq:v12-helicity}, we evaluated the \emph{flat, stationary} Weyl residual $\boldsymbol\tau\cdot\boldsymbol\nabla_{\!\mathrm{Cartesian}}\psi_\mathrm{trial}$, differentiating through the full oblate-to-Cartesian Jacobian. Its norm is nonzero: $0.0401770$ at $(r,a,\theta,\phi)=(2.5,0.4,0.8,0.3)$, $0.0313335$ at $(4,1.5,1.3,1.2)$, and $0.0340614$ even for the Schwarzschild-direction control $(3,0,1.1,0.5)$ in the stated length normalization. This \textbf{[D/C]} counterexample proves only that \emph{this specific spinor amplitude and phase} are not full stationary free Weyl solutions. Another complex amplitude, phase, time dependence or interaction may change the result. The weaker continuity equation cannot replace a check of the actual Dirac/Weyl operator.\par

\subsection{Failure and nonuniqueness tests; why this is not a hidden complete $\mathcal F_{\rm micro}$}
\textbf{(1) Non-universality of $H=M\nabla\cdot\boldsymbol n$ [D].} Let $\boldsymbol n=(\cos\kappa x,\sin\kappa x,0)$ in a fixed Cartesian Minkowski reference. It is a legitimate unit projective-spinor direction. The proposed weight gives $H=-M\kappa\sin(\kappa x)$. For $g=\eta+H(dt+\boldsymbol n\cdot d\boldsymbol x)^2$, the stationary \emph{linearized} vacuum equation would require $R_{tt}^{(1)}=-\tfrac12\nabla^2H=0$, but instead
\begin{equation}\label{eq:v12-counterexample}
\boxed{R_{tt}^{(1)}=-\frac12M\kappa^3\sin(\kappa x)\ne0\quad\text{generically}.}
\end{equation}
Thus arbitrary Pauli-direction fields do not solve EFE under the divergence rule. The congruence's optical/differential restrictions are essential. The same rule gives $H=0$ for a parallel null direction, whereas nontrivial pp-waves with parallel null rays exist; it cannot be a universal solution engine.

\textbf{(2) Conserved global charges alone are insufficient [D/C].} Two equal, static masses $M/2$ at positions $\pm d\hat z$ and at $\pm d\hat x$ have identical total $M$ and $J=0$, but their $zz$ trace-free quadrupole values are respectively $Q_{zz}=2Md^2$ and $Q_{zz}=-Md^2$ (when the weights sum to $M$). Their weak-field exterior metrics therefore differ. The Kerr rule $a=J/M$, $H=M\Theta$ cannot represent every physical source with the same charges; extra density-matrix multipoles and boundary structure are necessary. The special Kerr/no-hair multipole relation is a further conditional physical restriction, not a generic consequence of $M,J$.

\textbf{(3) Local exterior obstruction [D/C].} If an alleged $\mathcal F_{\rm micro}$ is a strictly pointwise finite-jet functional of local matter observables, and those observables coincide with the same vacuum values outside every compact source, then its output there is the same for every source. Such a map cannot distinguish Schwarzschild masses or Kerr angular momenta in a common exterior region. Consequently a successful map must violate at least one premise: involve nonlocal source-state information, a dynamically transmitted or constrained collective field, state-sensitive vacuum correlations, or other independent exterior/boundary data. This is an information argument, not a proof that a new fundamental force must be postulated.

\textbf{(4) Berry phase is not spacetime twist [D].} The phase connection of the chosen normalized spinor is
\begin{equation}\label{eq:v12-Berry}
\mathcal A_\xi=-i\xi^\dagger d\xi
=\sin^2(\theta/2)\left(d\phi+\frac{a}{r^2+a^2}dr\right),\quad
(d\mathcal A_\xi)_{\theta\phi}=\frac{\sin\theta}{2}.
\end{equation}
This Berry curvature remains nonzero at $a=0$ even when the radial null congruence is twist-free. It cannot be identified with the spacetime Riemann tensor or the Frobenius twist of $k$ without further structure.

\subsection{Resulting physical research criterion}
\begin{equation}\label{eq:v12-operational}
\boxed{\begin{gathered}
(S_{\rm micro},\widehat\varrho_{\rm source},\mathrm{BC})
\xrightarrow{\mathcal F_{\rm physical}}
(\boldsymbol n_{\rm projective},\mathcal M_{\rm optical})\\[3pt]
\longrightarrow[\Sigma]_{\rm physical}
\xrightarrow{\text{master metric}}g_{\mu\nu}.
\end{gathered}}
\end{equation}
On the restricted stationary Kerr branch a \emph{kinematic candidate} is to extract the source charges $M,J$, supply a projective field $Y=\tan(\theta/2)e^{i[\phi-\arctan(a/r)]}$ with $a=J/M$, set $H=M\nabla_{\rm flat}\cdot\boldsymbol n$ and $\sigma=\sqrt H\,k$. Its geometric output is exactly Kerr where the ansatz and vacuum conditions hold. What is \emph{not} derived: why the original interacting quantum state chooses this $Y$ among other projective fields, why its mass aspect obeys the vacuum source equation, how a vacuum exterior carries source information, whether the projective field solves the microscopic Dirac equation, how the general $\mathcal Q$ term behaves, and how universal matter species experience the same geometry. In particular an explicitly reconstructed Kerr spinor is not proof of microscopic gravitational emergence.

\textbf{Next concrete test.} Construct a \emph{state-derived} chiral two-point/coherence transport or stress-kinetic collective observable for a rotating source, with its causal exterior propagation/constraint equation from the same microscopic action. Test whether its projective eigen-direction is ray-transported, geodesic and shear-free, whether its mass aspect is constant in an exterior vacuum without using EFE as an input, and whether both Dirac and electromagnetic probes couple to the same physical coframe. For general sources include more than $M,J$: spatial multipoles and nonstationary source data cannot be dropped. The mathematical Kerr spinor and divergence identity are exact \emph{matching targets}; the physically selected $\mathcal F_{\rm micro}$ remains \textbf{open}.

\paragraph{Reproducibility and provenance.} Run \texttt{verify\_qm022.py} and \texttt{verify\_qm022\_optical.py} for symbolic Pauli and Gram identities, exact null-ray transport, the optical divergence and screen shear, the current-amplitude mismatch, the source quadrupole counterexample, and the failed general divergence prescription. General projective-spinor/null-congruence mathematics is known from the Kerr theorem and spinor formalism (see, e.g., \url{https://link.springer.com/article/10.12942/lrr-2012-1} and \url{https://doi.org/10.1140/epjc/s10052-014-2854-y}); the new content here is the explicit QGD-specific physical-selection audit and its restricted constitutive candidate. The prior QM020--QM021 Einstein calculations supply the independent stationary Kerr and wider-solution baselines; they are not recast as microscopic derivations.
% ===== END QM022 =====

% ===== QM023: CONSOLIDATION AND EXTERIOR COLLECTIVE FIELD RESEARCH CHARTER =====
\clearpage
\section{QM023: What the full QGD programme has accomplished and the newly fixed physical question}
\label{sec:v13-qm023}
\noindent\textbf{Scope.} This section consolidates existing research results from Chapters 1--8, QM001--QM022 and the Chapter~2 QM020--QM021 supplements. It adds a declared working hypothesis and explicit acceptance criteria, \emph{not} a new microscopic derivation or an independently verified QFT pole. The 1.2 manuscript is retained intact. Throughout, [D] denotes a proved identity or model-calculation under its declared assumptions, [C] a conditional effective-theory result, [A] a newly declared physical hypothesis, and [O] an open foundational claim.

\subsection{The decision that fixes our next physical research direction [A]}
\label{sec:v13-field-commitment}
\textbf{New user-selected physical working hypothesis (2 November 2026).} QGD should test whether the natural-invariant-producing structure is an \emph{actual collective quantum field extending into the exterior vacuum}, not merely a nonlocal label assigned to the source state or a reconstruction from an already selected solution of Einstein's equations. ``Actual'' refers to degrees of freedom or state-dependent collective observables with operationally definable correlations, source response and causal evolution. It does \emph{not} by definition imply a new fundamental elementary particle, an individual spinor expectation value, a prescribed massless pole, an a priori Einstein--Hilbert action or a nonzero Lorentz-breaking vacuum expectation value.

We introduce the \emph{placeholder}, not an invented operator identity,
\begin{equation}\label{eq:v13-physical-field}
\widehat{\mathfrak B}^{I}(x)\quad\text{or a constrained collection of composite correlators}\quad
\mathfrak b_{\rho}^{I}(x)=\big\langle\widehat{\mathfrak B}^{I}(x)\big\rangle_{\widehat\rho_{\rm source}},
\end{equation}
where $I$ is an unspecified collective/internal index and $x$ includes the exterior region. If symmetry requires $\langle\widehat{\mathfrak B}\rangle=0$, the gauge-invariant correlators or effective action, rather than this particular one-point function, must encode the field. Neither $I$ nor the microscopic definition of $\widehat{\mathfrak B}^{I}$ is claimed to be known. The field may be composite, emergent or constrained; microscopic origin must be stated in the final theory, not hidden in notation.

\begin{center}\fbox{\begin{minipage}{0.91\linewidth}
\textbf{CENTRAL QUESTION OF CHAPTER 6 --- QM023 [O]}\\[4pt]
\emph{Starting only from a specified microscopic quantum-matter action, source density operator and physical boundary/initial data, can we derive a collective quantum field with a causal state and nontrivial exterior-vacuum dynamics, whose microscopic selection equations determine the weighted directions and weights of the QGD natural-invariant equivalence class $[\Sigma]_{\rm physical}$, without inserting the target Einstein metric, Einstein equations or gravitational Green function as its premise?}\\[4pt]
\textbf{Decisive verification:} The resulting field must generate an operationally common effective geometry for distinct probes and reproduce the Einsteinian phenomenology in the appropriate macroscopic regime, including rotating twist, source multipoles, stationary/causal response and a homogeneous cosmological sector.
\end{minipage}}\end{center}

The original historical research arrow remains part of the programme:
\begin{equation}\label{eq:v13-historical-arrow}
\mathcal D_{\rm eff}\Psi=0\longrightarrow(\Psi,S,\mathsf R,j^\mu)
\overset{\mathcal F_{\rm micro}}{\longrightarrow}[\Sigma].
\end{equation}
But the adequate \emph{source-state formulation} must preserve additional information not contained in an arbitrary single probe spinor:
\begin{equation}\label{eq:v13-research-arrow}
\boxed{
(S_{\rm micro},\widehat\rho_{\rm source},\mathrm{BC})
\overset{\text{derive}}{\longrightarrow}
\mathfrak B_{\rm collective}(x)
\overset{\mathcal F_{\rm micro}}{\longrightarrow}
[\Sigma]_{\rm physical}
\overset{g[\Sigma]}{\longrightarrow}
 g_{\mu\nu}^{\rm physical}.}
\end{equation}
The derived microscopic field and the source-dependent map are open. The final metric arrow is the known QGD algebraic architecture. A second valid \emph{operational} ordering, if the quantum collective field first determines common kinetic cones, is $\mathfrak B\to g_{\rm micro}\to[\Sigma]$; these are not competing ontologies when they identify the same equivalence class in a regular metric-only branch.

\subsection{Why the weighted covectors work mathematically: exact statement, limited physical claim [D/C]}
\label{sec:v13-master-metric}
Do not replace the founding formula with its classical specialization. The complete master metric is
\begin{equation}\label{eq:v13-master}
\boxed{g_{\mu\nu}=T^a{}_{\mu}T^b{}_{\nu}
\left[\eta_{ab}-\sum_s\varepsilon_s\sigma_a^{(s)}\sigma_b^{(s)}
-\lambda\mathcal Q_{ab}\right],\qquad s_s=-\varepsilon_s.}
\end{equation}
For the \emph{declared} fixed-coframe, $\lambda\mathcal Q=0$ signed-sector specialization, the metric deformation is an indefinite Gram form, $h=\Sigma^\mathsf TJ\Sigma$. Thus the weighted covectors collectively determine the ten metric components. On a regular spanning eight-sector branch, $\mathrm{rank}(D_\Sigma g)=10$ and $\dim\ker(D_\Sigma g)=22$; the null directions coincide with the $O(4,4)$ orbit tangents when the action depends on the sectors only through this metric. This establishes a \emph{metric representation and quotient}, not the microscopic quantum origin of any particular sector representative. The counting changes on singular strata and generally when the historical $\mathcal Q_{ab}$ correction is active.

The covectors are \emph{weighted basis-like directions}, not vectors each carrying its own scalar Riemann curvature. The Levi--Civita connection and Riemann tensor are constructed from the \emph{complete} effective metric and its derivatives. The one-null-sector example $g=\bar g+Hk\otimes k$, $\sigma=\sqrt{H}\,k$ and $k^2_{\bar g}=0$ admits the exact coframe $e^a=\bar e^a+(H/2)k^ak$ and unchanged volume determinant. Kerr may have unchanged flat-reference volume determinant but nonzero Weyl curvature; Milne has varying coframe weights but vanishing four-dimensional curvature. Neither the absolute weight nor $\nabla H$ alone is a universal definition of gravity.

\subsection{Consolidated achievement ledger: positive theorems, constraints and genuine failures}
\label{sec:v13-ledger}
\textbf{Interpretation rule.} A verified negative result that excludes a tempting microscopic route is a scientific advance, but it is \emph{not} evidence that an untested alternative must succeed. We distinguish exact algebra [D], conditional Einsteinian/effective matches [C], sourced microscopic model tests [D/C], and a genuine missing selection [O].

{\small\setlength\tabcolsep{4pt}
\begin{longtable}{@{}p{.07\linewidth}p{.29\linewidth}p{.09\linewidth}p{.46\linewidth}@{}}
\toprule Chapter / test & Result & Status & Contribution to $\mathcal F_{\rm micro}$ and precise limit\\\midrule
\endfirsthead
\toprule Chapter / test & Result & Status & Contribution / limit\\\midrule
\endhead
Ch.~1 & Exact local Dirac two-block ratio $u_B=\mathsf Ru_A$ & D/C & Physical spinor orientation/momentum relation; no independent gravitational source-amplitude selection.\\
Ch.~1 & WKB, phase transport and conserved $r^2j^r$ & D/C & Quantum current/phase constraints; one phase cannot encode a twisting Kerr congruence.\\
Ch.~1 & Large-$N$ and collective semiclassical limits & D/C & Controlled macroscopic limits under specified scaling; generic many-fermion states retain correlations.\\
Ch.~1 & Four-mode frame reparametrization & D & Same spinor may admit different mode frames; a freely chosen spinor decomposition cannot fix $[\Sigma]$.\\
Ch.~1 & Mixed-state matter kinetic tensor can be full rank & D/C & Full source stress and correlations may encode more than mean current; rank is not a physical metric-selection law.\\
Ch.~2 & Independently varied spherical vacuum branch & C & $A^2=B^2=2G_NM/(c^2r)$ follows from adopted Einstein--Hilbert dynamics in the specified sector, not from flat quantum matter alone.\\
Ch.~2 & Fermion-star source and rotating geometry & C & Constructive matter-to-metric-to-$[\Sigma]$ pipeline \emph{using EFE}; source information beyond a single current matters.\\
QM020--21 & Kerr radial amplitude and broad Einstein solution checks & D/C & Rotating ansatz plus EFE selects $H=2Mr/\Sigma$; 41 checks demonstrate geometry-engine reach, not micro origin.\\
QM021 & Kerr $k\wedge dk\ne0$ & D & One scalar eikonal cannot generate the rotating null direction even after nonzero rescaling.\\
Ch.~3 & Independent $\alpha_G$ and state normalizations & D/C & Wavefunction normalization does not determine $G_N$ or the source's natural-invariant strength.\\
Ch.~3 & Spin-singlet source with $j^r=0$ and nonzero mass & C & Radial probability flux cannot universally serve as gravitational source variable.\\
Ch.~4 & Full metric pullback and Option C Euler equation & D/C & Mathematically exact after adopted metric-mediated action; stationarity alone not micro genesis.\\
Ch.~4 & Microscopic determinant/stress-correlator programme & D/C & Defined quantum observables and induced local terms; does not make the universal metric an output of the initial flat-matter model.\\
Ch.~5 & Stable collective static functional & C & If source-coupled $q$ exists, $G_{\rm eff}=g_c^2/(4\pi Z)$ and $H=2G_{\rm eff}M/(c^2r)$; existence of $q$ remains assumed.\\
Ch.~5 & Exact regular signed eight-sector lift & D & Any smooth metric deformation can be lifted locally on spanning branch; this is factorization, not selection.\\
Ch.~5 & Kinetic principal-symbol metric candidate & D/C & Common probe cones could operationally identify $g_{\rm micro}$; species dependence must be excluded.\\
QM006--07 & TT continuum, stress matrix and Ward contact & D/C & Actual microscopic tensor spectral information; no isolated universal spin-two pole in tested approximation.\\
QM009 & Inhomogeneous scalar source and finite response & D/C & Source-dependent Dirac mass/ratio changes; local Hartree--Fock kinetic principal symbol remains unchanged.\\
QM010--12 & Stationary $H=\sigma_{\rm amp}^2$ variable & D/C & Exact $1/r$ exterior from stipulated $H$ energy; naive canonical $|\nabla\sigma|^2$ gives wrong exponent.\\
QM011 & Actual Hubbard--Stratonovich scalar saddle & D/C & Microscopic stationary field exists; measured positive infrared curvature yields short-range response.\\
QM013 & Induced curvature term under premise (Q) & C/A & Correctly scoped one-loop induction; independent $\Sigma$ and universal metric coupling are supplied, not derived.\\
QM014 & $O(4,4)$ rank/gauge and two-mode induced checks & D/C & Physical metric quotient validated on regular branch; full nonlinear microscopic constraint algebra remains unproved.\\
QM015 & Spectral source-overlap theorem & D & Masslessness is insufficient unless the source excites the mode and the proposed geometric observable detects it.\\
QM015 & Exactly soluble source-mixed long-range toy model & D/C & Indirect source coupling can yield $1/r$; model's shift symmetry and coefficients are assumptions.\\
QM016 & Positive Hartree scalar curvature theorem & D/C & For positive-mass solutions with $\Lambda/\mu\ge e$, the tested sharp-cutoff scalar cannot be tuned massless by $G_4>0$.\\
QM016 & Coupled scalar--density Hessian and sum-rule checks & D/C & Quantitative thermodynamic susceptibility, not a new propagating gravitational degree of freedom.\\
QM017 & Screened-exchange kinetic correction & D/C & Momentum dependence can alter a quasiparticle's principal symbol; measured screening prevents exterior Newtonian tail in this model.\\
QM018 & Factorized spectator zero retarded response & D & A truly uncoupled species has no causal source response, even in initially entangled states.\\
QM019 & Species-specific Clifford/signed lift & D/C & Algebraic metric reconstruction may succeed for one species while universal geometry fails.\\
QM022 & Kerr projective-spinor null direction & D & Exact rotating weighted-direction reconstruction and shear-free ray transport, not independent spinor dynamics.\\
QM022 & Optical weight identity and its counterexample & D/C & $H=M\nabla\cdot\bm n$ works in Kerr--Schwarzschild family, fails generally; current scaling does not fix $H$.\\
Ch.~7 & Sector-orbit flat-vacuum cancellation and energy audit & D/C & Nonzero sector amplitudes can encode flat metric; pure-orbit gradient energy is not a geometric observable.\\
Ch.~8 & FLRW lifts and no-extra-metric-equation result & D/C & Representation does not select cosmic time evolution; static inverse Laplacian lacks homogeneous zero-mode dynamics.\\
\bottomrule
\end{longtable}}

\noindent\textbf{The net achievement.} QGD has (a) a valid \emph{representation} of geometric gravitational fields, (b) a well-defined \emph{conditional effective action} recovering Einsteinian dynamics, (c) exact source-conditioned \emph{static matching} if an appropriate collective mode is stipulated, (d) exact Dirac/optical \emph{kinematic bridges}, and (e) several \emph{genuine microscopic response calculations} which rule out specific scalar-only, screened, or species-specific shortcuts. No one item proves (f): a quantum-source-derived, universally coupled, exterior-vacuum collective weighted frame. The unexplained junction is chiefly (e)$\to$(f), not (f)$\to$Schwarzschild once EFE is adopted.

\subsection{What ``an actual collective quantum field in exterior vacuum'' obliges us to explain}
\label{sec:v13-vacuum-obligation}
Three separate questions must be answered without circularity.
\begin{enumerate}[label=\textbf{V\arabic*.},leftmargin=18mm]
\item \textbf{Existence and definition.} Which local or suitably quasilocal composite of the \emph{specified} microscopic operators becomes the collective field $\widehat{\mathfrak B}(x)$? What is its symmetry/gauge class, coarse-graining scale, norm and commutator or constrained phase space? Does it exist outside a region containing material quanta, rather than only in a chosen description of a matter wavefunction? A classical nonzero mean is not required if physical correlators or constraints carry the gravitational information.
\item \textbf{Source imprint and exterior dynamics.} What causal action or equations follow from the same microscopic interactions, source state and measure? The exterior may have no local material current, $J_{\rm matter}(x)=0$, yet a sourced field can be nonzero there if the theory supplies the necessary propagating/constraint/nonlocal dynamics and matching data. \emph{An algebraic pointwise rule} $\mathfrak b(x)=f(\langle\widehat O_{\rm matter}(x)\rangle)$ with vanishing exterior operator data cannot by itself generate a nonzero exterior field; vacuum polarization, nonlocal operators, boundary propagation or additional collective degrees must be explicitly justified.
\item \textbf{Weighted-frame selection and universality.} Which physical, invariant observables of $\mathfrak B$ determine covector directions and weights, including rotating twist and multipoles? Why do electrons, neutral matter and photons register one common metric instead of unrelated optical medium tensors? How do gauge/sector redundancies and an active historical $\mathcal Q_{ab}$ affect the observable class?
\end{enumerate}
The choice of a spatially extended quantum field therefore sharpens, rather than automatically solves, the microscopic task. It is not enough to say that the original four-mode spinor ``fills the vacuum''; one must derive a collective exterior quantum state and its field equations.

\subsection{A non-circular functional architecture: a target for future derivation, not its result}
\label{sec:v13-functional-architecture}
One possible \emph{working template} is to construct from $S_{\rm micro}$ and an initial/source density operator a causal collective effective functional (for real-time backgrounds ordinarily on a closed-time path)
\begin{equation}\label{eq:v13-gamma-target}
\Gamma_{\rm coll}[\mathfrak b^+,\mathfrak b^-;\widehat\rho_{\rm source},\mathrm{BC}],
\qquad
\left.\frac{\delta\Gamma_{\rm coll}}{\delta\mathfrak b^+_I(x)}\right|_{\mathfrak b^+=\mathfrak b^-}=0,
\end{equation}
with its source coupling and Ward identities derived from the \emph{same microscopic model}. The location $x$ must range through the exterior vacuum. These expressions are \emph{desired definitions of a calculational problem}, not an asserted microscopic functional already obtained.

The missing constitutive map then has the generic form
\begin{equation}\label{eq:v13-constitutive-target}
\big[\Sigma(x)\big]_{\rm physical}
=\mathcal C\big[\mathfrak b(x),\nabla\mathfrak b(x),
\langle\delta\widehat{\mathfrak B}(x)\delta\widehat{\mathfrak B}(y)\rangle,\ldots;\mathrm{BC}\big]
\end{equation}
with $\mathcal C$ to be independently derived and invariant under unphysical operator and sector coordinate changes. We do \emph{not} preassign $\mathcal C=\sqrt H\,k$ as a universal law: that is one useful classical null-sector representative, and Kerr's $k$ is twisting. The mapping may instead first yield common probe principal symbols, from which a regular signed representative can be constructed algebraically. In either case the full metric remains Eq.~\eqref{eq:v13-master}.

\textbf{Key logical firewall.} Defining the microscopic action \emph{at the outset} as $S_{\rm micro}[\psi;g(\Sigma)]$ and subsequently integrating out $\psi$ is the separate proposed premise (Q). Its induced Einsteinian action does not establish that $\Sigma$ is a collective observable of the \emph{original} flat-matter action. That route can be studied honestly as an extension, but not counted as completion of Eq.~\eqref{eq:v13-research-arrow}.

\subsection{Why the existing Kerr/Schwarzschild bridge is powerful but not yet an origin theorem}
\label{sec:v13-optical-lesson}
Chapter~2 and QM020--QM022 establish the following exact restricted chain:
\begin{equation}\label{eq:v13-kerr-chain}
\xi_{\rm Kerr}\to k_{\rm twisting}\to
\big(H=2Mr/\Sigma,\ \sigma=\sqrt H\,k\big)
\to g_{\rm Kerr}.
\end{equation}
Here $\xi_{\rm Kerr}$ is an explicitly \emph{constructed} normalized projective spinor, the congruence is null/geodesic/shear-free, its twist obeys $k\wedge\dd k\ne0$ when rotating, and $H=M\nabla_{\rm flat}\cdot\bm n$ holds \emph{on this Kerr/Schwarzschild family}. Chapter~2 independently selects the required radial amplitude through the \emph{adopted vacuum Einstein equations} when the rotating congruence is already specified. QM022 proves that $\xi$ and conserved current do not generally determine the Kerr weight: current density scales as $1/\Sigma$, not $r/\Sigma$. Nor is the simple geometric spinor automatically a free or interacting Dirac solution. A generic incoherent source has no unique pure null spinor current; $M$ and $J$ also do not determine generic source quadrupoles.

\textbf{Therefore:} a viable actual vacuum collective field must explain \emph{both} the source-state selection of Kerr-like optical directions and the source-state selection of its weight/multipoles. These cannot simply be asserted from a spinor reconstructed after the classical geometry has been chosen. The Kerr optical example is a precise \emph{target and kinematic possibility theorem}, not $\mathcal F_{\rm micro}$.

\subsection{Why a fundamental Newtonian force is not part of the hypothesis [A/C]}
\label{sec:v13-force-distinction}
The hypothesis is a sourced \emph{weighted-frame quantum field}. Classical gravitational accelerations, lensing and curvature arise from the effective metric, its compatible connection and matter propagation. A $1/r$ exterior and the corresponding weak-field $1/k^2$ response are useful \emph{phenomenological validation tests} in the ordinary Schwarzschild/Newtonian regime; they are not a separate primitive scalar-force ontology nor a universal requirement that the microscopic field be a simple Gaussian massless scalar. The analysis of QM011--QM019 remains relevant because it excludes concrete proposed \emph{source-selection mechanisms}, not because it proves every emergent geometric mechanism must resemble a Yukawa or Newtonian potential.

The strongest falsification is not simply ``does a scalar pole appear?'', but \emph{``can one and the same microscopic collective field state yield the geometry seen by independent species, including twisting vacuum configurations, consistent nonlinear and causal evolution, and the tested Einsteinian limit?''} Any explanation that only dresses one fermion's dispersion, changes its mass or invents a metric after the fact remains incomplete.

\subsection{Five non-negotiable tests and the status after QM022}
\label{sec:v13-acceptance}
{\small
\begin{tabularx}{\linewidth}{@{}p{.055\linewidth}X p{.24\linewidth}@{}}\toprule
& Acceptance condition for the actual exterior quantum field & Present status\\\midrule
\textbf{1}& Derive an operator/correlator and source-conditioned state from the original microscopic quantum dynamics; define physical equivalences and a causal exterior equation. & \textbf{Open}; HS scalar exists but is not the required long-range frame.\\
\textbf{2}& Derive physical weighted directions (including $k\wedge\dd k\ne0$ when rotating) and weights, with independent source multipoles and boundary data; avoid a one-phase or one-current identification. & \textbf{Open}; Kerr projective-spinor reconstruction only.\\
\textbf{3}& Obtain the complete physical metric through the full master map or an independently generated common kinetic metric, and show that the selected class is not a sector-gauge artifact. & \textbf{Map proved conditionally}; physical selection open.\\
\textbf{4}& Demonstrate causal cross-species universality and consistency with matter stress/Ward identities; an uncoupled spectator must not be declared to gravitate. & \textbf{Open}; QM018 supplies exact counterexample.\\
\textbf{5}& Recover static, rotating, radiative and homogeneous/cosmological Einsteinian behavior in the valid effective regime, without assuming the requested solution as a microscopic input. & \textbf{Conditional classical matches}; no micro dynamics yet.\\\bottomrule
\end{tabularx}}

\textbf{No double counting of results.} EFE-generated solutions from Chapter~2/QM021, $H$-energy minimization from Chapter~5, induced Einstein curvature from an externally supplied $g(\Sigma)$, and Clifford factorization of a one-species kinetic tensor are independent conditional achievements. None should be counted as evidence for a physical exterior collective field until the microscopic arrow is derived.

\subsection{A compact dependency map across all eight chapters}
\label{sec:v13-chapters}
\begin{center}\small
\begin{tabularx}{\linewidth}{@{}p{.13\linewidth}X X@{}}\toprule
Chapter & What is already supplied & Missing physical link\\\midrule
1 & Dirac spinors, many-body observables, current and semiclassical limits & Genuine field operator/coherence in sourced exterior vacuum\\
2 & Natural-invariant metric and independent classical solution machinery & Quantum selection of weights, twist and multipoles\\
3 & Empirical gravitational normalization and source/probe distinctions & Microscopic relation among wave amplitude, mass and frame weight\\
4 & Full action-pullback/Option C and quantum correlator formalism & Derivation of gravity-relevant effective functional from flat microscopic matter\\
5 & Exact conditional static collective-energy bridge & Quantum derivation of the collective degree and geometric coupling\\
6 & Quantitative microscopic tests, proof ledgers, Kerr spinor-optical bridge & Actual causal, common, exterior vacuum quantum field $\mathfrak B$ and its $[\Sigma]$ map\\
7 & Vacuum-orbit redundancies and geometric energy constraints & Gauge-invariant physical vacuum field/state and gravitational energy account\\
8 & Cosmological representations and selection/no-extra-equation limits & Microscopic homogeneous initial-state and expansion dynamics\\\bottomrule
\end{tabularx}\end{center}

\subsection{Research decision: what to calculate next, and what not to repeat}
\label{sec:v13-research-decision}
\textbf{Immediate decision.} Pause additional classical Einstein-solution enumeration and do not launch another scalar-only infrared scan merely to obtain a $1/r$ tail. First propose \emph{one actual} candidate composite quantum field or coherence/connection order parameter, identify its microscopic operator definition (or explicit new physical premise if none exists), then derive the source-conditioned exterior quantum functional and its correlation/causal transport. The candidate must have enough directional structure to encode Kerr twist and enough state information to select source-dependent weights and multipoles. Its coupling to distinct probe species must come from one microscopic theory; the Chapter~1 Dirac/WKB/current identities are constraints to preserve, not substitutes for the missing dynamics.

\textbf{Falsifiable branches.} (i) If a micro-derived exterior field exists and supplies common kinetic geometry, test its QGD quotient, EFE limit, stability and quantum corrections. (ii) If the original contact model has no collective exterior mode and no universal probe interaction, report a model-specific failure and identify explicitly the additional microscopic interaction or degrees of freedom needed. (iii) If a field alters only one species, classify it as a quantum medium/optical metric, not QGD's universal gravity. (iv) If the only way to obtain the field is to postulate independent $\Sigma$ and universal $g(\Sigma)$ coupling, declare a new axiom and test that distinct theory without retroactively claiming the original $\mathcal F_{\rm micro}$ has been derived.

\textbf{Central question retained for the next chapter increment.}
\begin{equation}\label{eq:v13-question-final}
\boxed{\begin{aligned}
(S_{\rm micro},\widehat\rho_{\rm source},\mathrm{BC})
&\xrightarrow{\rm derive}\mathfrak B_{\rho}(x),\qquad x\text{ in exterior vacuum},\\
\mathfrak B_{\rho}&\xrightarrow{\mathcal F_{\rm micro}}[\Sigma]_{\rm physical}
\xrightarrow{g[\Sigma]}g_{\mu\nu}^{\rm common}.
\end{aligned}}
\end{equation}
Answering this, not finding another coordinate form of Kerr, is the intended next foundational breakthrough. The present QM023 section is a synthesis and a research-ontology decision [A/O], not a new solution of the open map. No prior proof-status claim is upgraded by adding it.
% ===== END QM023 =====

% ==== QM024 research audit: 2 November 2026. Prior chapters/labels remain unchanged. ====
\clearpage
\section{QM024: Exact microscopic rewriting versus a physical pregeometric phase}
\label{sec:v14-qm024}
\textbf{Research status [D/N/C/O].} The original F1 and the new determinant-of-spinor-bilinears pregeometry M have been tested as \emph{different} microscopic theories. This section reports mathematical identities and restricted counterexamples. It does not supersede the QM023 exterior collective-field research question or promote the new theory M to an established closure of $\mathcal F_{\rm micro}$. The complete proof, numerical procedures, and references appear in the accompanying QM024 standalone research report.

\subsection{Route A: exact rewrite of original F1 [D/N]}
\label{sec:v14-routeA}
For two fermion pairs with $B=\bar\psi_1\psi_1+\bar\psi_2\psi_2$ and an exactly defined zero-dimensional Grassmann integration, a genuine collective identity is
\begin{equation}\label{eq:v14-HS}
\boxed{Z_F(J)=J^2+g=\frac1{\sqrt{2\pi g}}\int d\phi\,e^{-\phi^2/(2g)}(J+\phi)^2,\qquad g>0.}
\end{equation}
The identity preserves all source derivatives exactly. Yet at $J=0$ the bosonic auxiliary effective action has nonzero stationary points $\phi_\pm=\pm\sqrt{2g}$ while the exact expectation value is $\langle\phi\rangle=0$. This is an exact finite-regulator \emph{counterexample to inferring a spontaneously selected physical order parameter from a formal auxiliary saddle}. Thermodynamic, gauge and continuum limits require independent treatment.

An exact source-preserving rewrite cannot change the original retarded connected observables. In particular the QM018 decoupled-source/spectator result $\chi^R_{PS}=0$ persists under exact field redefinitions, including with entangled initial states. It does not follow that \emph{every} interacting F1 model lacks a collective gravity mode; the earlier scalar/stress tests provide restricted negative results, not a global nonperturbative F1 no-go theorem.

\subsection{Route B: the metric-free microscopic action is genuinely different [D]}
\label{sec:v14-routeB}
The proposed spinor-pregeometry action uses
\begin{equation}\label{eq:v14-composite}
\boxed{E^a{}_{\mu}=\sum_{i=1}^N\frac i2\bar\psi_i\gamma^a\overleftrightarrow{D_\mu}\psi_i,\qquad S_M=\kappa\int d^4x\det E.}
\end{equation}
The determinant of the \emph{sum} $E=\sum_i E_i$ contains mixed eight-fermion interactions among species, and is not the sum of individual free Dirac actions. M can therefore supply new common microscopic structure, which is physically promising, but it is not exactly equivalent to original F1. The formal identity $1=\int DE'\,\delta(E'-E)$ and Fourier multiplier $\Pi$ is a rewrite of \emph{M}, with source preserving properties only after the contours, measures and gauge treatment have been specified.

As a finite-regulator test, an eight-generator Grassmann model with $E_{a\mu}=\bar\psi_a\psi_\mu$ obeys
\begin{equation}\label{eq:v14-detGrass}
\det E=24\,\bar\psi_1\bar\psi_2\bar\psi_3\bar\psi_4\psi_1\psi_2\psi_3\psi_4,\qquad
Z_M(j)=j^4+24\kappa.
\end{equation}
Thus the composite determinant interaction is nonzero but $\langle E_{a\mu}\rangle=0$ at zero source in this finite model. It proves neither presence nor absence of a large-$N$ geometric phase. A nonzero ordinary-number pointwise unimodular equality $\det E=v\ne0$ is incompatible with finite-regulator nilpotency of a microscopic Grassmann composite; imposing unimodularity on a genuinely bosonic auxiliary field is a different proposal requiring derivation.

\subsection{Correct frame duality and quantum measure [D/N]}
\label{sec:v14-measure}
On the regular real $4\times4$ branch,
\begin{equation}\label{eq:v14-duality}
\boxed{\Pi=(\det e)e^{-1},\quad e=(\det\Pi)^{1/3}\Pi^{-1},\quad
\left|\det\frac{\partial\Pi}{\partial e}\right|=3|\det e|^8.}
\end{equation}
The last relation follows by splitting the $16$ fluctuations into one trace direction (eigenvalue $3$ at $e=I$) and $15$ traceless directions (eigenvalue $-1$), and extending by the cofactor map. Forty random-frame tests reproduce it with relative error below $10^{-12}$. The measure factor is required for a flat finite-dimensional component measure; it is not the full gauge-fixed continuum functional measure.

The dual field $e(\Pi)$ controls the common \emph{kinetic} fermion frame. The proposed microscopic mass term, however, is $m_i\det E\,\bar\psi_i\psi_i$. Its exact auxiliary representation contains $m_i\det E'\,\bar\psi_i\psi_i$. Off shell, the effective mass relative to the kinetic frame is $m_i\det E'/\det e$, and the two determinants need not agree. For $e=\operatorname{diag}(2,1,1,1)$ and $E'=I$ the ratio is $1/2$. A common conventional massive Dirac action can follow at an appropriate classical saddle, but not as the exact off-shell operator for all $E'$ and $\Pi$ configurations.

Further, the proposed real Euclidean auxiliary potential $\kappa\det E'$ is unbounded below. Along $E'=t\operatorname{diag}(-1,1,1,1)$ one has $\det E'=-t^4$, and $\exp(-\kappa\det E')=\exp(\kappa t^4)$ for $\kappa>0$. A different contour, a regulated Grassmann construction, or a stabilized collective functional must be specified before interpreting a standalone bosonic $E'$ path integral as exact. This issue does not invalidate the regulated \emph{fermionic} determinant-composite action by itself.

\subsection{The induced Einsteinian coefficient remains a conditional success [C/N]}
\label{sec:v14-Gind}
If a nondegenerate collective frame exists and the fermions have the standard covariant kinetic action there, the proper-time Dirac determinant gives the previously derived scheme-dependent induced coefficient. For $N$ identical four-component Dirac species,
\begin{equation}\label{eq:v14-Gind}
G_{\rm ind}=\frac{6\pi}{NI_1},\qquad
I_1=\Lambda_c^2e^{-m^2/\Lambda_c^2}-m^2E_1(m^2/\Lambda_c^2).
\end{equation}
At $N=100,m=0.01,\Lambda_c=1$, the numerical result is $G_{\rm ind}=0.18867731537$. Changing $\Lambda_c$ to $0.5$ or $2$ changes it to $0.75647771473$ or $0.04713687534$, respectively, at fixed matter content. This is a valid calculation \emph{in its specified proper-time convention}, not a parameter-free physical prediction or proof of existence of a healthy collective frame phase. For $E=O(N)$, extensivity of $\kappa\det E$ at $O(N)$ suggests $\kappa=O(N^{-3})$ in the simple normalization; a controlled large-$N$ derivation must specify the scaling.

Rerunning QM014 confirms the conditional induced-Einstein metric-operator ranks $(6,4)$ at generic and null momentum, and sector ranks $(6,4)$ with nullities $(26,28)$ on the regular metric-only branch. These are \emph{not} a complete microscopic Dirac--Bergmann constraint classification. The microscopic M action also contains no independently specified Maxwell kinetic term; equality of photon and fermion effective light cones remains an independent, uncompleted microscopic test.

\subsection{Outcome for the founding $\mathcal F_{\rm micro}$ programme}
\label{sec:v14-conclusion}
\textbf{Route A:} exact finite-regulator collective rewritings preserve physical source correlators, but do not generate new universal geometry absent from the original theory. \textbf{Route B:} an exact formal auxiliary rewriting of the \emph{new} spinor-pregeometry M and a correct local frame duality are demonstrated; the nondegenerate physical collective phase, its correct massive/gauge couplings, and its exterior causal field dynamics are not yet established.

The exact CTP object needed for a source-conditioned claim is a properly regulated collective functional\[\Gamma_{\rm coll}[e_{\pm},\omega_{\pm};\widehat\varrho_{\rm source}],\] with both contour branches, correct Jacobians and physical constraints. The geometric phase must be tested using gauge-invariant or relational observables and collective spectra, not inferred from a gauge-dependent frame expectation alone. Only after such an object is computed could one legitimately select the physical QGD class using the \emph{full} master metric Eq.~\eqref{eq:full-preserved} rather than a pure-algebraic lift.

\textbf{Central open question [O].} Does a metric-free, fully regulated microscopic quantum theory dynamically produce a stable exterior collective quantum field with an invertible physical weighted frame, causal source dependence and a common relativistic propagation structure for all relevant probes? If so, does its physical class $[\Sigma]$ recover the tested Einsteinian limit without postulating that limit in the microscopic action? QM024 has established the correct formal and algebraic prerequisites, \emph{not} that the required geometric phase exists.
% ==== END QM024 ====

% ==== QM025: Regulated collective-frame phase and exact state-selection tests ====
\clearpage
\section{QM025: Exact composite pushforward, dual-frame measure and exterior phase-selection test}
\label{sec:v15-qm025}
\textbf{Scope [D/N/C/O].} We now test whether the \emph{physical phase} can be inferred from the exact collective transformation of spinor pregeometry M. Unlike QM024, which checked the dual coframe algebra, QM025 examines the support of the exact finite-regulator composite distribution, its response to external sources, coordinate-dependent auxiliary saddles, the stationary equations of both auxiliary fields, and the limitations of a homogeneous covariant frame potential. The calculations are mathematical tests of the stated proposals, not a nonperturbative simulation of the four-dimensional theory. The canonical QM023 central question remains unchanged.

\subsection{A rigorous distinction between the composite $X$ and the dual coframe $e(\Pi)$ [D]}
\label{sec:v15-twofields}
Write $E^a{}_{\mu}[\psi,\omega]$ for the microscopic Grassmann-even spinor bilinear in Eq.~\eqref{eq:v14-composite}, $X^a{}_{\mu}$ for the constrained auxiliary image of this composite, and $\Pi^\mu{}_a$ for its Fourier/Lagrange multiplier. For a finite Grassmann regulator the formal exact identity is a distributional pushforward,
\begin{equation}\label{eq:v15-pushforward}
\begin{aligned}
1&=\int d^{16}X\,\delta^{(16)}\!\bigl(X-E[\psi,\omega]\bigr),\\
Z_M[J]&=\int D\psi D\omega\,e^{iS_M+iJ\cdot O}\\
&=\int d^{16}X\,D\psi D\omega\,
\delta^{(16)}(X-E)\,e^{iS_M(X)+iJ\cdot O}.
\end{aligned}
\end{equation}
The delta symbol means a \emph{finite derivative-of-delta distribution} when its shift has nilpotent entries, not a positive stochastic delta distribution over classical frames. It reproduces polynomial test functions and source derivatives exactly with consistent integration conventions.

\begin{proposition}[Support of a finite composite pushforward, not a no-go for condensation]\label{prop:v15-support}
At any finite lattice regulator the ordinary-number body of every matrix component of $E$ vanishes. Consequently $\delta^{(16)}(X-E)$ is a finite sum of partial derivatives of $\delta^{(16)}(X)$ and has distributional support at the degenerate point $X=0$. Replacing the \emph{exact} $X$-integration by a restriction to $\det X\ne0$ removes that support and is not an exact collective rewriting. This does not exclude nonzero renormalized composite expectations or an invertible \emph{effective} frame after a thermodynamic, large-$N$, continuum, or gauge-conditioned limit.
\end{proposition}
\begin{proof} Each $E^a{}_{\mu}$ is Grassmann even and nilpotent because it has no constant component in the finite Grassmann algebra. Expand the shifted distribution as
\[
\delta^{(16)}(X-E)=\sum_{|\alpha|\le K}\frac{(-E)^\alpha}{\alpha!}\,\partial_X^\alpha\delta^{(16)}(X),
\]
with the sum terminating by nilpotency. Every term has support at $X=0$. Its pairing with a smooth function supported away from $X=0$, including one supported inside $\mathrm{GL}(4,\mathbb R)$, is zero. Hence imposing $\det X\ne0$ at fixed regulator cannot be equated to the original exact identity. A distributional support statement is not a probability claim and does not commute automatically with a singular infinite-volume limit.\end{proof}

\textbf{Exactly solvable source diagnostic.} Let $\vartheta=\bar\psi\psi$ be one Grassmann-even nilpotent variable ($\vartheta^2=0$), normalized so $\int d\bar\psi d\psi\,\vartheta=1$. Then for $j\ne0$,
\begin{equation}\label{eq:v15-onepair}
Z(j)=\int d\bar\psi d\psi\,e^{j\vartheta}=j,\qquad
\frac1{Z(j)}\int d\bar\psi d\psi\,e^{j\vartheta}\delta(x-\vartheta)
=\delta(x)-\frac1j\delta'(x).
\end{equation}
The source-generated expectation is $\langle\vartheta\rangle=1/j$ even though the pushforward's distributional support is $x=0$, because its derivative-of-delta term gives $\langle x\rangle=1/j$ while $\langle x^2\rangle=0$. The connected $x$ second moment is $-1/j^2$. The pushforward is therefore \emph{not} a positive probability distribution; a seemingly nonzero expectation is not an ordinary-valued random invertible matrix at finite regulator.

For the previous exact four-pair determinant test (QM024), $Z(j)=j^4+24\kappa$ and the average diagonal bilinear obeys
\begin{equation}\label{eq:v15-fourpair}
\frac14\partial_j\log Z(j)=\frac{j^3}{j^4+24\kappa},\qquad
\left.\langle E_{aa}\rangle\right|_{j=0}=0,\qquad
\left.\frac{d\langle E_{aa}\rangle}{dj}\right|_{j=0}=0.
\end{equation}
This is a finite-volume source-response obstruction in the \emph{specific zero-dimensional model}; it does not extrapolate to the full geometric phase.

\subsection{The dual coframe $e(\Pi)$ can be invertible without $X$ being invertible [D/N]}
\label{sec:v15-dual}
The local duality reviewed in QM024 uses the \emph{multiplier}, not the composite:
\begin{equation}\label{eq:v15-frame-dual}
\Pi=\operatorname{cof}(e),\qquad e=(\det\Pi)^{1/3}(\Pi^{-1})^{\mathsf T},\qquad
\left|\det\frac{\partial\Pi}{\partial e}\right|=3|\det e|^8.
\end{equation}
Here the matrix index placement is chosen so $\operatorname{cof}(e)=(\det e)(e^{-1})^{\mathsf T}$; this is the same coframe density relation as QM024 with its contracted indices. Working on the invertible multiplier branch is locally permissible for the finite-dimensional coordinate change (the singular determinant hypersurface has zero Lebesgue measure) only when the oscillatory integral, gauge treatment and boundary contributions justify it. It does \emph{not} impose $\det X\ne0$ and does not by itself prove that the dual frame is a physical condensate.

\begin{proposition}[A Jacobian can manufacture a nonzero coordinate peak]\label{prop:v15-coordinatepeak}
A nonzero maximum of a dual-frame integration density is not, by itself, an invariant phase-selection criterion. Indeed, for $a>0$,
\begin{equation}\label{eq:v15-peak}
\int_0^\infty d\Pi\,e^{-a\Pi^2}
=\int_0^\infty 3b^2\,db\,e^{-ab^6},\qquad \Pi=b^3.
\end{equation}
The first integrand peaks at $\Pi=0$ while the transformed density peaks at $b=(3a)^{-1/6}>0$; the integrals and every transformed observable are identical. Thus a finite-dimensional auxiliary-frame saddle defined only as the pointwise maximum of the integration density can be a coordinate artifact.
\end{proposition}
The exact tests at $a=0.5,1,2$ agree to $10^{-12}$. The higher-dimensional cofactor Jacobian Eq.~\eqref{eq:v15-frame-dual} makes a consistent measure and gauge-invariant phase criterion particularly important; the one-dimensional example is a warning, not a claimed derivation of a specific QGD saddle.

\subsection{The actual coupled saddle equations, including the massive-source mismatch [D/C]}
\label{sec:v15-saddles}
For the Euclidean-sign formal first-order representation, let
\begin{equation}\label{eq:v15-aux}
S_{\rm aux}[\psi,\omega,X,\Pi]
=\int d^4x\left\{\kappa\det X+\Pi:(X-E[\psi,\omega])
+\sum_i m_i\det X\,\bar\psi_i\psi_i\right\}.
\end{equation}
After integrating fermions with an admissible initial-state prescription, the schematic effective functional is
\begin{equation}\label{eq:v15-Gamma}
\Gamma_{\rm coll}[X,\Pi,\omega;\rho]
=\int d^4x\{\kappa\det X+\Pi:X\}-\operatorname{Tr}_{\rho}\log\mathcal D[\Pi,X,\omega]
+\Gamma_{\rm measure}+\Gamma_{\rm other},
\end{equation}
where $\mathcal D$ includes a common kinetic density $-\Pi^\mu{}_a\gamma^aD_\mu$ and species masses proportional to $m_i\det X$. The trace-log alone is justified only for suitable Gaussian states and specified regulator/contour conditions; $\Gamma_{\rm other}$ denotes the required state and gauge terms, not an assumed Einstein action. In particular, the proper-time coefficient obtained by \emph{setting} $X=e$ and assuming a standard curved Dirac operator is not automatically the coefficient of the full off-shell $\Gamma_{\rm coll}$.

The two distinct stationary conditions read formally
\begin{equation}\label{eq:v15-eulers}
\boxed{X^a{}_{\mu}=\langle E^a{}_{\mu}\rangle_{\rho,\Pi,X},\qquad
\Pi^\mu{}_a=-\left(\kappa+\sum_i m_i\langle\bar\psi_i\psi_i\rangle_{\rho,\Pi,X}\right)\operatorname{cof}(X)^\mu{}_a,}
\end{equation}
up to counterterms and constraints in the regulated theory. These equations must be solved jointly with the source state and connection. A conditional aligned saddle $X=c e$ and $\Pi=\operatorname{cof}(e)$ implies
\begin{equation}\label{eq:v15-c}
\boxed{-\kappa_{\rm eff}c^3=1,\quad \kappa_{\rm eff}:=\kappa+\sum_i m_i\langle\bar\psi_i\psi_i\rangle,
\qquad m_{i,\rm eff}=m_i c^4.}
\end{equation}
This fixes a relation between the two frames \emph{if} an aligned saddle exists; it neither forces alignment nor supplies a stable vacuum. For $\kappa_{\rm eff}=-8,-1,-1/8$, respectively, the aligned positive solutions are $c=1/2,1,2$, and the common mass rescalings are $1/16,1,16$. The signs and contours require a physically specified Lorentzian/Euclidean continuation.

\subsection{A limited homogeneous phase-selection obstruction from diffeomorphism invariance [D/C]}
\label{sec:v15-homogeneous}
\begin{proposition}[Derivative-free local frame potential cannot alone stabilize a positive scale]\label{prop:v15-homogeneous}
For a \emph{single invertible} coframe with no extra background tensors or scalar order parameters, every local diffeomorphism- and local-Lorentz-invariant, zero-derivative density is proportional to $\det e$. Accordingly the local homogeneous ansatz $e^a{}_{\mu}=b\delta^a{}_{\mu}$, $b>0$, has
\begin{equation}\label{eq:v15-volume}
V_{\rm local}(b)=\Lambda_{\rm eff}b^4,\qquad
\frac{dV_{\rm local}}{db}=4\Lambda_{\rm eff}b^3.
\end{equation}
There is no isolated positive-$b$ stationary minimum: if $\Lambda_{\rm eff}\ne0$ none exists in this ansatz, and if $\Lambda_{\rm eff}=0$ the local scale direction is flat (subject to coordinate, boundary and gauge identifications). This does not rule out a geometric phase selected by nonlocal quantum effects, multiple collective fields, constraints, finite-density sources, boundary conditions or a nontrivial continuum limit.
\end{proposition}
\begin{proof} Without derivatives or other nontrivial fields, contractions of the coframe with its inverse and invariant tensors produce numerical constants; the only density carrying the coordinate Jacobian is the volume form $\epsilon_{abcd}e^a\wedge e^b\wedge e^c\wedge e^d$, proportional to $\det e\,d^4x$. The single-scale restriction and derivative follow immediately.\end{proof}
A pointwise unimodular \emph{bosonic} constraint $\det e=v>0$ kinematically excludes $e=0$, but it is an additional restriction, not a proof that the original unconstrained microscopic fermion theory dynamically preferred that phase. Nor can one equate $\det E=v$ pointwise in a finite Grassmann algebra, by QM024.

\subsection{What a source-conditioned causal collective action would actually require [C/O]}
\label{sec:v15-CTP}
For an initial density operator $\widehat\rho_{\rm src}$, the correct causal generating functional is the closed-time-path object
\begin{equation}\label{eq:v15-CTP}
Z_{\rho}[J_+,J_-]=\operatorname{Tr}\!\left(U_{J_+}\widehat\rho_{\rm src}U^\dagger_{J_-}\right),\qquad Z_{\rho}[J,J]=\operatorname{Tr}\widehat\rho_{\rm src}=1.
\end{equation}
Any exact $X_\pm,\Pi_\pm$ insertion must preserve this identity, the initial-state boundary terms, and all source differentiations. The causal mean-field equation is obtained from the difference-field variation of the CTP collective effective action at coincident forward and backward fields. In an actual gauge theory the physical phase should be characterized by gauge-invariant, relationally defined metric/composite correlators, not a gauge-unfixed local $\langle e^a{}_{\mu}\rangle$ alone.

\textbf{Required acceptance tests:} (i) a finite-regulator fermion and connection definition plus a consistent auxiliary distribution, contour and Jacobian; (ii) independent self-consistent $X,\Pi,\omega$ saddle equations including mass terms and source state; (iii) an observable nondegenerate geometric phase, free-energy comparison and proper $N$/volume limits; (iv) gauge-invariant low-energy spin-2 and torsion spectra and Ward identities; (v) causal exterior response to a specified source; (vi) one common kinetic geometry for fermions and photons with the microscopic gauge action specified; (vii) only then the regular QGD lift to $[\Sigma]$ using the founding metric, including $\mathcal Q_{ab}$ if active.

\subsection{QM025 result and physical decision [D/N/C/O]}
\label{sec:v15-verdict}
\textbf{Success:} The exact finite-regulator pushforward, its source-generated observables, and the invertible dual-frame coordinate algebra can be distinguished rigorously. The two coupled auxiliary Euler equations are now explicit, including the source-dependent mass/frame mismatch. The need for a nontrivial physical phase criterion follows from the coordinate-peak example and from the homogeneous covariant potential argument.\par
\textbf{Obstruction:} Neither a formal delta insertion, nor an invertible multiplier parameterization, nor a bosonic unimodular restriction proves that the original microscopic theory has dynamically selected a nondegenerate collective vacuum. The single-frame derivative-free potential cannot accomplish this alone. This is a failure of the proposed \emph{automatic closure proof}, not a no-go theorem for spinor pregeometry or QGD.\par
\textbf{Next decisive calculation:} Derive the \emph{regulated source-dependent} spinor/connection effective action with a specified diffeomorphism-respecting regulator, calculate a gauge-invariant physical frame/composite two-point spectrum and compare geometric to nongeometric phases while retaining both $X$ and $\Pi$. The transition and its massless spin-2 response must arise in the same microscopic theory. No claim of a stable phase is registered until that test succeeds.
% ==== END QM025 ====

% ===== QM026: exact determinant model and the physical geometric-phase gap =====
\clearpage
\section{QM026: Controlled large-$N$ collective ordering, source anisotropy and the missing geometric spectrum}
\label{sec:v16-qm026}
\textbf{Research status [D/N/C/O].} The QM025 obstruction to automatic geometric-phase selection must not be read as a prohibition against collective ordering. We now test a fully specified, finite Grassmann projection of the determinant interaction, derive its complete finite-$N$ generating function, and perform the two noncommuting source/large-$N$ limits exactly. The positive result is a composite \emph{order parameter}, not yet an invertible spacetime coframe in the original four-dimensional spinor-pregeometry model M. A second test proves that the determinant alone does not choose all the frame weights; a third makes explicit that a local ordered state need not have a propagating collective field. A separate conditional gauge-connection test examines photon universality in a nondegenerate induced-frame phase. All new scientific claims are scoped accordingly.

\subsection{A specified Grassmann determinant projection [D]}
\label{sec:v16-model}
Let $\chi_{ai},\bar\chi_{ai}$ be independent Grassmann generators, $a=1,\ldots,4$ and $i=1,\ldots,N$, with the Berezin integration convention $\int d\bar\chi_{ai}\,d\chi_{ai}\,\bar\chi_{ai}\chi_{ai}=1$ (an overall common sign is fixed consistently). Define the four Grassmann-even bilinears and the diagonal composite matrix
\begin{equation}\label{eq:v16-E}
 S_a=\sum_{i=1}^N\bar\chi_{ai}\chi_{ai},\qquad
 E=\operatorname{diag}(S_1,S_2,S_3,S_4),\qquad D=\det E=\prod_{a=1}^4S_a .
\end{equation}
The exactly finite generating function is
\begin{equation}\label{eq:v16-Zdef}
 Z_N(\mathbf j,\kappa)=\int D\chi D\bar\chi\;
 \exp\!\left[\kappa D+\sum_{a=1}^4j_a S_a\right].
\end{equation}
This is a \emph{zero-dimensional diagonal-channel projection} of the determinant-composite interaction. It has neither $\gamma^aD_\mu$ nor a Lorentz connection, spacetime derivative, photon, metric or spin-two propagator. It is not a complete regulator of $S_M=\kappa\int\det E[\psi,\omega]$; no result about the full four-dimensional geometric phase will be inferred without new calculations.

\begin{proposition}[Exact finite-$N$ determinant generating function]\label{prop:v16-Zexact}
For all integers $N\ge1$, and ordinary commuting sources $j_a$,
\begin{equation}\label{eq:v16-Zexact}
\boxed{Z_N(\mathbf j,\kappa)
 =\sum_{k=0}^N\frac{\kappa^k}{k!}
 \left[\frac{N!}{(N-k)!}\right]^4
 \left(\prod_{a=1}^4j_a\right)^{N-k}.}
\end{equation}
In particular $Z_N(\mathbf 0,\kappa)=\kappa^N(N!)^3$ for $\kappa\ne0$, and the zero-source expectations obey
\begin{equation}\label{eq:v16-moments}
 \left.\langle S_a\rangle\right|_{\mathbf j=0}=0,\qquad
 \left.\langle D\rangle\right|_{\mathbf j=0}=\frac N\kappa.
\end{equation}
\end{proposition}
\begin{proof}
Expand the determinant interaction in powers of $k$. Because $S_a^{N+1}=0$, $k\le N$. For a single channel,
$\int D\chi_aD\bar\chi_a\, S_a^k\exp(j_a S_a)=[N!/(N-k)!]j_a^{N-k}$: choose $k$ distinct bilinears from the $N$ available and saturate the others with source monomials. Multiply the four independent channel factors for fixed $k$. The zero-source partition is the $k=N$ term, and its log derivatives with respect to $j_a$ and $\kappa$ give Eq.~\eqref{eq:v16-moments}.
\end{proof}
A nonzero determinant expectation with vanishing one-bilinear means therefore follows from an exact finite-Grassmann calculation. It cannot by itself be interpreted as a nonsingular \emph{classical random matrix}: the underlying Berezin functional is not a positive distribution of ordinary c-number frames (QM025).

\subsection{An actual order-of-limits result at large $N$ [D/N]}
\label{sec:v16-largen}
Set $\kappa=g/N^3$, $g>0$, and for now choose all four sources $j_a=j>0$. The extensive scaling matches the schematic $E=O(N)$ and $\det E=O(N^4)$ of the original collective discussion. Denote the $k$th summand of Eq.~\eqref{eq:v16-Zexact} by $W_k$. Their exact ratio is
\begin{equation}\label{eq:v16-ratio}
 \frac{W_{k+1}}{W_k}=\frac{g}{N^3}\frac{(N-k)^4}{(k+1)j^4}.
\end{equation}
This is monotonically decreasing with $k$, so the positive-weight $k$-distribution is log-concave. With $x=k/N$, its unique bulk maximum concentrates in the $N\to\infty$ limit around
\begin{equation}\label{eq:v16-saddle}
\boxed{g(1-x)^4=xj^4,\quad 0<x<1.}
\end{equation}
Since $\partial_j\log Z_N=4(N-\langle k\rangle)/j$, the normalized bilinear response tends to
\begin{equation}\label{eq:v16-b}
 b(j;g):=\lim_{N\to\infty}\frac{\langle S_a\rangle_{N,j}}N=\frac{1-x(j;g)}{j}.
\end{equation}
Writing $1-x=j y$ in Eq.~\eqref{eq:v16-saddle} and taking $j\to0^+$ gives $g y^4=1$. Thus
\begin{equation}\label{eq:v16-noncommute}
\boxed{\lim_{j\to0^+}\lim_{N\to\infty}\frac{\langle S_a\rangle_{N,j}}N=g^{-1/4},\qquad
\lim_{N\to\infty}\lim_{j\to0^+}\frac{\langle S_a\rangle_{N,j}}N=0.}
\end{equation}
This is a genuine source-selected large-$N$ branch in the \emph{defined projection}; it is not a demonstration of a gauge-invariant physical frame in $M$. For $j<0$ the limiting response has the opposite sign, so the extensive source-dependent free energy has a cusp at the origin (the exact finite-$N$ polynomial is analytic). The numerical verifier calculates Eq.~\eqref{eq:v16-Zexact} using log-sum-exp for $N=4,16,64,256,1024$ and confirms the convergence and source derivatives.

\begin{figure}[htbp]
 \centering\includegraphics[width=.88\linewidth]{phase_order_limits.png}
 \caption{Exact finite-$N$ determinant-model bilinear response (solid lines) compared with the analytic large-$N$ branch (dashed). The small-source and large-$N$ limits do not commute; this establishes ordering only in the named zero-dimensional projection.}
 \label{fig:v16-order}
\end{figure}
The determinant is already nonzero as an invariant composite at zero source:
\begin{equation}\label{eq:v16-det}
\left.\left\langle\frac{\det E}{N^4}\right\rangle\right|_{\mathbf j=0}=
\frac{1}{g},\qquad \left.\langle E\rangle\right|_{\mathbf j=0}=0\quad(N<\infty).
\end{equation}
This explicitly demonstrates why a source-free one-point frame expectation is not an adequate finite-system definition of a collective geometric phase.

\subsection{Important failure: the determinant does not select a unique weighted frame [D/N]}
\label{sec:v16-anisotropy}
The exact polynomial Eq.~\eqref{eq:v16-Zexact} depends on independent sources only through $\prod_a j_a$. With $j_a=t c_a>0$ and $\kappa=g/N^3$, its large-$N$ saddle obeys $g(1-x)^4=x t^4\prod_a c_a$. Taking the large-$N$ limit first and then $t\to0^+$ yields
\begin{equation}\label{eq:v16-anisotropy}
\boxed{b_a=\lim_{t\to0^+}\lim_{N\to\infty}\frac{\langle S_a\rangle}{N}
=\frac{(\prod_{b=1}^4c_b/g)^{1/4}}{c_a},\qquad
\prod_{a=1}^4b_a=\frac1g.}
\end{equation}
For $g=1$, source directions $(c_a)=(1,1,1,1)$ and $(1,2,1/2,1)$ select respectively $(b_a)=(1,1,1,1)$ and $(1,1/2,2,1)$, both with the same determinant.
\begin{figure}[htbp]\centering
\includegraphics[width=.80\linewidth]{anisotropic_frame_selection.png}
\caption{Different channel-source directions select different weighted-frame shapes in the exact large-$N$ projection. The determinant tends to the same value as $t\to0^+$ in both branches.}\label{fig:v16-anisotropy}
\end{figure} Exact $N=4096$ evaluations approach these limits as the sources decrease. Thus the determinant interaction can fix a volume-like collective product without selecting all frame-shape data. Some constant-frame shape freedom may be related to coordinate or gauge transformations in a fully covariant spacetime model; the reduced projection lacks that gauge structure and cannot classify which shapes are physical. The \emph{universal metric selection problem remains open}.

\subsection{Connected source response and the missing spacetime propagator [D/N/O]}
\label{sec:v16-response}
At equal sources $j>0$ define the exact probability weights $p_k=W_k/\sum_lW_l$ as a computational device for the \emph{coefficient sum} (not a positive probability measure on fermion fields). For two bilinear sources belonging to different channels and different species, the static mixed source derivative is
\begin{equation}\label{eq:v16-cross}
\boxed{\frac{\partial^2\log Z_N}{\partial j_{ai}\partial j_{b\ell}}
=\frac{\operatorname{Var}_{p}(k)}{N^2j^2}>0,\qquad a\ne b,\ i\ne\ell.}
\end{equation}
At $g=1,j=0.2$ and $N=4,16,64,256,1024$, the values are approximately $0.345657$, $0.0713862$, $0.0175760$, $0.00438056$ and $0.00109435$. The microscopic determinant therefore creates a \emph{shared static interaction across otherwise distinct constituent species}; the individual cross-derivative scales as $O(1/N)$ at fixed positive source. This is not a causal retarded two-point function.

A spatial array of \emph{independent copies} of this exact projected model has $Z=\prod_x Z_{N,x}$ and therefore
\begin{equation}\label{eq:v16-ultralocal}
\boxed{\langle D(x)D(y)\rangle_c=0\quad (x\ne y)}
\end{equation}
regardless of the nonzero local determinant condensate. Its connected Fourier spectrum has no $k^{-2}$ spin-two pole. This is an exact limitation of the independent-cell extension, \emph{not} of the full four-dimensional $M$, whose derivative spinor composite and Lorentz connection can couple neighboring cells. A physically propagating exterior QGD field requires an independently computed spatial/causal collective kernel, not merely a local condensate.

\subsection{Gauge-orbit check: a vanishing mean frame does not imply absent geometry [D/N]}
\label{sec:v16-gauge}
As a separate Euclidean local-Lorentz illustration, let $e=R e_0$ with $R\in SO(4)$ Haar-distributed and $e_0$ invertible. Group averaging yields
\begin{equation}\label{eq:v16-haar}
\boxed{\langle e^a{}_{\mu}\rangle_R=0,\qquad
\left\langle e^a{}_{\mu}e^b{}_{\nu}\right\rangle_R
=\frac{\delta^{ab}}4(e_0^{\mathsf T}e_0)_{\mu\nu},\qquad
\langle e^{\mathsf T}e\rangle_R=e_0^{\mathsf T}e_0.}
\end{equation}
A 12,000-sample Haar check has metric error below $1.2\times10^{-12}$ and average frame norm $\simeq0.022$. Thus physical phase criteria must be expressed through gauge-invariant or relational composites; $\langle e\rangle=0$ in an unfixed local gauge is not evidence of a missing geometric phase. This group-orbit calculation is separate from the zero-dimensional determinant model and cannot establish its local-Lorentz symmetry.

\subsection{A conditional route to a shared photon/fermion coframe [D/C/O]}
\label{sec:v16-gaugefield}
The microscopic proposal M originally includes spinors and a Lorentz connection but not an independently specified electromagnetic gauge connection. A \emph{specified extension} replaces $D_\mu\psi_i$ in its defining composite by $D_\mu\psi_i-iq_i A_\mu\psi_i$ while preserving the same composite determinant interaction. The exact auxiliary kinetic term $-\Pi:E[\psi,\omega,A]$ then supplies a common coframe to every charged Dirac species, with its usual $q_iA_\mu$ coupling. If a nondegenerate frame phase exists and the fermion determinant admits a covariant low-energy expansion, the gauge-field part of the Euclidean heat-kernel coefficient of a four-component Dirac species is
\begin{equation}\label{eq:v16-a4}
\boxed{\operatorname{tr}\left(\frac1{12}\Omega_{\mu\nu}\Omega^{\mu\nu}
+\frac12\mathsf E_F^2\right)
=\frac23 q_i^2 F_{\mu\nu}F^{\mu\nu},\quad
\Omega_{\mu\nu}=-iq_iF_{\mu\nu}I_4,\quad
\mathsf E_F=-\frac{iq_i}{2}\gamma^{\mu\nu}F_{\mu\nu}.}
\end{equation}
The independent gamma-matrix calculation reproduces this expression to relative error $6\times10^{-16}$. Its logarithmic fermionic determinant therefore induces a Maxwell-type $F^2$ action using the \emph{same} chosen coframe at the stated derivative order. This is a constructive conditional common-metric mechanism, \emph{not} proof of a physically realized photon sector: the gauge field is an added microscopic variable, measure and spin-connection treatment remain open, and the interacting two-field/massive-sector mismatches from QM024--QM025 persist.

\subsection{QM026 result and the next genuine acceptance test}
\label{sec:v16-verdict}
\textbf{Positive result [D/N].} The determinant-composite interaction itself can generate a source-selected nonzero bilinear order in a rigorously computed large-$N$ limit of an explicitly regulated finite Grassmann projection. Both its exact finite-$N$ source generating functional and the $N\to\infty$ saddle are known. Shared static cross-species responses and a nonzero composite determinant expectation follow without inserting any metric or Einstein action. This invalidates the inference that finite-Grassmann nilpotency or a vanishing finite-volume $\langle E\rangle$ forbids collective geometric ordering at large $N$.

\textbf{Failure/limitation [D/O].} The projection has no spacetime derivative, spin connection, gauge dynamics or transverse-traceless correlator. Its determinant does not uniquely select frame weights, and independent cells have identically vanishing nonlocal correlations. A separate conditional gauge-extension shows how a shared photon kinetic tensor may be induced if a physical coframe is already realized, but that realization remains open. \emph{The nondegenerate, dynamically preferred exterior quantum frame of the actual four-dimensional microscopic model M has not been demonstrated.}

\textbf{Next essential computation [O].} Construct a finite, diffeomorphism-respecting discretization of the \emph{actual} operator $E^a{}_{\mu}=\sum_i\frac i2\bar\psi_i\gamma^a\overleftrightarrow{D_\mu}\psi_i$, integrate the spin connection and all fermions or simulate with a demonstrably correct determinant/Pfaffian measure, and scan source-conditioned, gauge-invariant collective two-point spectra over increasing four-dimensional volumes. Determine whether a nondegenerate phase and its source-visible, transverse-traceless long-distance dynamics arise on the \emph{same} branch. Then check that gauge connections coupled through the same microscopic composite yield a common effective geometry. Only after those physical tests may the regular QGD map $g\to[\Sigma]$ be promoted from an algebraic lift to microscopic selection. The full founding metric with $T^a{}_{\mu}$ and $\mathcal Q_{ab}$ remains the theory's master equation.
% ===== END QM026 =====


\input{QM027_section.tex}
\input{QM031_section.tex}
\input{QM034_section.tex}
\input{QM037_section.tex}
\input{QM038_section.tex}
\clearpage
\input{QM039_section.tex}
\input{QM032_036_retrospective_section.tex}
\input{QM040_section.tex}
\input{QM041_Chapter6_Addendum.tex}
\input{QM042_status_and_inverse_synthesis.tex}
\input{QM043_coarsegraining_statistical_basis_inverse_map.tex}
\input{QN001_QN005_Action_Integration_Chapter6.tex}
\input{Chapter6_Q052_Q054_Quantum_And_DE.tex}

\clearpage


\section{Q061: final quantum-action, symmetry, and energy-status consolidation}
\noindent\textbf{Which action is being quantized?} The operational gravitational theory is the renormalized metric-mediated action $\Gamma_{\rm met}^{\rm ren}[g(\Sigma)]$ of Chapter~4 on a regular, complete, $\lambda\mathcal Q=0$ field-map branch. Its matter-induced local $R^2,C^2,E_4$ coefficients are consistently renormalized independent EFT inputs; QN004 determines a restricted Euclidean one-loop nonlocal response in a declared proper-time scheme. The same local loop contribution must not be counted both in $c_i(\mu)$ and in the remaining nonlocal/state action.

\noindent\textbf{Physical quantum response.} On a metric-shell background the Hessian obeys $H_\Sigma=L^*H_gL$ with $L=D_\Sigma g$. Nonzero quantum corrections alter the observable metric combinations of the natural invariants; local metric-preserving $O(p,q)$ changes do not automatically carry quantum energy. The historical $\mathcal Q_{ab}$ in Axiom~I has no unique off-shell value from this functional, and a calculated state-dependent $q_{\mu\nu}^{\rm back}$ is not an extra operator (Q051--Q052). Chapter~7 reconstructs explicit Boulware-state spherical backreaction; Chapter~8 reconstructs the conformal anomaly in FLRW, both within their domains.

\noindent\textbf{Causality.} A time-dependent quantum source requires the closed-time-path (in--in) response for retarded metric backreaction. QN004 Euclidean form factors are not, by themselves, such a kernel. Finite higher-derivative couplings are treated as perturbative EFT corrections; no extra propagating high-frequency gravitational mode is certified by truncating the derivative expansion.

\noindent\textbf{Dark energy.} Signed-sector symmetry, universal metric coupling, and boundary conditions cannot select a metric-null dark-energy amplitude. Q055--Q058 show that a direct positive-kinetic sector action can be internally consistent and support acceleration, but the physical selector and energy scale are new dynamical postulates, not outcomes of Axiom I or II. The simplest shift-symmetric condensate faces a strong early-radiation constraint, whereas a bounded hilltop remains a conditional quintessence-type extension. Chapter~8 contains the detailed tests. Negative static Boulware vacuum energy in Chapter~7 must not be conflated with homogeneous vacuum-like $w=-1$ energy. The original QM research archive is preserved without changing its proof status.


\end{document}
